% =====================================================================
%  Trabajo de Fin de Grado -- De la medida neutral al riesgo al
%  Teorema de Girsanov en integración estocástica. Valoración de
%  productos financieros
%  Maquetación: A4, 11 pt, interlineado 1.5, márgenes 2.5 cm.
% =====================================================================

\documentclass[11pt,a4paper,oneside]{report}
\usepackage{graphicx}  % Logo de la universidad
\usepackage{xcolor}
% --------- Codificación e idioma ------------------------------------
\usepackage[utf8]{inputenc}
\usepackage[T1]{fontenc}
\usepackage[spanish,es-noshorthands]{babel}


% --------- Maquetación ----------------------------------------------
\usepackage[a4paper,top=2.5cm,bottom=2.5cm,left=2.5cm,right=2.5cm]{geometry}
\usepackage{setspace}
\onehalfspacing                        % interlineado 1.5

% --------- Matemáticas ----------------------------------------------
\usepackage{amsmath,amssymb,amsthm,mathtools}
\usepackage{bm}
% Indicadores 1_A: se definen con un comando propio (véase más abajo).

% --------- Otros recursos -------------------------------------------
\usepackage{graphicx}
\usepackage{tikz}
\usetikzlibrary{arrows.meta,positioning,calc}
\usepackage{enumitem}
\usepackage{booktabs}
\usepackage{caption}
\usepackage[colorlinks=true,linkcolor=black,citecolor=blue,urlcolor=blue]{hyperref}
\hypersetup{
  pdftitle={De la medida neutral al riesgo al Teorema de Girsanov en integración estocástica. Valoración de productos financieros},
  pdfauthor={Roberto Cepeda Rocandio},
  pdfsubject={Trabajo de Fin de Grado en Matemáticas, Universidad de Oviedo},
  pdfkeywords={no-arbitraje, martingala, integral de Itô, teorema de Girsanov, representación de martingalas, fórmula de Black-Scholes},
  pdflang={es-ES}
}

% --------- Entornos de teoremas -------------------------------------
\theoremstyle{plain}
\newtheorem{teorema}{Teorema}[chapter]
\newtheorem{proposicion}[teorema]{Proposición}
\newtheorem{lema}[teorema]{Lema}
\newtheorem{corolario}[teorema]{Corolario}

\theoremstyle{definition}
\newtheorem{definicion}[teorema]{Definición}
\newtheorem{ejemplo}[teorema]{Ejemplo}

\theoremstyle{remark}
\newtheorem{observacion}[teorema]{Observación}

% --------- Atajos de notación ---------------------------------------
\newcommand{\R}{\mathbb{R}}
\newcommand{\N}{\mathbb{N}}
\newcommand{\Z}{\mathbb{Z}}
\newcommand{\E}{\mathbb{E}}
\newcommand{\PP}{\mathbb{P}}
\newcommand{\Q}{\mathbb{Q}}
\newcommand{\F}{\mathcal{F}}
\newcommand{\G}{\mathcal{G}}
\newcommand{\Hh}{\mathcal{H}}
\newcommand{\HH}{\mathcal{H}}
\newcommand{\1}{\mathbf{1}}
\newcommand{\dd}{\mathrm{d}}
\newcommand{\eqdef}{\coloneqq}


% =====================================================================
\begin{document}

% =====================================================================
%   PORTADA
% =====================================================================
\begin{titlepage}
\centering

\vspace*{1.5cm}

% Logo de la Universidad de Oviedo
\includegraphics[width=0.73\textwidth]{logo_uniovi.png}

\vspace{3cm}

% Título del TFG
{\LARGE\bfseries De la medida neutral al riesgo\\[0.3cm]
        al Teorema de Girsanov en integración estocástica.\\[0.3cm]
        Valoración de productos financieros\par}

\vspace{1.5cm}

% Autor
{\Large Roberto Cepeda Rocandio\par}

\vspace{2.5cm}

% Tutor
{\large Supervisado por: José Santos Domínguez Menchero\par}

\vspace{1.2cm}

% Bloque institucional
{\sc Universidad de Oviedo}\\[0.2cm]
{Facultad de Ciencias}\\[0.2cm]
{Grado en Matemáticas}

\vspace{1.2cm}

% Fecha
{Curso 2026-2027}

\vfill
\end{titlepage}

\clearpage

% Numeración romana en el front matter
\pagenumbering{roman}
% =====================================================================
%   RESUMEN / ABSTRACT
% =====================================================================
\chapter*{Resumen}
\addcontentsline{toc}{chapter}{Resumen}
\vspace*{-0.75cm}

\itshape

El presente trabajo desarrolla, de manera autocontenida y matemáticamente
rigurosa, la teoría de valoración de opciones europeas, culminando en
la demostración de la fórmula de Black-Scholes para la opción
\emph{call} europea. El enfoque adoptado se lleva a cabo de forma unificada:
plantear primero el problema en el contexto discreto del modelo
binomial multiperíodo, donde las ideas conceptuales centrales
(principio de no arbitraje, replicación, medida neutral al riesgo,
representación de martingalas) son accesibles con herramientas
elementales de probabilidad; y a continuación trasladar el problema y
su resolución al contexto continuo del modelo de Black-Scholes,
desarrollando para ello el aparato necesario del cálculo estocástico.
El enfoque elegido, vía martingalas en lugar de la vía clásica
de las ecuaciones diferenciales parciales, vuelve transparente la
traducción de las ideas del caso discreto al caso continuo, y permite
presentar la fórmula de Black-Scholes como el límite natural de la
fórmula binomial multiperíodo cuando los pasos temporales se hacen
infinitamente pequeños. Así. el desarrollo requerirá introducir y construir con detalle buena parte del cálculo estocástico necesario para el caso continuo: la integral de Itô, el teorema de Girsanov y el teorema de representación de martingalas. Estas herramientas no solo son necesarias para los fines de este trabajo, sino que ocupan un lugar central en la teoría de procesos estocásticos. La memoria recoge además la derivación de la
EDP clásica de Black-Scholes a partir de la fórmula obtenida vía martingalas, mostrando así
la consistencia entre ambas formulaciones. 

\upshape

\vspace{0.25cm}

\section*{Abstract}

\itshape

This work develops, in a self-contained and mathematically rigorous
manner, the theory of European option pricing, culminating in the
proof of the Black-Scholes formula for the European \emph{call}
option. The approach adopted is carried out in a unified way: we
first formulate the problem in the discrete setting of the
multi-period binomial model, where the central conceptual ideas
(no-arbitrage principle, replication, risk-neutral measure,
martingale representation) are accessible with elementary probability
tools; we then translate the problem and its resolution to the
continuous setting of the Black-Scholes model, developing along the
way the necessary apparatus of stochastic calculus. The chosen
approach, via martingales rather than the classical route of partial
differential equations, makes transparent the translation of the
ideas from the discrete case to the continuous case, and allows the
Black-Scholes formula to be presented as the natural limit of the
multi-period binomial formula when the time steps become infinitely
small. Thus, the development requires introducing and constructing in detail a substantial part of the stochastic calculus needed for the continuous setting: the Itô integral, Girsanov’s theorem and the martingale representation theorem. These tools are not only necessary for the purposes of this work, but also play a central role in the theory of stochastic processes. The dissertation also includes the derivation of the classical Black-Scholes PDE from the formula obtained via martingales, thus showing the consistency between the two formulations.

\upshape

\vspace{0.3cm}
\noindent\textbf{Palabras clave:} no-arbitraje, martingala, integral
de Itô, teorema de Girsanov, representación de martingalas, fórmula
de Black-Scholes.

\vspace{0.15cm}
\noindent\textbf{Keywords:} no-arbitrage, martingale, Itô integral,
Girsanov's theorem, martingale representation, Black-Scholes formula.

% =====================================================================
%   ÍNDICE
% =====================================================================
\setcounter{tocdepth}{2}    % Índice: capítulos + secciones + subsecciones
\tableofcontents

% =====================================================================
%   NOTACIÓN
% =====================================================================
\chapter*{Notación}
\addcontentsline{toc}{chapter}{Notación}

A lo largo del presente trabajo se utilizará la siguiente notación.

\vspace{0.3cm}

\noindent\textbf{Espacios de probabilidad y filtraciones}
\begin{center}
\begin{tabular}{@{}p{4cm}p{10cm}@{}}
$(\Omega, \F, \PP)$ & Espacio de probabilidad. \\
$(\F_t)_{t \geq 0}$ & Filtración: codifica el flujo de información. \\
$\F_t^W$ & Filtración natural del browniano $\{W_t\}$. \\
$\PP, \Q$ & Medida del mercado y medida neutral al riesgo. \\
$\Q \sim \PP$ & Las medidas $\Q$ y $\PP$ son equivalentes. \\
$d\Q/d\PP$ & Derivada de Radon-Nikodym de $\Q$ respecto a $\PP$. \\
\end{tabular}
\end{center}

\vspace{0.2cm}

\noindent\textbf{Esperanzas e indicadores}
\begin{center}
\begin{tabular}{@{}p{4cm}p{10cm}@{}}
$\E_\PP[X]$ & Esperanza de $X$ bajo $\PP$. \\
$\E_\Q[X]$ & Esperanza de $X$ bajo $\Q$. \\
$\E[X \mid \G]$ & Esperanza condicional de $X$ respecto a $\G$. \\
$\1_A$ & Función indicadora del suceso $A$. \\
\end{tabular}
\end{center}

\vspace{0.2cm}

\noindent\textbf{Procesos estocásticos}
\begin{center}
\begin{tabular}{@{}p{4cm}p{10cm}@{}}
$\{W_t\}_{t \geq 0}$ & Movimiento browniano estándar bajo $\PP$. \\
$\{\widetilde W_t\}$ & Movimiento browniano bajo $\Q$ (Girsanov). \\
$\{S_t\}$ & Precio del activo subyacente. \\
$\widetilde S_t = e^{-rt} S_t$ & Precio descontado del subyacente. \\
$B_t = e^{rt}$ & Precio del bono libre de riesgo. \\
$\{V_t\}$ & Valor del portfolio. \\
$(\phi_t, \psi_t)$ & Cantidades de subyacente y bono del portfolio. \\
$\{L_t\}$ & Proceso de densidad de Radon-Nikodym. \\
\end{tabular}
\end{center}

\vspace{0.2cm}

\noindent\textbf{Cálculo estocástico}
\begin{center}
\begin{tabular}{@{}p{4cm}p{10cm}@{}}
$\int_0^t \phi_s\, dW_s$ & Integral estocástica de Itô. \\
$\HH^2[0, T]$ & Procesos con $\E[\int_0^T \phi_s^2\, ds] < \infty$. \\
$\langle X, Y\rangle_t$ & Covariación cuadrática entre $X$ e $Y$. \\
$dX_t = a_t dt + b_t dW_t$ & Notación diferencial para procesos de Itô. \\
\end{tabular}
\end{center}

\newpage

\noindent\textbf{Parámetros del modelo de Black-Scholes}
\begin{center}
\begin{tabular}{@{}p{4cm}p{10cm}@{}}
$r$ & Tipo de interés constante libre de riesgo. \\
$\mu, \sigma$ & Rendimiento esperado y volatilidad del subyacente. \\
$T,\ K$ & Vencimiento y \emph{strike} de la opción. \\
$\theta = (\mu - r)/\sigma$ & Precio de mercado del riesgo. \\
$N(\cdot)$ & Función de distribución de la normal estándar. \\
$d_1, d_2$ & Coeficientes de la fórmula de Black-Scholes. \\
$L^2$ & Espacio de variables de cuadrado integrable. \\
$\eqdef$ & Igualdad por definición; c.s. abrevia ``casi seguramente''. \\
\end{tabular}
\end{center}

% Inicio del cuerpo del trabajo: numeración arábiga
\pagenumbering{arabic}
\setcounter{page}{1}

% =====================================================================
%   CAPÍTULO 1 — INTRODUCCIÓN
% =====================================================================
\chapter{Introducción}\label{cap:intro}

\section{Motivación y contexto}
\label{sec:motivacion}

Una \emph{opción financiera} es un contrato que otorga a su tenedor el
derecho, que no la obligación, de comprar o vender un activo subyacente
(una acción, un índice, una divisa, una materia prima) a un precio
prefijado en una fecha futura. La opción más sencilla, la \emph{call
europea}, otorga el derecho de compra a un precio $K$ (el
\emph{strike}) en una fecha $T$ (el \emph{vencimiento}). El tenedor de
la opción ejercerá el derecho si y solo si, llegada la fecha $T$, el
precio del subyacente es superior a $K$; en tal caso obtendrá un
beneficio igual a la diferencia. La cuestión central, desde el punto
de vista matemático, es la siguiente: ¿cuál es el precio justo de
este contrato en el instante de su emisión?

La pregunta admite una formulación matemática precisa, anclada en el
\emph{principio de no arbitraje}: el precio del contrato debe ser tal
que no permita a ningún operador del mercado obtener un beneficio
seguro sin asumir riesgo alguno. Bajo este principio, sumado a una
descripción matemática de la dinámica del subyacente, el precio queda
unívocamente determinado. La determinación de este precio, y la
construcción de la estrategia que permite a quien emite la opción
\emph{replicar} su valor, es el problema que articula el presente
trabajo.

\medskip

\noindent\textbf{Apunte histórico.} La primera formulación matemática
del problema se debe a Louis Bachelier, quien en su tesis doctoral
\emph{Théorie de la spéculation} \cite{bachelier} modeló la dinámica
de los precios bursátiles mediante el movimiento browniano. La obra
de Bachelier, ignorada en su tiempo, fue redescubierta a mediados del
siglo XX por Paul Samuelson y otros autores, quienes refinaron el
modelo sustituyendo el movimiento browniano aritmético por el
\emph{movimiento browniano geométrico}, que asegura la positividad de
los precios. El problema fue resuelto definitivamente en 1973 por
Fischer Black, Myron Scholes y Robert Merton, en dos artículos
seminales \cite{black-scholes,merton} reconocidos en 1997 con el
Premio Nobel de Economía concedido a Scholes y Merton (Black había
fallecido en 1995). La fórmula de Black-Scholes, que constituye el
resultado central del presente trabajo, es hoy uno de los desarrollos
matemáticos más influyentes del último medio siglo.

\medskip

\noindent\textbf{Por qué es matemáticamente interesante.} Más allá
de su relevancia práctica, el problema requiere un aparato
matemático rico: probabilidad avanzada (martingalas, esperanza
condicional), análisis estocástico (browniano, integral y fórmula
de Itô), análisis funcional (densidad en $L^2$) y, en su versión
EDP, ecuaciones parabólicas. La fórmula de Black-Scholes es punto
de convergencia de varias áreas centrales de la matemática moderna.

\medskip

\noindent\textbf{Un planteamiento más general de lo que parece.} Conviene
anticipar una idea que está presente a lo largo de todo el trabajo y que va más allá del
problema financiero que la motiva. La noción de precio justo nos llevará,
tanto en el modelo discreto como en el continuo, a sustituir la medida
de probabilidad del mercado por otra equivalente (la \emph{medida
neutral al riesgo}) bajo la cual el precio descontado del subyacente
se convierte en una \emph{martingala}, uno de los conceptos
fundamentales de la probabilidad moderna. El cambio de medida que
produce esa martingala descansa sobre dos resultados de carácter general:
el teorema de Radon-Nikodym, que garantiza la existencia de la densidad
entre ambas medidas, y el teorema de Girsanov, que describe dicho cambio
en el contexto del movimiento browniano. Ninguna de estas herramientas
es exclusiva de la valoración: el principio de sustituir una medida por
otra equivalente para transformar un problema difícil en uno tratable es
de uso corriente en probabilidad y en sus aplicaciones, por ejemplo,
en la reducción de varianza de las simulaciones Monte Carlo de sucesos
poco probables, mediante la técnica de \emph{importance sampling}
(véase, e.g., \cite{karatzas-shreve}). En este sentido, el problema de valoración permite introducir herramientas matemáticas que tienen aplicaciones mucho más amplias.

\section{Objetivos y enfoque del trabajo}
\label{sec:objetivos}

El presente trabajo se propone articular, de manera autocontenida y
matemáticamente rigurosa, la teoría de valoración de opciones
europeas, culminando con la demostración de la fórmula de
Black-Scholes para la \emph{call} europea. La articulación se llevará a cabo de forma unificada: presentar primero el problema en el
contexto \emph{discreto} del modelo binomial multiperíodo, donde las
ideas conceptuales son accesibles con herramientas elementales de
probabilidad, y a continuación trasladar el problema y su resolución
al contexto \emph{continuo} del modelo de Black-Scholes, mediante el
desarrollo del cálculo estocástico necesario.

\medskip

\noindent\textbf{Justificación del plan discreto-continuo.} El orden
expositivo elegido es pedagógico, no histórico: cronológicamente, la
fórmula de Black-Scholes (1973) es anterior al modelo binomial
multiperíodo (Cox, Ross y Rubinstein, 1979). Sin embargo, desde el
punto de vista matemático el modelo binomial es estructuralmente más
simple: permite articular las ideas centrales de la teoría (cambio de
medida, replicación, no arbitraje) con herramientas de probabilidad
elemental, sin necesidad del aparato del cálculo estocástico. Este
hecho es lo que justifica el plan adoptado: una vez asimilada la lógica del
problema en el contexto discreto, el paso al continuo se reduce a una
traducción técnica de las herramientas, conservando intactas las ideas
conceptuales.

\medskip

\noindent\textbf{La vía de las martingalas.} La teoría continua de
Black-Scholes admite dos vías clásicas: martingalas (cambio de
medida y teoremas de representación) y EDP (cobertura instantánea).
Optamos por la primera pues admite traducción directa desde el
caso discreto. Esta elección tiene, además, una ventaja que va más allá de la mera exposición: las
herramientas que exige, la integral de Itô, el teorema de Girsanov y
el teorema de representación de martingalas, son objetos matemáticos
con interés y utilidad propios, y constituyen buena parte del contenido técnico del trabajo. La vía EDP se recupera al final
(Sección~\ref{sec:edp_bs}) como consecuencia de la formulación
martingala.

\medskip

\noindent\textbf{Lo que el trabajo no aborda.} Por razones de extensión
y de coherencia con el plan elegido, el trabajo se limita a opciones
\emph{europeas} sobre un único subyacente, en el modelo de Black-Scholes
con parámetros constantes. Quedan, por consiguiente, fuera del alcance
del trabajo varias extensiones naturales —opciones americanas,
modelos con volatilidad estocástica o saltos, métodos numéricos para
\emph{payoffs} sin fórmula cerrada, calibración a datos de mercado—,
recogidas brevemente en el Capítulo~\ref{cap:conclusiones} como líneas
de continuación.

\section{Fuentes y aportación propia}
\label{sec:aportacion}

El hilo expositivo de este trabajo sigue esencialmente el de
A.~Etheridge~\cite{etheridge}, del que tomamos tanto la estructura
global (el paso del modelo binomial discreto al modelo de
Black-Scholes por la vía de las martingalas) como el orden en que se
introducen las principales ideas: la motivación de la integral estocástica a partir
del no-arbitraje, el cambio de medida como paso crucial de la valoración
y el teorema de representación de martingalas. La introducción al cálculo estocástico, que se limita al movimiento browniano, sigue igualmente esta estructura.

En cuanto a la parte técnica, esta se apoya en varios textos introductorios sobre cálculo estocástico y teoría de la probabilidad. Para el movimiento browniano (construcción, principales propiedades y filtración aumentada) seguimos principalmente a S.~E.~Shreve~\cite{shreve}. Los resultados centrales del cálculo estocástico que utilizamos (la construcción de la integral de Itô y su isometría, la fórmula de Itô y el teorema de Girsanov) se presentan directamente para el movimiento browniano, siguiendo los textos introductorios de la materia~\cite{oksendal,shreve}
y el ya mencionado Etheridge~\cite{etheridge}. Una formulación más general, válida para martingalas continuas de cuadrado integrable y semimartingalas, puede consultarse en I.~Karatzas y
S.~E.~Shreve~\cite{karatzas-shreve}. Los resultados que utilizamos en el caso browniano se obtienen así como caso particular, con $\langle W\rangle_t = t$. Para la teoría de la medida y de martingalas (esperanza condicional, teorema de Radon-Nikodym y martingalas en tiempo discreto) seguimos a D.~Williams~\cite{williams}. Por último, las referencias históricas~\cite{bachelier,black-scholes,merton} se utilizan únicamente para contextualizar el problema.

Dado que los resultados que se presentan en este trabajo son clásicos y bien conocidos, la aportación no está tanto en obtener resultados nuevos, cuanto en su selección, organización y desarrollo. Algunas demostraciones que las fuentes presentan de forma resumida, dan por conocidas o formulan en un nivel de generalidad mayor del que aquí se necesita, se han desarrollado con más detalle y adaptadas al problema que nos ocupa. Este es el caso de la construcción de la integral de Itô del Capítulo~\ref{cap:itocalculus}, con la isometría demostrada paso a paso, la extensión a $\mathcal{H}^2$ y la continuidad de la integral indefinida; de la fórmula de Itô, detallando las herramientas de análisis que intervienen en su demostración; y del argumento constructivo de no arbitraje que motiva la introducción de una nueva integral (Sección~\ref{sec:motivacion_ito}). Además, se ha procurado mantener una notación uniforme entre las distintas fuentes y relacionar de forma clara los casos discreto y continuo. Así, el cambio de medida de Girsanov (Sección~\ref{sec:girsanov}), por ejemplo, se presenta como la versión continua del ajuste de probabilidades de las distintas trayectorias del modelo binomial, y no como un resultado independiente.

Algunos de los aspectos que requirieron más trabajo fueron, ya en el modelo discreto, la propia idea de \emph{no-arbitraje} (Capítulo~\ref{cap:discreto}). Entender por qué la ausencia de beneficios seguros sin inversión permite determinar un precio, y cómo este principio se convierte en la base de la teoría, requiere un cambio de perspectiva que no es inmediato. En el modelo continuo, la diferencia entre la filtración browniana natural y su aumentación (Capítulo~\ref{cap:brown}) también requirió especial atención, ya que los resultados posteriores precisaban de una filtración que cumpliera determinadas propiedades técnicas. Por último, la aparición del \emph{término
de Itô} hizo necesario explicar por qué la variación cuadrática del movimiento browniano no puede despreciarse y de dónde surge la corrección de segundo orden $(dW_t)^2 = dt$, que distingue el cálculo estocástico del cálculo ordinario.

\section{Estructura de la memoria}
\label{sec:estructura}

La memoria se estructura en seis capítulos: introducción, parte
discreta (Capítulo~\ref{cap:discreto}), fundamentos continuos
(Capítulos~\ref{cap:brown} y~\ref{cap:itocalculus}), aplicación a
Black-Scholes (Capítulo~\ref{cap:bs}) y conclusiones.

\medskip

\noindent\textbf{Capítulo~\ref{cap:discreto}: el modelo binomial.}
Plantea el problema de valoración en el contexto discreto del modelo
binomial multiperíodo. Articula las herramientas elementales de
probabilidad necesarias (filtraciones, esperanza condicional,
martingalas, integral estocástica discreta) y resuelve el problema
mediante una receta de tres pasos: cambio de medida, construcción de
la martingala asociada al \emph{payoff} descontado, y aplicación del
teorema de representación binomial para identificar el portfolio
replicante. El capítulo concluye con una breve introducción al contexto continuo, anticipando la convergencia del random walk binomial al
movimiento browniano.

\medskip

\noindent\textbf{Capítulo~\ref{cap:brown}: el movimiento browniano.}
Introduce el movimiento browniano como objeto matemático, establece su
existencia como teorema (sin demostrar, citando bibliografía) y
estudia las propiedades de sus trayectorias: continuidad, no
diferenciabilidad, variación total infinita y variación cuadrática
finita e igual a la longitud del intervalo. Esta última propiedad es
la pieza técnica que motivará el desarrollo del cálculo estocástico
del capítulo siguiente.

\medskip

\noindent\textbf{Capítulo~\ref{cap:itocalculus}: cálculo estocástico.}
Constituye el corazón técnico del trabajo. Construye la integral
estocástica de Itô primero sobre procesos simples y luego, por
densidad, sobre el espacio $\mathcal{H}^2$. Demuestra la fórmula de
Itô para funciones del browniano y enuncia su versión general. Articula
el cambio de medida en tiempo continuo mediante el teorema de
Girsanov, y enuncia el teorema de representación de martingalas
brownianas. Cada una de estas piezas es la traducción al contexto
continuo de un objeto introducido en el caso discreto en el
Capítulo~\ref{cap:discreto}.

\medskip

\noindent\textbf{Capítulo~\ref{cap:bs}: el modelo de Black-Scholes.}
Aplica las herramientas del Capítulo~\ref{cap:itocalculus} al problema
de valoración en el modelo de Black-Scholes, articulando la receta de
tres pasos en su versión continua. Obtiene la fórmula general del
precio de un derivado europeo y la particulariza al \emph{payoff} de
la \emph{call} europea, derivando explícitamente la fórmula clásica
de Black-Scholes con sus coeficientes $d_1$ y $d_2$. El capítulo
concluye derivando la EDP de Black-Scholes a partir de la fórmula
obtenida, mostrando la consistencia con la vía clásica.

\medskip

\noindent\textbf{Capítulo~\ref{cap:conclusiones}: conclusiones.}
Recapitula el trabajo realizado, expone la convergencia de la
fórmula binomial multiperíodo a la fórmula de Black-Scholes (como culminación de la relación entre ambos modelos en términos cuantitativos) y
discute brevemente las líneas de continuación naturales del trabajo.

% =====================================================================
%   CAPÍTULO 2 — VALORACIÓN DE OPCIONES EN TIEMPO DISCRETO
% =====================================================================

\chapter{Valoración de opciones en tiempo discreto}\label{cap:discreto}

\section{Conceptos financieros y principio de no-arbitraje}
\label{sec:conceptos_no_arbitraje}

\subsection{Definiciones básicas de finanzas}
\label{subsec:def_basicas}

A lo largo del trabajo nos situaremos en un mercado idealizado en el que
se negocian dos tipos básicos de instrumentos financieros: un activo
libre de riesgo y un activo con riesgo. Sobre ellos se construirán los
contratos derivados, cuyo valor depende en última instancia del
comportamiento futuro del activo subyacente.

\medskip

\noindent\textbf{El bono libre de riesgo.} Denotaremos con $B_t$ el precio
en el instante $t \geq 0$ del activo libre de riesgo, con la
normalización $B_0 = 1$. Este activo representa una inversión sin
incertidumbre (en la práctica, un depósito remunerado o un bono del Estado) cuyo valor crece de manera determinista a un tipo de interés
fijo. Antes de escribir su expresión, conviene detenerse en cómo se
modela ese crecimiento.

Supongamos que una unidad monetaria se coloca a un tipo de interés $r$
por unidad de tiempo y que los intereses se reinvierten (se
\emph{capitalizan}) $n$ veces en cada unidad de tiempo. Cada subperiodo,
de longitud $1/n$, devenga el tipo proporcional $r/n$; al reinvertirse
los intereses, el capital se multiplica en cada subperiodo por el factor
$1 + r/n$, de modo que, transcurrida una unidad de tiempo ($n$
subperiodos), la unidad monetaria inicial se habrá convertido en
\begin{equation*}
  \left(1 + \frac{r}{n}\right)^{n}
\end{equation*}
Esto es la \emph{capitalización compuesta} con frecuencia $n$.

La \emph{capitalización continua} surge al reinvertir los intereses de
forma instantánea, esto es, al hacer la frecuencia $n$ arbitrariamente
grande. Tomando el límite se obtiene la expresión clásica
\begin{equation}\label{eq:limite_e}
  \lim_{n \to \infty}\left(1 + \frac{r}{n}\right)^{n} = e^{r}
\end{equation}
El factor de
crecimiento por unidad de tiempo es, por tanto, $e^{r}$, y el
correspondiente a un horizonte $t > 0$ cualquiera se obtiene por
composición:
\begin{equation*}
  B_t = \left(e^{r}\right)^{t} = e^{rt}.
\end{equation*}
Adoptaremos, pues,
\begin{equation}\label{eq:bono}
  B_t = e^{rt}, \qquad r > 0, \qquad B_0 = 1,
\end{equation}
como modelo del activo libre de riesgo. La denominación \emph{libre de
riesgo} se justifica en que el rendimiento es conocido con certeza en el
instante inicial: invertir una unidad monetaria en $t = 0$ asegura
$e^{rT}$ unidades en $t = T$, sin que exista incertidumbre sobre el resultado.

La elección de la capitalización continua no supone pérdida de
generalidad frente a los regímenes discretos habituales: es solo una
forma cómoda de expresar el mismo fenómeno. En efecto, si un régimen
discreto ofrece un tipo efectivo $r'$ por unidad de tiempo, de modo
que una unidad monetaria se convierte en $1 + r'$ al cabo de una unidad
de tiempo, el tipo continuo $r$ que produce el mismo factor de
crecimiento es el que verifica $e^{r} = 1 + r'$, es decir,
\begin{equation*}
  r = \log(1 + r').
\end{equation*}
Por ejemplo, un tipo efectivo anual del $5\%$ ($r' = 0{,}05$) equivale a
un tipo continuo $r = \log(1{,}05) \approx 0{,}0488$, esto es, un
$4{,}88\%$. En lo sucesivo, $r$ denotará siempre el tipo en
capitalización continua.

\medskip

\noindent\textbf{La acción con riesgo.} En cambio, el precio del activo
con riesgo (al que llamaremos genéricamente \emph{acción}) se modela
como un proceso estocástico $\{S_t\}_{t \geq 0}$ definido sobre un
espacio de probabilidad adecuado. La dinámica concreta de $\{S_t\}$ se
irá especificando a lo largo del trabajo: en este capítulo tomará
valores en un conjunto finito en cada instante (modelo binomial),
mientras que en los Capítulos~\ref{cap:brown} y \ref{cap:bs}
consideraremos modelos en tiempo continuo.

\medskip

\noindent\textbf{Contratos derivados.} Un \emph{contrato derivado} es un
acuerdo entre dos partes cuyo \emph{payoff} (la cantidad que han de
intercambiar en una fecha futura) depende del precio del activo
subyacente. A lo largo del trabajo aparecerán recurrentemente los tres
derivados que recogen las siguientes definiciones, en las que $T > 0$
representa la fecha de vencimiento y $K > 0$ el precio acordado en el
contrato.

\begin{definicion}[Contrato a plazo o \emph{forward}]
Un contrato a plazo, o \emph{forward}, sobre la acción $S$ con
vencimiento $T$ y precio de entrega $K$ es un acuerdo por el que el
comprador está obligado a adquirir el activo $S$ en el instante $T$ al
precio $K$, fijado de antemano. El \emph{payoff} para el comprador en
$T$ es la variable aleatoria
\begin{equation*}
  S_T - K.
\end{equation*}
\end{definicion}

El \emph{forward} permite a las dos partes fijar hoy el precio de una
operación que se ejecutará en el futuro, eliminando la incertidumbre
sobre el precio al que se comprará o venderá el activo en $T$. Es, por
ello, una herramienta habitual de cobertura: quien vaya a necesitar el
activo más adelante se protege de una subida de su precio, y quien vaya a
desprenderse de él, de una bajada.

El \emph{forward} es un contrato \emph{de obligación}: ambas partes
deben cumplir lo pactado, de modo que el \emph{payoff} puede ser
negativo si $S_T < K$. La pregunta natural es qué valor de $K$ hace
que el contrato sea \emph{justo}, en el sentido de poder firmarse
sin pago en $t = 0$.

\medskip

Conviene, antes de definirlas, explicar para qué sirven las dos opciones
que siguen. A diferencia del \emph{forward}, que compromete a ambas
partes, una opción funciona en la práctica como un \emph{seguro} frente a
los movimientos del precio del subyacente. Una \emph{call} interesa a
quien prevé comprar el activo en el futuro y quiere protegerse de una
subida: le permite asegurarse hoy un precio de compra, aprovechando el
alza si se produce, pero sin quedar obligado si el precio baja. Por su parte, una
\emph{put} interesa a quien posee el activo o planea
venderlo y teme una caída: le garantiza un precio de venta, protegiendo
sus ganancias frente a las bajadas. Junto a este uso de cobertura, las
opciones se emplean también con fines especulativos, pues permiten
apostar por la dirección del precio con un desembolso inicial reducido.

\begin{definicion}[Opción \emph{call} europea]
Una \emph{opción call} europea sobre $S$ con vencimiento $T$ y precio
de ejercicio $K$ otorga al comprador el derecho (no la obligación) de
adquirir el activo en $T$ al precio $K$. Su \emph{payoff} en $T$ es la
variable aleatoria
\begin{equation*}
  (S_T - K)^+ \eqdef \max\{S_T - K,\, 0\}.
\end{equation*}
\end{definicion}

\begin{definicion}[Opción \emph{put} europea]
Una \emph{opción put} europea sobre $S$ con vencimiento $T$ y precio
de ejercicio $K$ otorga al comprador el derecho (no la obligación) de
vender el activo en $T$ al precio $K$. Su \emph{payoff} en $T$ es la
variable aleatoria
\begin{equation*}
  (K - S_T)^+ \eqdef \max\{K - S_T,\, 0\}.
\end{equation*}
\end{definicion}

A diferencia del \emph{forward}, la \emph{call} y la \emph{put} son
contratos \emph{de derecho}: el tenedor decide al vencimiento si
ejerce o no, eligiendo lo más favorable. Su \emph{payoff} es por
tanto no negativo, y a diferencia del \emph{forward}, cuyo precio
justo en $t = 0$ resulta ser nulo, una opción tiene siempre valor
estrictamente positivo al emitirse.

\medskip

\noindent\textbf{Vocabulario.} En la terminología estándar, el activo
$S$ se llama \emph{underlying}, el precio $K$ \emph{strike}, y $T$
\emph{maturity}. El calificativo \emph{europea} indica que el
ejercicio solo es posible en $T$, en contraste con las opciones
americanas (ejercibles en cualquier $t \in [0, T]$), que quedan
fuera del alcance de este trabajo.

\subsection{El principio de no-arbitraje}
\label{subsec:noarbitraje}

Antes de entrar en el modelo concreto que ocupará el resto del
capítulo, conviene detenerse en el principio sobre el que se sustenta
la teoría de valoración. Dicho principio, conocido como \emph{ausencia
de arbitraje} o \emph{principio de no-arbitraje}, no es un resultado
matemático que vayamos a demostrar, sino una hipótesis sobre el mercado
que asumiremos como punto de partida.

\medskip

\noindent\textbf{La idea informal.} Una \emph{oportunidad de arbitraje}
es, intuitivamente, una estrategia de inversión que, partiendo de
capital inicial nulo, garantiza al inversor un beneficio no negativo en
todo escenario posible y estrictamente positivo en al menos uno de
ellos. Dicho de otro modo: una posibilidad de obtener algo a partir de
nada, sin asumir riesgo alguno.

Asumiremos que el mercado no admite oportunidades de arbitraje. La
justificación es económica: cualquier oportunidad real sería
detectada y aprovechada de inmediato, haciéndola desaparecer; los
arbitrajes, si existen, son fenómenos transitorios. La formalización
matemática se diferirá hasta la Sección~\ref{sec:binario}, una vez
especificado el modelo.

\medskip

\noindent\textbf{El principio de replicación.} La hipótesis de no
arbitraje puede reformularse de manera más operativa. El siguiente
enunciado, aunque informal, recoge su idea fundamental.

\begin{quote}
\emph{Si dos carteras de activos producen exactamente los mismos pagos
futuros en todos los escenarios posibles, entonces deben tener el mismo
valor en cualquier instante anterior.}
\end{quote}

La justificación es inmediata: si en un instante $t < T$ una de las
carteras valiera más que la otra, vender la cara y comprar la barata
produciría una ganancia segura sin asumir riesgo, contradiciendo la
hipótesis. Llamaremos a esta consecuencia \emph{principio de
replicación}.

Este principio sugiere una estrategia general para valorar contratos
derivados: dada una opción, o, más en general, cualquier contrato cuyo
\emph{payoff} en $T$ sea una variable aleatoria $C_T$ conocida, basta
con construir una cartera (formada por activos cuyos precios sí
conozcamos) cuyo valor en $T$ coincida con $C_T$ en todos los
escenarios. Diremos entonces que la cartera \emph{replica} el contrato.
Por el principio de replicación, el precio justo del derivado en
cualquier instante anterior es necesariamente igual al valor de esta
cartera replicante. Esta es la idea central de toda la teoría que
desarrollaremos: tanto en el contexto discreto como en el continuo, el
problema de hallar el precio de un derivado se reducirá, en última
instancia, al problema de construir una cartera replicante.

\medskip

\noindent\textbf{Sobre el papel de las probabilidades.} Conviene hacer
una observación que irá cobrando sentido a lo largo del capítulo: el
argumento de replicación, tal como lo hemos planteado, no utiliza en
ningún momento las probabilidades asociadas a cada escenario futuro.
Solo intervienen los \emph{posibles} valores del activo, no la
probabilidad de que se produzcan. Como veremos, este hecho tiene
consecuencias profundas: el precio justo de un derivado es independiente
de las creencias subjetivas del mercado sobre el futuro. Es una de las
ideas más sorprendentes y bonitas de la teoría.

\subsection{Valoración del forward como caso de calentamiento}
\label{subsec:forward}

Como primera aplicación del principio de replicación, abordamos el
problema más sencillo posible: la valoración de un \emph{forward}. A lo
largo de la subsección asumiremos que los dos activos del mercado son
negociables sin restricciones en cualquier instante $t \in [0, T]$, sin
costes de transacción y en cantidades fraccionarias arbitrarias. Estas
hipótesis son idealizaciones, pero serán suficientes para que los
argumentos que siguen tengan validez.

\medskip

\noindent\textbf{El precio justo de un \emph{forward} con \emph{strike}
$K$.} Consideremos un \emph{forward} sobre la acción $S$ con
vencimiento $T$ y precio de entrega $K$. Su \emph{payoff} para el
comprador en $T$ es $S_T - K$. Buscamos el precio justo $V_0$ del
contrato en el instante inicial: la cantidad que el comprador debe
abonar al emisor en $t = 0$ a cambio de recibir $S_T - K$ en $t = T$.

Consideremos la siguiente cartera, constituida en $t = 0$:
\begin{itemize}
  \item se compra una unidad de la acción, al precio $S_0$;
  \item se toma prestada (es decir, se emite) una cantidad $K e^{-rT}$
        del activo libre de riesgo.
\end{itemize}
El coste neto de constituirla en $t = 0$ es $\Pi_0 = S_0 - K e^{-rT}$.
En $t = T$, la unidad de acción habrá pasado a valer $S_T$, mientras
que la deuda contraída habrá crecido hasta $K e^{-rT} \cdot e^{rT} = K$.
Por consiguiente, el valor de la cartera en $t = T$ es
\begin{equation*}
  \Pi_T = S_T - K,
\end{equation*}
que coincide con el \emph{payoff} del \emph{forward}. Por el principio
de replicación,
\begin{equation}\label{eq:precio_forward}
  V_0 = S_0 - K e^{-rT}.
\end{equation}

Conviene verificar que el argumento es consistente, en el sentido de que
cualquier desviación del precio \eqref{eq:precio_forward} produciría una
oportunidad de arbitraje. Si el contrato se vendiera a $V_0' > S_0 - K
e^{-rT}$, un inversor podría: en $t = 0$, vender el \emph{forward}
(recibiendo $V_0'$), constituir la cartera replicante (a un coste $S_0
- K e^{-rT}$) y guardar el remanente $V_0' - (S_0 - K e^{-rT}) > 0$. En
$t = T$, los pagos asociados al \emph{forward} y a la cartera se
cancelan, de modo que el inversor habría obtenido un beneficio cierto
sin asumir riesgo. La situación simétrica $V_0' < S_0 - K e^{-rT}$
conduce a un arbitraje análogo.

\medskip

\noindent\textbf{El strike justo y el precio \emph{forward}.} En la
práctica, los \emph{forwards} se firman habitualmente con $V_0 = 0$.
Imponer esta condición en \eqref{eq:precio_forward} permite despejar el
precio de entrega que hace el contrato justo:
\begin{equation}\label{eq:precio_forward_just}
  K = S_0 e^{rT}.
\end{equation}
Esta cantidad recibe el nombre de \emph{precio forward} de la acción
para el vencimiento $T$, y se interpreta como el valor en $t = T$ que
tendría, por capitalización al tipo libre de riesgo, una unidad de
acción comprada hoy.

\medskip

\noindent\textbf{Dos lecciones generales.}

\begin{observacion}[Independencia de la dinámica del subyacente]
La derivación del precio del \emph{forward} no ha hecho uso alguno del
modelo concreto que rige la evolución del precio de la acción: no hemos
supuesto ni que $\{S_t\}$ sigue un proceso particular, ni siquiera que
toma valores en un conjunto específico. La fórmula
\eqref{eq:precio_forward} es válida en cualquier mercado libre de
arbitraje en el que el bono y la acción sean negociables sin
restricciones. La razón última es que el \emph{payoff} del
\emph{forward} depende del subyacente de manera \emph{lineal}, y por
ello la replicación puede llevarse a cabo mediante una estrategia
estática. Esta universalidad es, sin embargo, una excepción: en cuanto
el \emph{payoff} del derivado es no lineal en $S_T$ (como ocurre con la
\emph{call} y la \emph{put}) el argumento de replicación deja de ser
estático y pasa a requerir un modelo concreto para la dinámica del
subyacente.
\end{observacion}

\begin{observacion}[Irrelevancia de las probabilidades subjetivas]
En la fórmula \eqref{eq:precio_forward} no aparece ninguna probabilidad
asociada a la evolución futura de la acción. El precio justo del
\emph{forward} no depende de cuán probable sea, según las creencias del
mercado, que la acción suba o baje, ni de su rendimiento medio
esperado, ni de su volatilidad. Esta independencia respecto de las
probabilidades subjetivas reaparecerá, en una forma más sutil, en los
modelos no triviales que estudiaremos a continuación.
\end{observacion}

Tomando este ejemplo inicial como referencia, podemos abordar el
problema general en su versión más simple, aunque no por ello trivial.

\section{El modelo binario uniperíodo}
\label{sec:binario}

El modelo más simple que conserva los ingredientes esenciales del
problema es el \emph{binario uniperíodo}: el mercado se observa solo
en $t = 0$ y $t = T$, y el precio del subyacente toma solo dos
valores posibles en $T$.

\subsection{Descripción del modelo y formalización del no-arbitraje}
\label{subsec:setup_binario}

\noindent\textbf{Descripción del modelo.} El mercado consta de un bono
libre de riesgo, con $B_t = e^{rt}$ y $r > 0$, y una acción cuyo precio
inicial $S_0 > 0$ es conocido y cuyo precio final $S_T$ es una variable
aleatoria que toma únicamente dos valores. Concretamente, fijamos dos
constantes $u, d$ con $0 < d < u$ y suponemos que
\begin{equation*}
  S_T = \begin{cases}
    S_0\, u & \text{con probabilidad } p, \\[2pt]
    S_0\, d & \text{con probabilidad } 1 - p,
  \end{cases}
\end{equation*}
para cierto $p \in (0, 1)$. Llamaremos a $u$ y $d$ los \emph{factores
de subida} y \emph{de bajada}, y a $p$ la \emph{probabilidad subjetiva}
de subida.

Para fijar ideas, formulamos el modelo en el lenguaje
de la teoría de la probabilidad. Sea $\Omega = \{\omega_u, \omega_d\}$
el espacio muestral, $\F = 2^{\Omega}$ y $\PP$ la probabilidad sobre
$(\Omega, \F)$ definida por $\PP(\{\omega_u\}) = p$ y $\PP(\{\omega_d\})
= 1 - p$. Sobre $(\Omega, \F, \PP)$, el precio final de la acción es la
variable aleatoria
\begin{equation*}
  S_T(\omega_u) = S_0\, u, \qquad S_T(\omega_d) = S_0\, d.
\end{equation*}
Llamaremos a $\PP$ la \emph{medida de mercado} o \emph{medida física}.

Aunque incluyamos $p$ en la descripción, su papel es solo especificar
qué escenarios consideramos \emph{posibles} (los de probabilidad
estrictamente positiva). El valor cuantitativo de $p$ no intervendrá
en ningún cálculo.

\medskip

\noindent\textbf{Formalización del concepto de arbitraje.}

\begin{definicion}[Estrategia o cartera]
Una \emph{estrategia} en el modelo binario uniperíodo es un par
$(\phi, \psi) \in \R^2$, donde $\phi$ representa el número de unidades
de la acción y $\psi$ el número de unidades del bono que se mantienen
en cartera durante el período $[0, T]$.
\end{definicion}

El valor de la estrategia en cada instante es
\begin{align*}
  V_0(\phi, \psi)   &= \phi\, S_0 + \psi, \\
  V_T(\phi, \psi)   &= \phi\, S_T + \psi\, e^{rT}.
\end{align*}
Permitimos que $\phi$ y $\psi$ tomen cualquier valor real, incluidos
negativos: $\phi < 0$ corresponde a \emph{ir corto} en la acción, y
$\psi < 0$ a tomar prestada del banco una cantidad $|\psi|$. Esta
posibilidad es esencial para la teoría: sin ella, no toda opción podría
ser replicada.

\begin{definicion}[Oportunidad de arbitraje]
\label{def:arbitraje}
Una estrategia $(\phi, \psi)$ es una \emph{oportunidad de arbitraje} si:
\begin{enumerate}[label=\textnormal{(\roman*)}]
  \item $V_0(\phi, \psi) = 0$;
  \item $V_T(\phi, \psi) \geq 0$ casi seguramente bajo $\PP$;
  \item $\PP\bigl(\{V_T(\phi, \psi) > 0\}\bigr) > 0$.
\end{enumerate}
Diremos que el modelo es \emph{libre de arbitraje} si no admite ninguna
estrategia de este tipo.
\end{definicion}

\medskip

\begin{teorema}[Caracterización del no arbitraje en el modelo binario]
\label{teo:noarb_binario}
El modelo binario uniperíodo es libre de arbitraje si y solo si
\begin{equation}\label{eq:cond_noarb_binario}
  d < e^{rT} < u.
\end{equation}
\end{teorema}

\begin{proof}
$(\Rightarrow)$ Procederemos por contrarrecíproco: si la
condición~\eqref{eq:cond_noarb_binario} falla, exhibiremos una
oportunidad de arbitraje.

Supongamos $e^{rT} \leq d$. Consideremos la estrategia $(\phi, \psi) =
(1, -S_0)$, consistente en comprar una unidad de acción financiándose
con bono. Su valor inicial es $V_0 = S_0 - S_0 = 0$, y su valor final
es
\begin{equation*}
  V_T = S_T - S_0\, e^{rT} = S_0\bigl(S_T/S_0 - e^{rT}\bigr),
\end{equation*}
que vale $S_0(u - e^{rT}) > 0$ en $\omega_u$ y $S_0(d - e^{rT}) \geq 0$
en $\omega_d$. La estrategia constituye una oportunidad de arbitraje.
Análogamente, si $e^{rT} \geq u$, la estrategia $(\phi, \psi) = (-1,
S_0)$ produce un arbitraje.

$(\Leftarrow)$ Supongamos~\eqref{eq:cond_noarb_binario}. Sea $(\phi,
\psi)$ con $V_0 = 0$, esto es, $\psi = -\phi\, S_0$. Su valor final es
\begin{align*}
  V_T(\omega_u) &= \phi\, S_0\bigl(u - e^{rT}\bigr), \\
  V_T(\omega_d) &= \phi\, S_0\bigl(d - e^{rT}\bigr).
\end{align*}
Por~\eqref{eq:cond_noarb_binario}, los factores $u - e^{rT}$ y $d -
e^{rT}$ tienen signos opuestos. El signo de $V_T$ es por tanto opuesto
en los dos estados, salvo en el caso trivial $\phi = 0$ que implica
$V_T \equiv 0$. No existe estrategia $(\phi, \psi) \neq (0, 0)$ con
$V_0 = 0$, $V_T \geq 0$ y $V_T > 0$ en al menos un estado.
\end{proof}

La condición~\eqref{eq:cond_noarb_binario} admite una interpretación
económica natural: el rendimiento de la acción debe poder superar al
del bono libre de riesgo (de ahí $u > e^{rT}$) pero también debe poder
quedar por debajo (de ahí $d < e^{rT}$). En otro caso surgirían
arbitrajes triviales. La hipótesis exige que la acción sea genuinamente
\emph{arriesgada} respecto al bono. La supondremos en lo sucesivo.

\subsection{Construcción del portfolio replicante}
\label{subsec:replicante_binario}

Consideremos un derivado con \emph{payoff} $C_T$ definido sobre
$(\Omega, \F, \PP)$. Por construcción, $C_T$ toma a lo sumo dos
valores:
\begin{equation*}
  C_u \eqdef C_T(\omega_u), \qquad C_d \eqdef C_T(\omega_d).
\end{equation*}
Buscamos una estrategia $(\phi, \psi) \in \R^2$ tal que
\begin{equation}\label{eq:replicante_binario}
  V_T(\phi, \psi)(\omega) = C_T(\omega) \quad \text{para todo } \omega \in \Omega.
\end{equation}
Una estrategia que satisfaga \eqref{eq:replicante_binario} se denomina
\emph{cartera replicante} del derivado.

La condición~\eqref{eq:replicante_binario}, escrita explícitamente en
cada estado, se traduce en el sistema lineal
\begin{equation}\label{eq:sistema_replicante}
  \left\{
  \begin{aligned}
    \phi\, S_0\, u + \psi\, e^{rT} &= C_u, \\
    \phi\, S_0\, d + \psi\, e^{rT} &= C_d.
  \end{aligned}
  \right.
\end{equation}
La matriz de coeficientes tiene determinante $S_0\, e^{rT}\, (u - d) >
0$, y por tanto el sistema admite una única solución:
\begin{equation}\label{eq:phi_psi}
  \phi = \frac{C_u - C_d}{S_0\,(u - d)}, \qquad
  \psi = e^{-rT}\, \frac{u\, C_d - d\, C_u}{u - d}.
\end{equation}

\begin{teorema}[Replicación en el modelo binario uniperíodo]
\label{teo:replicacion_binario}
En el modelo binario uniperíodo libre de arbitraje, todo derivado con
\emph{payoff} $C_T$ admite una única cartera replicante $(\phi, \psi)
\in \R^2$, dada por~\eqref{eq:phi_psi}. El precio justo del derivado en
$t = 0$ es
\begin{equation}\label{eq:precio_binario}
  V_0 = \phi\, S_0 + \psi
       = \frac{C_u - C_d}{u - d}
       + e^{-rT}\, \frac{u\, C_d - d\, C_u}{u - d}.
\end{equation}
\end{teorema}

\begin{proof}
La existencia y unicidad de $(\phi, \psi)$ se han justificado en la
discusión anterior. Para ver que $V_0 = \phi\, S_0 + \psi$ es el
precio justo, supongamos que el derivado se negocia a $V_0' > V_0$.
Vender el derivado y constituir la cartera replicante requiere una
inversión inicial $V_0 - V_0' < 0$, dejando un remanente positivo. En
$t = T$, el \emph{payoff} del derivado y el valor de la cartera se
cancelan, produciendo un beneficio seguro. Análogamente con $V_0' <
V_0$. Ambas situaciones contradicen la hipótesis de no arbitraje.
\end{proof}

\medskip

\noindent\textbf{Dos observaciones.} La fórmula
\eqref{eq:precio_binario} no contiene la probabilidad subjetiva $p$:
el precio depende solo de $S_0, u, d, r, T, C_u, C_d$. Por otra
parte, la cantidad $\phi$ obtenida en~\eqref{eq:phi_psi} es el
cociente entre la variación del \emph{payoff} y la del precio del
subyacente: mide la sensibilidad del derivado a movimientos del
subyacente y recibe en la literatura financiera el nombre de
\emph{delta}.

\subsection{Aparición de la medida neutral al riesgo}
\label{subsec:medida_riesgo}

La fórmula~\eqref{eq:precio_binario} admite una manipulación algebraica
elemental que hace emerger un objeto nuevo (y central) de la teoría.
Reagrupando los términos en $C_u$ y $C_d$,
\begin{equation*}
  V_0 = e^{-rT}\!\left[\frac{e^{rT} - d}{u - d}\, C_u
                       + \frac{u - e^{rT}}{u - d}\, C_d\right].
\end{equation*}
Definiendo
\begin{equation}\label{eq:def_q}
  q \eqdef \frac{e^{rT} - d}{u - d},
\end{equation}
y observando que $1 - q = (u - e^{rT})/(u - d)$, la fórmula adopta la
forma compacta
\begin{equation}\label{eq:precio_binario_q}
  V_0 = e^{-rT}\,\bigl[\,q\, C_u + (1 - q)\, C_d\,\bigr].
\end{equation}

\medskip

\noindent\textbf{Propiedades del cociente $q$.} La condición de no
arbitraje $d < e^{rT} < u$ implica que tanto $e^{rT} - d > 0$ como $u -
d > 0$, y que $u - e^{rT} > 0$. Por consiguiente, $q \in (0, 1)$. La
afirmación recíproca también se verifica con facilidad. Esta
equivalencia entre $q \in (0,1)$ y la condición de no
arbitraje~\eqref{eq:cond_noarb_binario} es importante: lo que el
hecho $q \in (0,1)$ implica es que los coeficientes que aparecen en la
fórmula~\eqref{eq:precio_binario_q} son números no negativos que suman
uno, es decir, \emph{son las probabilidades de una medida sobre $\Omega$}.

\medskip

\noindent\textbf{Definición de la medida neutral al riesgo.} La
observación anterior justifica la introducción, sobre el espacio
$(\Omega, \F)$, de una nueva medida de probabilidad
\begin{equation}\label{eq:def_Q}
  \Q(\{\omega_u\}) = q, \qquad \Q(\{\omega_d\}) = 1 - q.
\end{equation}

\begin{definicion}[Medida neutral al riesgo]
\label{def:medida_riesgo}
La probabilidad $\Q$ definida en~\eqref{eq:def_Q} se denomina
\emph{medida neutral al riesgo} (o \emph{medida martingala
equivalente}) del modelo binario uniperíodo.
\end{definicion}

Adviértase que $\Q$ y $\PP$ asignan probabilidad estrictamente positiva
a los mismos sucesos. En el lenguaje de la teoría de la medida, son
\emph{equivalentes}, y se denota $\Q \sim \PP$. Esta equivalencia,
trivial ahora, será una condición delicada en los modelos continuos del
Capítulo~\ref{cap:bs}.

\medskip

\noindent\textbf{La propiedad fundamental de $\Q$.} Calculemos la
esperanza (bajo $\Q$) del precio descontado de la acción en $t = T$:
\begin{align*}
  \E_{\Q}\bigl[e^{-rT}\, S_T\bigr]
  &= e^{-rT}\, S_0\,\bigl[q\, u + (1 - q)\, d\bigr].
\end{align*}
Sustituyendo $q = (e^{rT} - d)/(u - d)$, un cálculo directo muestra que
$q\, u + (1 - q)\, d = e^{rT}$, y por tanto
\begin{equation}\label{eq:martingala_uniperiodo}
  \E_{\Q}\bigl[e^{-rT}\, S_T\bigr] = S_0.
\end{equation}

La identidad~\eqref{eq:martingala_uniperiodo} contiene una idea central. Si introducimos la notación $\widetilde
S_t \eqdef e^{-rt} S_t$ para el precio descontado, se reescribe
$\E_{\Q}[\widetilde S_T] = \widetilde S_0$, propiedad que en la
Sección~\ref{subsec:martingalas_discretas} reconoceremos como la
versión preliminar de la propiedad de \emph{martingala}. Por este
motivo $\Q$ recibe también el nombre de medida martingala equivalente.

En cuanto al adjetivo \emph{neutral al riesgo}, su origen es el
siguiente. Reescribiendo \eqref{eq:martingala_uniperiodo} sin
descuento, $\E_\Q[S_T] = S_0\, e^{rT}$: bajo $\Q$, el rendimiento medio
de la acción coincide con el del bono libre de riesgo. Es decir, $\Q$
es la medida bajo la cual a un inversor le sería indiferente mantener
el activo arriesgado o el activo libre de riesgo, pues ambos le
ofrecerían el mismo rendimiento esperado. En la realidad (bajo $\PP$)
los inversores son aversos al riesgo y exigen prima por mantener
activos arriesgados; la medida $\Q$ elimina dicha prima por
construcción.

\medskip

\noindent\textbf{La fórmula del precio como esperanza descontada.} Por
la definición~\eqref{eq:def_Q} de $\Q$, la expresión $q\, C_u + (1-q)\,
C_d$ es exactamente la esperanza de $C_T$ bajo $\Q$. La fórmula del
precio adopta entonces la forma definitiva
\begin{equation}\label{eq:precio_esperanza}
  \boxed{\,V_0 = e^{-rT}\, \E_{\Q}[C_T].\,}
\end{equation}
El precio justo de un derivado en $t = 0$ es el valor descontado de la
esperanza de su \emph{payoff} bajo la medida neutral al riesgo $\Q$. La
fórmula reduce el problema de valoración a un cálculo de esperanzas,
una vez identificada la medida apropiada.

Insistamos en una característica notable de
\eqref{eq:precio_esperanza}: hemos llegado a ella sin haber utilizado
en ningún momento la probabilidad subjetiva $p$.

\subsection{Independencia de las probabilidades subjetivas}
\label{subsec:indep_subjetiva}

\noindent\textbf{Una intuición que conviene corregir.} Para un lector
con formación en probabilidad, la primera respuesta natural a la
pregunta ``¿cuál es el precio justo de un derivado con \emph{payoff}
aleatorio $C_T$?'' sería: \emph{el valor esperado descontado de su
\emph{payoff}}, esto es,
\begin{equation}\label{eq:precio_ingenuo}
  V_0^{\text{ing}} \eqdef e^{-rT}\, \E_{\PP}[C_T]
                  = e^{-rT}\,\bigl[p\, C_u + (1-p)\, C_d\bigr].
\end{equation}
Sin embargo, basta comparar~\eqref{eq:precio_ingenuo} con el precio
correcto~\eqref{eq:precio_esperanza} para advertir que ambos coinciden
si y solo si $p = q$. Salvo en este caso fortuito, los dos precios son
distintos.

¿Qué tiene de erróneo el razonamiento intuitivo? La
fórmula~\eqref{eq:precio_ingenuo} solo sería razonable si los
inversores fueran indiferentes al riesgo, en cuyo caso evaluarían
cualquier activo por su valor esperado descontado. En la práctica, los
inversores son aversos al riesgo y exigen primas. Si se aplicara
\eqref{eq:precio_ingenuo} para fijar el precio del derivado, surgirían
oportunidades de arbitraje\footnote{Si $p > q$, entonces
$V_0^{\text{ing}} > V_0$ (en el caso $C_u > C_d$), y un inversor
podría vender el derivado al precio $V_0^{\text{ing}}$, constituir la
cartera replicante a coste $V_0$, y embolsarse la diferencia sin asumir
riesgo. La situación $p < q$ produce un arbitraje simétrico.}.

\medskip

\noindent\textbf{Valoración por replicación, no por esperanza.} El
argumento que nos llevó a~\eqref{eq:precio_esperanza} no descansa en
ningún cálculo de esperanzas bajo la medida del mercado, sino en la
construcción explícita de una cartera replicante. Lo que determina el
precio del derivado es la posibilidad de \emph{reproducir} su
\emph{payoff} mediante una combinación de bonos y acciones; la
probabilidad subjetiva $p$ no interviene porque la cartera replicante
reproduce el \emph{payoff} en \emph{todos los escenarios posibles}, sin
necesidad de ponderarlos por sus probabilidades.

\medskip

\noindent\textbf{El papel de las probabilidades en la teoría.} Lo
anterior podría sugerir que las probabilidades desaparecen del problema
de valoración, pero no es así: las probabilidades \emph{relevantes} no
son las del mercado, sino las de la medida neutral al riesgo $\Q$. Bajo
$\Q$, el precio del derivado se recupera como esperanza descontada del
\emph{payoff}. La intuición probabilística inicial no era errónea en su
forma (el precio sí es una esperanza descontada) sino en la elección de
la medida bajo la cual debe tomarse.

Esta sustitución de $\PP$ por $\Q$ es la idea que vertebrará el
desarrollo posterior, tanto en el modelo binomial multiperíodo como
en Black-Scholes.

\subsection{El modelo ternario y la incompletitud}
\label{subsec:ternario}

La generalización natural del modelo binario consiste en permitir
más valores del subyacente. Estudiamos brevemente el caso ternario
para mostrar que la replicación se rompe, fenómeno que motiva la
noción de \emph{completitud} de un mercado.

\medskip

Mantenemos los dos activos del modelo y suponemos que en $t = T$ el
precio toma valores en $\{S_0\, u_1, S_0\, u_2, S_0\, u_3\}$ con $0 <
u_3 < u_2 < u_1$, todos con probabilidad estrictamente positiva.
Asumimos que el modelo es libre de arbitraje (lo cual requiere $u_3 <
e^{rT} < u_1$).

Dado un derivado con \emph{payoff} $C_T$ de valores $C_i \eqdef
C_T(\omega_i)$ para $i = 1, 2, 3$, la condición de replicación
$V_T(\phi, \psi) = C_T$ se traduce en un sistema lineal de \emph{tres}
ecuaciones con dos incógnitas:
\begin{equation*}
  \phi\, S_0\, u_i + \psi\, e^{rT} = C_i, \quad i = 1, 2, 3.
\end{equation*}
Para un \emph{payoff} arbitrario, el vector $(C_1, C_2, C_3)^\top$ no
pertenece en general al subespacio bidimensional de $\R^3$ generado por
las columnas de la matriz del sistema, y por tanto este no admite
solución. La conclusión es directa: en el modelo ternario, el método de
replicación falla, y el principio de no arbitraje, por sí solo, no
proporciona un precio único.

\begin{definicion}[Mercado completo, versión preliminar]
\label{def:completitud_preliminar}
Un mercado se dice \emph{completo} si todo derivado sobre los activos
del mercado admite una cartera replicante.
\end{definicion}

En un mercado incompleto, el principio de no arbitraje sigue
restringiendo el precio de un derivado, pero no de manera unívoca: deja
un \emph{intervalo} de precios libres de arbitraje, en general no
degenerado. Para fijar un precio dentro del intervalo es necesario
invocar consideraciones adicionales (preferencias del inversor,
modelos de equilibrio, criterios de utilidad) que escapan del puro
principio de no arbitraje.

A lo largo del trabajo solo abordaremos modelos completos. Tanto el
binomial multiperíodo (Sección~\ref{sec:multiperiodo}) como el de
Black-Scholes (Capítulo~\ref{cap:bs}) lo serán; en el contexto
continuo, la completitud quedará garantizada por el teorema de
representación de martingalas brownianas
(Capítulo~\ref{cap:itocalculus}).

\section{Modelo binomial multiperíodo y maquinaria probabilística}
\label{sec:multiperiodo}

El modelo binario uniperíodo es insuficiente como modelo realista de
mercado. Su generalización natural, el modelo \emph{binomial
multiperíodo}, permite observar el mercado en varios instantes
intermedios. Su tratamiento exige introducir filtraciones, esperanzas
condicionales, martingalas y procesos predecibles, conceptos que
formalizan ideas financieras naturales del modelo.

\subsection{Descripción del modelo binomial multiperíodo}
\label{subsec:setup_multiperiodo}

\noindent\textbf{Discretización del horizonte temporal.} Fijemos $T >
0$ y $N \geq 1$ entero. Subdividimos $[0, T]$ en $N$ subintervalos de
longitud $\delta t \eqdef T/N$, y consideramos los instantes $t_k \eqdef
k\, \delta t$ para $k = 0, 1, \ldots, N$. El mercado se observa
únicamente en estos instantes. El bono evoluciona de manera
determinista: $B_k \eqdef e^{r\, t_k}$.

\medskip

\noindent\textbf{La dinámica del subyacente.} La acción tiene precio
inicial $S_0 > 0$ conocido. Fijamos constantes $u, d$ con $0 < d < u$ y
suponemos que entre dos instantes consecutivos el precio se multiplica
por $u$ con probabilidad $p$ o por $d$ con probabilidad $1-p$, e
independientemente del comportamiento en los pasos anteriores:
\begin{equation*}
  S_{k+1} = \begin{cases}
    S_k\, u & \text{con probabilidad } p, \\[2pt]
    S_k\, d & \text{con probabilidad } 1 - p,
  \end{cases}
\end{equation*}
La estructura se visualiza con un árbol binomial recombinante
(Figura~\ref{fig:arbol_binomial}): en el instante $t_k$, el precio
puede tomar $k+1$ valores distintos, $S_0\, u^j\, d^{k-j}$ para $j = 0,
\ldots, k$.

\begin{figure}[h]
\centering
\begin{tikzpicture}[
  scale=1,
  every node/.style={font=\small},
  edge/.style={->, >=stealth, thick}
]
  \node (S0) at (0,0) {$S_0$};
  \node (S1u) at (2,1) {$S_0 u$};
  \node (S1d) at (2,-1) {$S_0 d$};
  \node (S2uu) at (4,2) {$S_0 u^2$};
  \node (S2ud) at (4,0) {$S_0 ud$};
  \node (S2dd) at (4,-2) {$S_0 d^2$};
  \node (S3uuu) at (6,3) {$S_0 u^3$};
  \node (S3uud) at (6,1) {$S_0 u^2 d$};
  \node (S3udd) at (6,-1) {$S_0 u d^2$};
  \node (S3ddd) at (6,-3) {$S_0 d^3$};
  \draw[edge] (S0) -- (S1u);
  \draw[edge] (S0) -- (S1d);
  \draw[edge] (S1u) -- (S2uu);
  \draw[edge] (S1u) -- (S2ud);
  \draw[edge] (S1d) -- (S2ud);
  \draw[edge] (S1d) -- (S2dd);
  \draw[edge] (S2uu) -- (S3uuu);
  \draw[edge] (S2uu) -- (S3uud);
  \draw[edge] (S2ud) -- (S3uud);
  \draw[edge] (S2ud) -- (S3udd);
  \draw[edge] (S2dd) -- (S3udd);
  \draw[edge] (S2dd) -- (S3ddd);
  \node at (0,-3.7) {$t_0$};
  \node at (2,-3.7) {$t_1$};
  \node at (4,-3.7) {$t_2$};
  \node at (6,-3.7) {$t_3$};
\end{tikzpicture}
\caption{Árbol binomial multiperíodo con $N = 3$ pasos. La estructura
recombinante (las trayectorias $ud$ y $du$ confluyen en el mismo nodo
$S_0\,ud$) reduce el número de nodos de $2^k$ a $k+1$.}
\label{fig:arbol_binomial}
\end{figure}

\medskip

\noindent\textbf{El espacio probabilístico.} Tomamos como espacio
muestral el conjunto de todas las trayectorias posibles:
\begin{equation*}
  \Omega \eqdef \{u, d\}^N,
\end{equation*}
donde cada $\omega = (\omega_1, \ldots, \omega_N)$ con $\omega_j \in
\{u, d\}$ representa una historia completa del mercado. Tomamos $\F
\eqdef 2^{\Omega}$ y, sobre $(\Omega, \F)$, definimos la medida de
mercado $\PP$ mediante
\begin{equation*}
  \PP(\{\omega\}) \eqdef p^{\#\{j: \omega_j = u\}}\,
                        (1-p)^{\#\{j: \omega_j = d\}},
\end{equation*}
que codifica la independencia de los pasos. El precio en el instante
$t_k$ es la variable aleatoria $S_k(\omega) \eqdef S_0\, \prod_{j=1}^k
f(\omega_j)$, donde $f(u) = u$ y $f(d) = d$.

\subsection{Espacios de probabilidad filtrados}
\label{subsec:filtraciones}

Asumiremos en lo que sigue familiaridad con las nociones básicas de
probabilidad (espacio medible, $\sigma$-álgebra, variable aleatoria,
esperanza) estudiadas en cualquier grado \emph{STEM}.

\medskip

\noindent\textbf{$\sigma$-álgebras como información.} La idea es que una $\sigma$-álgebra $\G \subseteq \F$ sobre $(\Omega, \F)$
puede interpretarse como el conjunto de aquellos sucesos cuya
ocurrencia un observador es capaz de determinar: si $A \in \G$, el
observador puede contestar a la pregunta ``¿ha ocurrido $A$?''. La
$\sigma$-álgebra completa $\F$ corresponde al observador omnisciente;
sub-$\sigma$-álgebras $\G \subseteq \F$ corresponden a observadores con
información parcial.

\medskip

\noindent\textbf{La $\sigma$-álgebra $\F_k$ del modelo binomial.}
Ilustremos esta lectura sobre el modelo binomial multiperíodo. Para
cada $k \in \{0, 1, \ldots, N\}$, definimos $\F_k$ como la
$\sigma$-álgebra generada por la partición de $\Omega$ que agrupa las
trayectorias que coinciden en sus primeras $k$ coordenadas:
\begin{equation*}
  A_{(\omega_1, \ldots, \omega_k)} \eqdef
  \bigl\{\omega' \in \Omega : \omega'_j = \omega_j,\ j = 1, \ldots, k\bigr\}.
\end{equation*}
Hay $2^k$ de tales bloques (los \emph{átomos} de $\F_k$). En los
extremos, $\F_0 = \{\emptyset, \Omega\}$ es la $\sigma$-álgebra trivial
(sin haber observado nada, no se distingue ningún suceso del mundo
entero) y $\F_N = \F$ (tras observar la trayectoria completa, se conoce
todo). Sobre el árbol de la Figura~\ref{fig:arbol_binomial}, $A \in
\F_k$ si y solo si determinamos $\omega \in A$ a partir de los primeros
$k$ pasos.

La sucesión $(\F_k)_{k=0}^{N}$ es \emph{creciente}: $\F_k \subseteq
\F_{k+1}$ para todo $k$. Esto motiva la siguiente definición.

\begin{definicion}[Filtración]
\label{def:filtracion}
Sea $(\Omega, \F, \PP)$ un espacio de probabilidad. Una
\emph{filtración} sobre este espacio es una sucesión $(\F_k)_{k=0}^{N}$
de sub-$\sigma$-álgebras de $\F$ con $\F_k \subseteq \F_{k+1} \subseteq
\F$ para todo $k$. La cuádrupla $(\Omega, \F, (\F_k), \PP)$ se denomina
\emph{espacio de probabilidad filtrado}.
\end{definicion}

La filtración del modelo binomial es, desde el punto de vista
financiero, la \emph{filtración del mercado}.

\medskip

\begin{definicion}[Proceso adaptado]
\label{def:adaptado}
Un \emph{proceso estocástico} es una sucesión $(X_k)_{k=0}^{N}$ de
variables aleatorias sobre $(\Omega, \F)$. El proceso es
\emph{adaptado} a la filtración $(\F_k)$ si, para cada $k$, $X_k$ es
$\F_k$-medible.
\end{definicion}

La medibilidad respecto a $\F_k$ formaliza que el valor de $X_k$ es
conocido tras observar lo ocurrido hasta el instante $t_k$. En el
modelo binomial, $S_k$ es por construcción función de las primeras $k$
coordenadas y por tanto $\F_k$-medible: el proceso del precio es
adaptado.

\medskip

\begin{definicion}[Filtración natural]
\label{def:filtracion_natural}
La \emph{filtración natural} de un proceso $(X_k)_{k=0}^{N}$ es
$(\F_k^X)_{k=0}^{N}$ con $\F_k^X \eqdef \sigma(X_0, X_1, \ldots, X_k)$,
la $\sigma$-álgebra más pequeña que hace medibles las variables $X_0,
\ldots, X_k$.
\end{definicion}

Es la menor filtración respecto a la cual $(X_k)$ es adaptado. En el
modelo binomial multiperíodo, la filtración del mercado coincide con la
filtración natural del precio, y ambas serán identificadas y denotadas
simplemente por $(\F_k)$.

\subsection{Esperanza condicional}
\label{subsec:esperanza_condicional}

\noindent\textbf{Punto de partida.} El lector recordará la noción
elemental de esperanza condicional respecto a un suceso $B \in \F$ con
$\PP(B) > 0$: el número $\E[X \mid B] \eqdef \E[X\, \1_B]/\PP(B)$, que
representa el valor medio de $X$ en aquellos escenarios en los que $B$
ha ocurrido. La noción que necesitaremos es de naturaleza distinta: no
queremos condicionar respecto a un suceso individual, sino respecto a
toda una $\sigma$-álgebra $\G \subseteq \F$, y el resultado dejará de
ser un número para convertirse en una variable aleatoria.

\medskip

\noindent\textbf{Motivación sobre el árbol.} Sea $X$ una variable
aleatoria $\F_3$-medible sobre el árbol binomial con $N = 3$ pasos, y
preguntémonos por su valor medio ``dada la información $\F_2$''.
Intuitivamente, situados en el instante $t_2$ el observador ha llegado
a uno de los $4$ átomos de $\F_2$; en cada uno de ellos hay dos
trayectorias compatibles, correspondientes a los dos resultados del
tercer paso. La esperanza condicional dada $\F_2$ debe asignar a cada
átomo el valor medio de $X$ sobre las dos trayectorias compatibles. Es
una variable aleatoria $\F_2$-medible (constante sobre cada átomo de
$\F_2$). La definición axiomática captura esta idea sin recurrir a la
estructura particular del modelo.

\medskip

\noindent\textbf{La definición de Kolmogorov.}

\begin{definicion}[Esperanza condicional]
\label{def:esperanza_condicional}
Sea $X$ una variable aleatoria integrable sobre $(\Omega, \F, \PP)$ y
$\G \subseteq \F$ una sub-$\sigma$-álgebra. Una variable aleatoria $Y$
es una \emph{esperanza condicional de $X$ dado $\G$} si:
\begin{enumerate}[label=\textnormal{(\roman*)}]
  \item $Y$ es $\G$-medible;
  \item para todo $A \in \G$, $\displaystyle \int_A Y \, d\PP = \int_A X \, d\PP$.
\end{enumerate}
\end{definicion}

\begin{teorema}[Existencia y unicidad]
\label{teo:existencia_esperanza_condicional}
Sea $X$ integrable y $\G \subseteq \F$. Existe una esperanza condicional
de $X$ dado $\G$, y dos esperanzas condicionales cualesquiera coinciden
$\PP$-casi seguramente.
\end{teorema}

La demostración se apoya en el teorema de Radon-Nikodym aplicado a la
medida $\nu(A) \eqdef \int_A X \, d\PP$ sobre $(\Omega, \G)$, que es
absolutamente continua respecto a $\PP\!\restriction_\G$ y admite una
densidad $Y$ que satisface (i) y (ii). Puede consultarse en
\cite{williams}. Dado que la esperanza condicional es única salvo
igualdad casi segura, hablaremos de ``\emph{la}'' esperanza condicional
de $X$ dado $\G$, denotada $\E[X \mid \G]$, entendiendo que las
identidades en las que aparezca lo serán $\PP$-casi seguramente.

La condición (i) formaliza que la esperanza condicional es determinable
a partir de la información $\G$. La (ii) es la \emph{conservación de
los promedios}: sobre cualquier $A \in \G$, las integrales de $\E[X
\mid \G]$ y $X$ coinciden. Tomando $A = \Omega$, $\E[\E[X \mid \G]] =
\E[X]$.

\medskip

\noindent\textbf{Propiedades fundamentales.} En lo que sigue, $X, Y$
son variables aleatorias integrables y $\G, \HH \subseteq \F$ son
sub-$\sigma$-álgebras.

\begin{proposicion}[Propiedades de la esperanza condicional]
\label{prop:propiedades_esperanza_cond}
Se cumplen las siguientes propiedades $\PP$-casi seguramente.
\begin{enumerate}[label=\textnormal{(\roman*)}]
  \item \emph{Linealidad}: $\E[a X + b Y \mid \G] = a\, \E[X \mid \G]
        + b\, \E[Y \mid \G]$ para cualesquiera $a, b \in \R$.
  \item \emph{Variables medibles ya conocidas}: si $X$ es
        $\G$-medible, entonces $\E[X \mid \G] = X$.
  \item \emph{Sacar lo conocido}: si $Y$ es $\G$-medible y $XY$ es
        integrable, entonces $\E[XY \mid \G] = Y\, \E[X \mid \G]$.
  \item \emph{Propiedad de la torre}: si $\HH \subseteq \G$, entonces
        $\E\!\bigl[\E[X \mid \G] \,\big|\, \HH\bigr] = \E[X \mid \HH]$.
  \item Independencia: si $X$ es
independiente de $\G$, entonces $\E[X \mid \G] = \E[X]$.
\end{enumerate}
\end{proposicion}

A modo de ilustración (y porque será la propiedad más utilizada)
demostramos la propiedad de la torre.

\begin{proof}[Demostración de la propiedad de la torre]
Sea $Z \eqdef \E\bigl[\E[X \mid \G] \,\big|\, \HH\bigr]$. Hay que ver
que $Z$ satisface las dos condiciones de la
Definición~\ref{def:esperanza_condicional} aplicada a $X$ y $\HH$. La
primera es inmediata: $Z$ es $\HH$-medible por construcción. Para la
segunda, sea $A \in \HH$. Como $\HH \subseteq \G$, también $A \in \G$.
Aplicando (ii) primero a $Z$ con $\HH$ y luego a $\E[X \mid \G]$ con
$\G$,
\begin{equation*}
  \int_A Z \, d\PP = \int_A \E[X \mid \G] \, d\PP = \int_A X \, d\PP.
\end{equation*}
Por unicidad, $Z = \E[X \mid \HH]$ casi seguramente.
\end{proof}

La lectura intuitiva de la propiedad de la torre es la siguiente: condicionar
primero respecto de una información rica $\G$ y luego respecto de
una más pobre $\HH$ equivale a condicionar directamente respecto de
$\HH$.

\medskip

\noindent\textbf{Notación abreviada.} Si $Y_1, \ldots, Y_k$ son
variables aleatorias, escribiremos $\E[X \mid Y_1, \ldots, Y_k] \eqdef
\E[X \mid \sigma(Y_1, \ldots, Y_k)]$.

\subsection{Martingalas en tiempo discreto}
\label{subsec:martingalas_discretas}

La propiedad $\E_\Q[\widetilde S_T] = \widetilde S_0$
(ec.~\eqref{eq:martingala_uniperiodo}) que caracterizó a la medida
neutral al riesgo del modelo uniperíodo admite una generalización
natural: el concepto de \emph{martingala}.

\medskip

\noindent\textbf{Definición y primeras propiedades.} Trabajamos sobre
un espacio de probabilidad filtrado $(\Omega, \F, (\F_k), \PP)$.

\begin{definicion}[Martingala]
\label{def:martingala}
Un proceso $(M_k)_{k=0}^{N}$ adaptado a $(\F_k)$ es una
\emph{martingala} respecto de $\PP$ si:
\begin{enumerate}[label=\textnormal{(\roman*)}]
  \item $M_k$ es integrable para todo $k$;
  \item $\E[M_{k+1} \mid \F_k] = M_k$ casi seguramente, para $k = 0,
        \ldots, N-1$.
\end{enumerate}
\end{definicion}

La condición (ii) admite la siguiente lectura: la mejor predicción del valor del
proceso en el instante $k+1$, dada la información hasta $k$, es el
valor actual $M_k$. Una martingala es un proceso ``imparcial''.
Iterando (ii) y aplicando la propiedad de la torre, para $0 \leq j \leq
k \leq N$,
\begin{equation}\label{eq:martingala_iterada}
  \E[M_k \mid \F_j] = M_j \quad \text{c.s.}
\end{equation}
En particular, $\E[M_k] = \E[M_0]$ para todo $k$: la esperanza de una
martingala es constante en el tiempo.

\medskip

La propiedad de martingala depende de la medida: cambiar $\PP$ por
una medida $\Q$ adecuada hará del precio descontado del subyacente
una martingala.

\medskip

\noindent\textbf{La medida neutral al riesgo en multiperíodo.} La
condición de no arbitraje se traduce, tomando $\delta t = T/N$, en
\begin{equation}\label{eq:cond_noarb_multi}
  d < e^{r \delta t} < u,
\end{equation}
que asumimos en lo sucesivo. Definimos
\begin{equation}\label{eq:def_q_multi}
  q \eqdef \frac{e^{r \delta t} - d}{u - d}\, \in (0, 1),
\end{equation}
en analogía con el caso uniperíodo. Sobre $(\Omega, \F)$, definimos la
medida $\Q$ mediante
\begin{equation}\label{eq:def_Q_multi}
  \Q(\{\omega\}) \eqdef q^{\#\{j: \omega_j = u\}}\,
                       (1-q)^{\#\{j: \omega_j = d\}}.
\end{equation}
Bajo $\Q$, los pasos son independientes con probabilidad $q$ de subida
en cada uno. Como $q \in (0, 1)$ y $p \in (0, 1)$, las medidas $\Q$ y
$\PP$ asignan probabilidad estrictamente positiva a las mismas
trayectorias, esto es, $\Q \sim \PP$.

\medskip

\noindent\textbf{El precio descontado como martingala bajo $\Q$.} Sea
$\widetilde S_k \eqdef e^{-r\, t_k}\, S_k$ el precio descontado.

\begin{teorema}[Martingala del precio descontado]
\label{teo:martingala_precio_descontado}
En el modelo binomial multiperíodo, bajo la medida $\Q$
de~\eqref{eq:def_Q_multi}, el proceso $(\widetilde S_k)_{k=0}^{N}$ es
una martingala respecto a $(\F_k)$.
\end{teorema}

\begin{proof}
La adaptación es inmediata pues $\widetilde S_k$ es función de $S_k$,
$\F_k$-medible. La integrabilidad es trivial. Para la propiedad
martingala, escribamos $S_{k+1} = S_k\, f_{k+1}$, donde $f_{k+1}$ vale
$u$ con probabilidad $q$ y $d$ con probabilidad $1-q$ bajo $\Q$, e
independientemente de $\F_k$. Entonces
\begin{equation*}
  \widetilde S_{k+1} = e^{-r \delta t}\, \widetilde S_k\, f_{k+1}.
\end{equation*}
Aplicando ``sacar lo conocido'' (Proposición
\ref{prop:propiedades_esperanza_cond}) y la independencia,
\begin{align*}
  \E_\Q[\widetilde S_{k+1} \mid \F_k]
  &= e^{-r \delta t}\, \widetilde S_k\, \E_\Q[f_{k+1} \mid \F_k]
   = e^{-r \delta t}\, \widetilde S_k\, \E_\Q[f_{k+1}] \\
  &= e^{-r \delta t}\, \widetilde S_k\,\bigl[q\, u + (1 - q)\, d\bigr]
   = e^{-r \delta t}\, \widetilde S_k\, e^{r \delta t}
   = \widetilde S_k.
\end{align*}
\end{proof}

La medida $\Q$ recibe los mismos nombres del caso uniperíodo:
\emph{medida neutral al riesgo} o \emph{medida martingala
equivalente}. Su unicidad se demostrará en la
subsección~\ref{subsec:representacion_binomial}.

\subsection{Procesos predecibles e integral estocástica discreta}
\label{subsec:predecibles_int_disc}

\medskip

\begin{definicion}[Proceso predecible]
\label{def:predecible}
Un proceso $(H_k)_{k=1}^{N}$ es \emph{predecible} respecto a la
filtración $(\F_k)$ si, para cada $k = 1, \ldots, N$, $H_k$ es
$\F_{k-1}$-medible.
\end{definicion}

La diferencia respecto a la noción de proceso adaptado
(Definición~\ref{def:adaptado}) es de un solo índice, pero su contenido
es esencial: un proceso adaptado puede usar la información disponible
\emph{hasta} $t_k$, mientras que uno predecible solo puede usar la
disponible \emph{antes} de $t_k$. En lenguaje financiero, $H_k$ se
elige en el instante $t_{k-1}$, antes de saber cómo se moverá el
mercado en el siguiente paso. Es esta restricción la que excluye
estrategias ``con visión del futuro''.

\medskip

\begin{definicion}[Integral estocástica discreta]
\label{def:integral_discreta}
Sean $(M_k)_{k=0}^{N}$ adaptado y $(H_k)_{k=1}^{N}$ predecible. La
\emph{integral estocástica discreta} de $H$ respecto a $M$ es el
proceso
\begin{equation}\label{eq:integral_discreta}
  (H \cdot M)_n \eqdef \sum_{k=1}^{n} H_k\, (M_k - M_{k-1}),
  \qquad n = 0, 1, \ldots, N,
\end{equation}
con $(H \cdot M)_0 \eqdef 0$.
\end{definicion}

La fórmula admite una lectura financiera sugerente: si $M_k$ es el
precio de un activo y $H_k$ la cantidad mantenida durante $(t_{k-1},
t_k]$, entonces $H_k\, (M_k - M_{k-1})$ es la ganancia generada en ese
período, y $(H \cdot M)_n$ es la ganancia total hasta $t_n$.

El nombre ``integral estocástica discreta'' anticipa el papel de este
objeto en el contexto continuo: sustituyendo la suma por una integral y
los incrementos $M_k - M_{k-1}$ por incrementos infinitesimales $dM_t$,
la fórmula se transforma en $\int_0^t H_s\, dM_s$, la
\emph{integral de Itô} del Capítulo~\ref{cap:itocalculus}.

\medskip

\begin{teorema}[Estabilidad de las martingalas bajo integración estocástica]
\label{teo:integral_preserva_martingala}
Sea $(M_k)$ una $\PP$-martingala respecto a $(\F_k)$ y $(H_k)$
predecible y acotado. Entonces $\bigl((H \cdot M)_n\bigr)_{n=0}^{N}$ es
una $\PP$-martingala respecto a la misma filtración.
\end{teorema}

\begin{proof}
\emph{Adaptación e integrabilidad.} Para cada $n$, $(H \cdot M)_n$ es
suma de productos $H_k\, (M_k - M_{k-1})$ con $k \leq n$, todos
$\F_n$-medibles. La integrabilidad se sigue de la acotación de $H$ y la
integrabilidad de $M$.

\emph{Propiedad de martingala.} Como $(H \cdot M)_{n+1} = (H \cdot
M)_n + H_{n+1}(M_{n+1} - M_n)$, por linealidad
\begin{equation*}
  \E\bigl[(H \cdot M)_{n+1} \,\big|\, \F_n\bigr]
  = (H \cdot M)_n + \E\bigl[H_{n+1}(M_{n+1} - M_n) \,\big|\, \F_n\bigr],
\end{equation*}
donde hemos usado que $(H\cdot M)_n$ es $\F_n$-medible. Como $H_{n+1}$
es $\F_n$-medible (por previsibilidad), la propiedad ``sacar lo
conocido'' da
\begin{equation*}
  \E\bigl[H_{n+1}(M_{n+1} - M_n) \,\big|\, \F_n\bigr]
  = H_{n+1}\, \E[M_{n+1} - M_n \mid \F_n] = 0,
\end{equation*}
donde la última igualdad usa la propiedad de martingala de $(M_k)$. Se
concluye $\E[(H \cdot M)_{n+1} \mid \F_n] = (H \cdot M)_n$.
\end{proof}

El hecho de que previsibilidad del integrando + martingala del
integrador $\Rightarrow$ martingala de la integral es la idea
estructural que reaparecerá en el contexto continuo. Como
consecuencia inmediata, $\E[(H \cdot M)_n] = 0$ para todo $n$: la
ganancia esperada de cualquier estrategia predecible sobre una
martingala es nula.

\subsection{Teorema de representación binomial}
\label{subsec:representacion_binomial}

Cerramos la sección con el resultado que articula toda la teoría: el
\emph{teorema de representación binomial}. Establece que toda
$\Q$-martingala admite ser escrita como integral estocástica discreta
respecto al precio descontado del subyacente. Es la versión discreta y
constructiva del \emph{teorema de representación de martingalas
brownianas} del Capítulo~\ref{cap:itocalculus}: la demostración que
daremos aquí proporciona una fórmula explícita del integrando, mientras
que en el contexto continuo solo dispondremos de su existencia.

Trabajamos sobre el modelo binomial multiperíodo dotado de la medida
neutral al riesgo $\Q$ y la filtración $(\F_k)$. Denotamos $\widetilde
S_k = e^{-r k\, \delta t} S_k$.

\begin{teorema}[Representación binomial]
\label{teo:representacion_binomial}
Sea $(M_k)_{k=0}^{N}$ una $\Q$-martingala respecto a $(\F_k)$. Existe
un único proceso predecible $(\phi_k)_{k=1}^N$ tal que
\begin{equation}\label{eq:representacion_binomial}
  M_n = M_0 + \sum_{k=1}^{n} \phi_k\, (\widetilde S_k - \widetilde S_{k-1}),
  \qquad n = 0, 1, \ldots, N.
\end{equation}
\end{teorema}

\begin{proof}
Construimos $\phi$ nodo a nodo. Sea $k \in \{1, \ldots, N\}$ y $A$ un
átomo de $\F_{k-1}$. Por la estructura del modelo, $A$ se descompone
como unión disjunta de dos átomos $A_u, A_d$ de $\F_k$. Sobre cada uno
las variables $\F_k$-medibles toman un valor constante; denotemos
$M_{k-1}^A$, $M_k^{A_u}$, $M_k^{A_d}$, $\widetilde S_{k-1}^A$,
$\widetilde S_k^{A_u}$, $\widetilde S_k^{A_d}$ a los valores
correspondientes. Queremos $\phi_k$ constante sobre $A$ (denotemos
$\phi_k^A$) tal que $\eqref{eq:representacion_binomial}$ se satisfaga
sobre $A_u$ y $A_d$:
\begin{equation*}
\left\{
\begin{aligned}
  M_k^{A_u} &= M_{k-1}^A + \phi_k^A\,
              \bigl(\widetilde S_k^{A_u} - \widetilde S_{k-1}^A\bigr), \\
  M_k^{A_d} &= M_{k-1}^A + \phi_k^A\,
              \bigl(\widetilde S_k^{A_d} - \widetilde S_{k-1}^A\bigr).
\end{aligned}
\right.
\end{equation*}
Restando se elimina $M_{k-1}^A$:
\begin{equation*}
  M_k^{A_u} - M_k^{A_d}
  = \phi_k^A\, \bigl(\widetilde S_k^{A_u} - \widetilde S_k^{A_d}\bigr).
\end{equation*}
Como $\widetilde S_k^{A_u} \neq \widetilde S_k^{A_d}$,
\begin{equation}\label{eq:phi_explicito}
  \phi_k^A = \frac{M_k^{A_u} - M_k^{A_d}}{\widetilde S_k^{A_u}
                                          - \widetilde S_k^{A_d}}.
\end{equation}
La fórmula~\eqref{eq:phi_explicito} determina $\phi_k$ unívocamente
sobre cada átomo de $\F_{k-1}$. La primera ecuación queda satisfecha
por construcción; resta verificar la segunda. Multiplicando la primera
por $q$ y la segunda por $1 - q$ y sumando,
\begin{equation*}
  q\, M_k^{A_u} + (1-q)\, M_k^{A_d}
  = M_{k-1}^A + \phi_k^A\,
    \bigl[q\, \widetilde S_k^{A_u} + (1-q)\, \widetilde S_k^{A_d}
          - \widetilde S_{k-1}^A\bigr].
\end{equation*}
El lado izquierdo es $\E_\Q[M_k \mid \F_{k-1}]\restriction_A = M_{k-1}^A$
(propiedad de martingala de $M$); el corchete del lado derecho es
$\E_\Q[\widetilde S_k \mid \F_{k-1}]\restriction_A - \widetilde S_{k-1}^A
= \widetilde S_{k-1}^A - \widetilde S_{k-1}^A = 0$ por la propiedad de
martingala de $(\widetilde S_k)$
(Teorema~\ref{teo:martingala_precio_descontado}). La identidad se
verifica trivialmente, y combinándola con la primera ecuación se
recupera la segunda. Repitiendo sobre cada átomo de cada $\F_{k-1}$ se
obtiene un proceso predecible $(\phi_k)_{k=1}^N$ que
satisface~\eqref{eq:representacion_binomial}. La unicidad es inmediata
de~\eqref{eq:phi_explicito}.
\end{proof}

\medskip

\noindent\textbf{Lectura del integrando.} La
fórmula~\eqref{eq:phi_explicito} expresa $\phi_k$ como cociente de la
variación de $M$ sobre la variación del precio descontado entre los dos
hijos del átomo. Si identificamos $M$ con el valor descontado de un
derivado, $\phi_k$ es la cantidad de acción que hay que mantener para
absorber su variación posible: es la \emph{delta} del derivado en cada
nodo.

\medskip

\noindent\textbf{Dos consecuencias.} La primera es la \emph{unicidad de
la medida neutral al riesgo}, anticipada en
\ref{subsec:martingalas_discretas}. Sea $\Q'$ otra medida bajo la cual
$(\widetilde S_k)$ es martingala. Para cada átomo $A$ de $\F_{k-1}$, la
condición de martingala
\begin{equation*}
  \Q'(A_u \mid A)\, \widetilde S_k^{A_u}
  + \Q'(A_d \mid A)\, \widetilde S_k^{A_d}
  = \widetilde S_{k-1}^A,
\end{equation*}
es lineal en $\Q'(A_u \mid A) \in (0, 1)$ con coeficientes positivos, y
admite solución única, que coincide con $q$. Por consiguiente $\Q' =
\Q$. La segunda consecuencia es la \emph{completitud del modelo
binomial multiperíodo}, que encontrará justificación plena en la
subsección~\ref{subsec:tres_pasos_discreto}: como toda $\Q$-martingala
se representa como integral estocástica respecto a $\widetilde S$, la
$\Q$-martingala asociada al valor descontado de un derivado dará lugar,
vía representación, al integrando que constituye la cantidad de acción
del portfolio replicante.

\section{Solución del problema en tiempo discreto}
\label{sec:solucion_discreto}

Reunimos ahora los resultados desarrollados en la Sección~\ref{sec:multiperiodo} para obtener la fórmula del precio justo y la estrategia replicante. La construcción se lleva a cabo en tres pasos.

\subsection{Los tres pasos hacia la replicación}
\label{subsec:tres_pasos_discreto}

\noindent\textbf{Estrategias self-financing.} La idea es que una
estrategia self-financing es aquella en la que el rebalanceo de la
cartera no requiere inyectar ni retirar capital del exterior.

\begin{definicion}[Estrategia self-financing]
\label{def:self_financing_discreto}
Una \emph{estrategia} en el modelo binomial multiperíodo es un par de
procesos $\bigl((\phi_k), (\psi_k)\bigr)$ predecibles respecto a
$(\F_k)$. Su \emph{valor} en $t_k$ es $V_k \eqdef \phi_k\, S_k +
\psi_k\, B_k$ (con $V_0 \eqdef \phi_1\, S_0 + \psi_1\, B_0$). La
estrategia es \emph{self-financing} si, para todo $k = 1, \ldots, N - 1$,
\begin{equation}\label{eq:cond_self_financing_discreto}
  \phi_k\, S_k + \psi_k\, B_k = \phi_{k+1}\, S_k + \psi_{k+1}\, B_k.
\end{equation}
\end{definicion}

\begin{proposicion}[Caracterización vía proceso descontado]
\label{prop:caracterizacion_self_financing}
Sea $\bigl((\phi_k), (\psi_k)\bigr)$ una estrategia y $\widetilde V_k
\eqdef e^{-r\, t_k} V_k$. La estrategia es self-financing si y solo si
\begin{equation}\label{eq:self_financing_descontado}
  \widetilde V_n = V_0 + \sum_{k=1}^{n} \phi_k\,
                   \bigl(\widetilde S_k - \widetilde S_{k-1}\bigr),
  \qquad n = 0, 1, \ldots, N.
\end{equation}
\end{proposicion}

\begin{proof}
$(\Rightarrow)$ Calculemos $\widetilde V_k - \widetilde V_{k-1}$. Como
$\widetilde V_k = \phi_k\, \widetilde S_k + \psi_k$,
\begin{align*}
  \widetilde V_k - \widetilde V_{k-1}
  &= \phi_k\, (\widetilde S_k - \widetilde S_{k-1})
   + (\phi_k - \phi_{k-1})\, \widetilde S_{k-1}
   + (\psi_k - \psi_{k-1}).
\end{align*}
La condición de self-financing
\eqref{eq:cond_self_financing_discreto} en $t_{k-1}$, dividida por
$B_{k-1}$, implica $(\phi_k - \phi_{k-1}) \widetilde S_{k-1} + (\psi_k
- \psi_{k-1}) = 0$. Por consiguiente $\widetilde V_k - \widetilde
V_{k-1} = \phi_k(\widetilde S_k - \widetilde S_{k-1})$. Sumando se
obtiene~\eqref{eq:self_financing_descontado}.

$(\Leftarrow)$ Si se cumple~\eqref{eq:self_financing_descontado},
restando los casos $n$ y $n-1$ y combinando con $\widetilde V_k =
\phi_k\, \widetilde S_k + \psi_k$ se recupera
\eqref{eq:cond_self_financing_discreto}.
\end{proof}

La proposición relaciona la noción
financiera de ``estrategia self-financing'' con la noción matemática de
``integral estocástica discreta''. En particular, combinando con el
Teorema~\ref{teo:integral_preserva_martingala}, el valor descontado de
toda estrategia self-financing es una $\Q$-martingala.

\medskip

\noindent\textbf{La receta en tres pasos.} Sea $C_T$ una variable
aleatoria $\F_N$-medible (el \emph{payoff} del derivado).

\begin{quote}
\textbf{Receta de valoración por replicación.}
\begin{enumerate}[label=\textbf{Paso \arabic*.}, leftmargin=*]
  \item \emph{Cambio de medida.} Determinar la única medida $\Q$ bajo
        la cual $(\widetilde S_k)$ es martingala
        (Sección~\ref{subsec:martingalas_discretas}).
  \item \emph{Construcción del valor descontado.} Definir
        \begin{equation}\label{eq:def_M_replicante_discreto}
          M_k \eqdef \E_\Q\bigl[e^{-rT}\, C_T \,\big|\, \F_k\bigr].
        \end{equation}
        Por construcción, $(M_k)$ es una $\Q$-martingala con $M_N =
        e^{-rT} C_T$.
  \item \emph{Aplicación de la representación.} Por el
        Teorema~\ref{teo:representacion_binomial}, existe un único
        $(\phi_k)$ predecible tal que $M_k = M_0 + (\phi \cdot
        \widetilde S)_k$. Definir $\psi_k \eqdef M_k - \phi_k
        \widetilde S_k$.
\end{enumerate}
\end{quote}

\begin{teorema}[Valoración por replicación en el modelo binomial multiperíodo]
\label{teo:valoracion_binomial}
En el modelo binomial multiperíodo libre de arbitraje, todo derivado
con \emph{payoff} $C_T$ admite una estrategia self-financing
replicante. El precio justo en $t_n$ es
\begin{equation}\label{eq:precio_binomial_multi}
  V_n = \E_\Q\bigl[e^{-r(T - t_n)}\, C_T \,\big|\, \F_n\bigr],
\end{equation}
y en $t = 0$, $V_0 = \E_\Q[e^{-rT}\, C_T]$.
\end{teorema}

\begin{proof}
\emph{Paso 1.} La existencia y unicidad de $\Q$ se establecieron en la
Sección~\ref{subsec:martingalas_discretas}.

\emph{Paso 2.} El proceso $(M_k)$ definido
en~\eqref{eq:def_M_replicante_discreto} es una $\Q$-martingala por la
propiedad de la torre: para $j \leq k$,
\begin{equation*}
  \E_\Q[M_k \mid \F_j]
  = \E_\Q\bigl[\E_\Q[e^{-rT} C_T \mid \F_k] \,\big|\, \F_j\bigr]
  = \E_\Q[e^{-rT} C_T \mid \F_j] = M_j.
\end{equation*}

\emph{Paso 3.} Por el Teorema~\ref{teo:representacion_binomial} existe
$(\phi_k)$ predecible tal que $M_n = M_0 + \sum_{k=1}^n \phi_k(\widetilde
S_k - \widetilde S_{k-1})$. Definimos $\psi_k \eqdef M_k - \phi_k\,
\widetilde S_k$, que es predecible (combinación de cantidades
$\F_{k-1}$-medibles), y $V_n \eqdef e^{r t_n} M_n$. Por la
Proposición~\ref{prop:caracterizacion_self_financing}, la representación
de $M_n$ como integral estocástica equivale a la condición
self-financing con $V_0 = M_0$ y $\widetilde V_n = M_n$. La estrategia
es self-financing y replica el derivado, pues $V_N = e^{rT}\, M_N =
e^{rT}\, e^{-rT}\, C_T = C_T$.

Finalmente, por el principio de no arbitraje, el precio justo del
derivado en $t_n$ es el valor del portfolio en ese instante:
\begin{equation*}
  V_n = e^{r t_n}\, M_n = \E_\Q\bigl[e^{-r(T - t_n)}\, C_T \mid \F_n\bigr].
\end{equation*}
\end{proof}

La fórmula~\eqref{eq:precio_binomial_multi} generaliza
\eqref{eq:precio_esperanza} al multiperíodo: el precio justo es el
valor descontado de la esperanza del \emph{payoff} bajo la medida
neutral al riesgo, condicionada a la información en el instante de
la valoración.

\subsection{Transición al continuo}
\label{subsec:obertura_continuo}

Cerramos el capítulo con una mirada hacia adelante. El modelo binomial
multiperíodo posee una característica artificial: el mercado se
observa solo en un número finito de instantes equiespaciados. La
presente subsección argumenta, de manera deliberadamente heurística,
que el refinamiento progresivo de la partición temporal (hacer $N \to
\infty$ con $T$ fijo y $\delta t = T/N \to 0$) conduce a un modelo en
tiempo continuo cuyo subyacente sigue un movimiento browniano
geométrico. La construcción rigurosa de este modelo, así como la
fórmula de valoración asociada (la fórmula de Black-Scholes), serán el
objeto de los capítulos siguientes.

\medskip

\noindent\textbf{Calibración de los factores $u$ y $d$.} Si tomáramos
$u, d$ constantes (independientes de $\delta t$), las trayectorias del
precio se volverían dramáticamente erráticas a medida que aumentamos
$N$. Una elección sensata exige que los rendimientos logarítmicos en
cada paso sean compatibles con un comportamiento estable en el límite.
Concretamente, sean $\sigma > 0$ y $\nu \in \R$ dos constantes que
representarán la \emph{volatilidad} y la \emph{tendencia} del precio en
el modelo límite. Calibremos
\begin{equation}\label{eq:calibracion_ud}
  u = e^{\sigma \sqrt{\delta t}}, \qquad
  d = e^{-\sigma \sqrt{\delta t}},
\end{equation}
y la probabilidad subjetiva como
$p = \tfrac{1}{2}(1 + \tfrac{\nu}{\sigma}\sqrt{\delta t})$. Bajo esta
calibración, el rendimiento logarítmico $\log(S_{k+1}/S_k)$ vale
$+\sigma\sqrt{\delta t}$ con probabilidad $p$ y $-\sigma\sqrt{\delta t}$
con probabilidad $1-p$. Su esperanza es $\nu\, \delta t$ y su varianza
$\sigma^2\, \delta t + O(\delta t^2)$.

\medskip

\noindent\textbf{Convergencia al browniano geométrico.} Tenemos
\begin{equation*}
  S_N = S_0\, \exp\!\left(\sum_{k=1}^{N} \log\frac{S_k}{S_{k-1}}\right).
\end{equation*}
La suma del exponente es de variables independientes idénticamente
distribuidas, cada una con esperanza $\nu \delta t$ y varianza $\sigma^2
\delta t + O(\delta t^2)$. La suma total tiene esperanza $\nu T$ y
varianza tendiente a $\sigma^2 T$ cuando $N \to \infty$. Por el Teorema
Central del Límite,
\begin{equation*}
  \log\!\frac{S_T}{S_0} \xrightarrow{d} \mathcal{N}(\nu T,\, \sigma^2 T).
\end{equation*}
Equivalentemente,
\begin{equation}\label{eq:gbm_limite}
  S_T = S_0\, \exp\bigl(\nu T + \sigma\, W_T\bigr),
\end{equation}
con $W_T \sim \mathcal{N}(0, T)$. El argumento se extiende a cualquier
instante intermedio $t \in [0, T]$, dando
\begin{equation*}
  S_t = S_0\, \exp\bigl(\nu t + \sigma\, W_t\bigr),
\end{equation*}
con $W_t \sim \mathcal{N}(0, t)$. La familia $\{W_t\}_{t \geq 0}$ que
emerge en el límite es lo que se conoce como \emph{movimiento
browniano}: un proceso estocástico en tiempo continuo, con incrementos
gaussianos independientes y trayectorias continuas. Su construcción
rigurosa se abordará en el Capítulo~\ref{cap:brown}. El proceso $S_t$
definido por~\eqref{eq:gbm_limite} se denomina \emph{movimiento
browniano geométrico}, y constituirá el modelo canónico del precio de
la acción en el contexto continuo.

Bajo la medida neutral al riesgo $\Q$ del modelo discreto, un cálculo
análogo (sustituyendo $p$ por $q$ definido
en~\eqref{eq:def_q_multi} y desarrollando en serie en $\delta t$)
muestra que la deriva del logaritmo del precio bajo $\Q$ se convierte
en $r - \tfrac{1}{2}\sigma^2$ en lugar de $\nu$. Este \emph{cambio de
deriva} inducido por el cambio de medida es la versión discreta y
heurística del fenómeno que en el contexto continuo dará lugar al
\emph{teorema de Girsanov}, una de las piezas centrales del
Capítulo~\ref{cap:itocalculus}.

\medskip

\noindent\textbf{Recapitulación del capítulo.} Partiendo de la valoración estática de un \emph{forward}
mediante un argumento de replicación, hemos progresado hasta conseguir en el modelo binomial multiperíodo una fórmula completa y rigurosa
para el precio de cualquier derivado europeo, expresada como esperanza
descontada del \emph{payoff} bajo una medida neutral al riesgo, junto
con una estrategia self-financing que lo replica. La articulación final
en tres pasos es el esquema que reaparecerá en el contexto continuo de
Black-Scholes. Cada uno de los conceptos matemáticos introducidos
(filtraciones, esperanza condicional, martingalas, procesos
predecibles, integral estocástica discreta) admite una versión continua
que emplearemos en el Capítulo~\ref{cap:itocalculus}. Por último, la introducción hecha en la presente subsección sugiere que el
subyacente del contexto continuo será un proceso con incrementos
gaussianos independientes y trayectorias continuas: el \emph{movimiento
browniano}, al que dedicamos el Capítulo~\ref{cap:brown}.

% =====================================================================
%   CAPÍTULO 3 — MOVIMIENTO BROWNIANO
% =====================================================================

\chapter{Movimiento browniano}\label{cap:brown}

Este capítulo introduce el \emph{movimiento browniano}, base estocástica de la teoría continua de valoración. La motivación
financiera queda ahora en segundo plano: las herramientas que
necesitaremos requieren cierto desarrollo autónomo antes de volver
a aplicarlas.

\section{Definición y existencia}
\label{sec:browniano_definicion}

\subsection{Definición formal del movimiento browniano}
\label{subsec:browniano_def}

\begin{definicion}[Movimiento browniano estándar]
\label{def:browniano}
Un \emph{movimiento browniano estándar} sobre un espacio de
probabilidad $(\Omega, \F, \PP)$ es una familia $\{W_t\}_{t \geq 0}$ de
variables aleatorias reales que satisface:
\begin{enumerate}[label=\textnormal{(\roman*)}]
  \item $W_0 = 0$ casi seguramente;
  \item para cualesquiera $0 \leq s < t$, el incremento $W_t - W_s$
        tiene distribución $\mathcal{N}(0,\, t - s)$;
  \item los incrementos sobre intervalos disjuntos son
        independientes, esto es, para todo $n \geq 1$ y toda elección
        de instantes $0 \leq t_0 < t_1 < \cdots < t_n$, las variables
        aleatorias
        \begin{equation*}
          W_{t_1} - W_{t_0},\ W_{t_2} - W_{t_1},\ \ldots,\
          W_{t_n} - W_{t_{n-1}}
        \end{equation*}
        son mutuamente independientes;
  \item las trayectorias $t \mapsto W_t(\omega)$ son funciones
        continuas casi seguramente.
\end{enumerate}
\end{definicion}

Las condiciones (ii)–(iii) heredan, vía el TCL, la naturaleza
gaussiana e independiente de los pasos del random walk reescalado;
la (i) es una mera normalización (cualquier proceso que cumpla
(ii)–(iv) puede trasladarse para partir de cero) y la (iv) es la
condición que distingue al browniano cualitativamente: trayectorias
continuas, sin saltos.

Conviene observar que las condiciones (i)–(iii) son
\emph{distribucionales}, especifican la ley conjunta del proceso,
mientras que (iv) es \emph{trayectorial}: requiere que cada
realización $t \mapsto W_t(\omega)$ sea una función continua.

\newpage

\subsection{Existencia del movimiento browniano}
\label{subsec:browniano_existencia}

La definición anterior plantea inmediatamente la cuestión de si el
objeto matemático que define existe efectivamente. La respuesta es
afirmativa, pero su demostración no es trivial: requiere construir un
espacio de probabilidad $(\Omega, \F, \PP)$ y una familia de variables
aleatorias $\{W_t\}_{t \geq 0}$ sobre él que satisfaga todas las
propiedades simultáneamente, lo cual no se sigue de hechos elementales
de probabilidad.

\begin{teorema}[Existencia del movimiento browniano]
\label{teo:existencia_browniano}
Existe un espacio de probabilidad $(\Omega, \F, \PP)$ y una familia de
variables aleatorias $\{W_t\}_{t \geq 0}$ sobre él que satisface las
condiciones (i)--(iv) de la Definición~\ref{def:browniano}.
\end{teorema}

La demostración rigurosa de este teorema fue proporcionada
originariamente por Norbert Wiener en 1923 y, en una versión más
elemental basada en el sistema de funciones de Schauder, por Paul
Lévy en 1939. La idea esencial de la construcción de Lévy es la
siguiente: se definen los valores de $W$ sobre los racionales diádicos
del intervalo $[0, 1]$ mediante una serie infinita de ajustes
sucesivos, controlando que las trayectorias parciales converjan
uniformemente con probabilidad uno; el proceso así construido se
extiende por continuidad al intervalo completo. La demostración
detallada puede consultarse en \cite{karatzas-shreve} o \cite{williams}.
A partir de este punto, asumiremos como cierta la existencia del
movimiento browniano, y nos centraremos en estudiar sus propiedades.

\begin{observacion}[Sobre la condición de continuidad en la definición]
\label{obs:continuidad_definicion}
Muchos autores definen el movimiento browniano
exigiendo únicamente~(i)--(iii), que son de naturaleza distribucional:
como ya se ha señalado, especificar~(i)--(iii) equivale a fijar las
distribuciones finito-dimensionales del proceso (las leyes conjuntas de
$(W_{t_1}, \dots, W_{t_n})$, que resultan gaussianas de media nula y
covarianza $\E[W_s W_t] = \min(s, t)$). Una vez fijadas estas distribuciones, el \emph{teorema de extensión de Kolmogorov}
garantiza la existencia de un proceso con dichas distribuciones. Ahora bien, este resultado no garantiza que sus trayectorias sean continuas: la construcción de Kolmogorov permite obtener el proceso a partir de sus distribuciones finito-dimensionales, pero no proporciona por sí misma información sobre la regularidad de las trayectorias
$t \mapsto W_t(\omega)$.

Para obtener trayectorias continuas se necesita un paso más. Por el \emph{teorema de Doob}, el proceso admite una versión separable, esto es, un proceso con las mismas distribuciones finito-dimensionales cuyo comportamiento queda determinado por los valores que este toma sobre un conjunto denso numerable de instantes. Pues bien, se puede demostrar que esta versión tiene trayectorias continuas casi seguramente. Esta es la forma habitual de establecer la existencia de un movimiento browniano continuo en un curso de procesos estocásticos y constituye una alternativa a la construcción explícita de Lévy presentada anteriormente.

Exigir~(iv) desde la definición, como aquí hacemos, no es por tanto una
condición extra: la continuidad se obtiene sin coste adicional al
pasar a la versión separable, y de hecho varios autores optan por
incorporarla ya en la definición. Esta elección resulta además cómoda
para nuestros fines, pues todo el cálculo estocástico que desarrollamos a
partir del Capítulo~\ref{cap:itocalculus}, empezando por la propia
construcción de la integral de Itô, depende de la continuidad de las trayectorias del browniano.
\end{observacion}

\section{Propiedades de las trayectorias}
\label{sec:propiedades_brown}

Estudiamos en esta sección tres propiedades trayectoriales del
browniano: la no diferenciabilidad, la variación no acotada y el
cómputo explícito de la variación cuadrática.

\subsection{No diferenciabilidad de las trayectorias}
\label{subsec:no_diferenciabilidad}

La continuidad trayectorial del browniano podría sugerir una imagen
de funciones suaves; la realidad es la opuesta: las trayectorias
son tan irregulares que no admiten derivada en ningún punto.

\medskip

\noindent\textbf{Argumento heurístico.} Supongamos, por reducción al absurdo, que la trayectoria de un
movimiento browniano fuese diferenciable en algún instante $t_0$. Eso
significaría que el cociente incremental
\begin{equation*}
  \frac{W_{t_0 + h} - W_{t_0}}{h}
\end{equation*}
admite un límite finito cuando $h \to 0$. Pero por la condición (ii) de
la Definición~\ref{def:browniano}, $W_{t_0+h} - W_{t_0}$ tiene
distribución $\mathcal{N}(0, h)$, y por tanto el cociente incremental
tiene distribución $\mathcal{N}(0,\, 1/h)$. Su varianza diverge cuando
$h \to 0$: el cociente no puede converger a nada finito.

El argumento es heurístico (pasamos por alto la diferencia entre
convergencia en distribución y convergencia casi segura) pero captura
el fenómeno de fondo. La demostración rigurosa requiere refinamientos
técnicos, en particular el lema de Borel-Cantelli aplicado a estimaciones
sobre los incrementos en escalas diádicas, y puede consultarse en
\cite{karatzas-shreve}.

\begin{teorema}[No diferenciabilidad de las trayectorias brownianas]
\label{teo:no_diferenciabilidad}
Sea $\{W_t\}_{t \geq 0}$ un movimiento browniano estándar. Con
probabilidad uno, la trayectoria $t \mapsto W_t(\omega)$ no es
diferenciable en ningún punto $t \geq 0$.
\end{teorema}

Las trayectorias del browniano son, por tanto, continuas en todas partes
pero no diferenciables en ninguna, combinación que motivará la
construcción del cálculo \emph{ad hoc} del
Capítulo~\ref{cap:itocalculus}.

\subsection{Variación no acotada de las trayectorias}
\label{subsec:variacion_no_acotada}

Demostramos a continuación una propiedad cuantitativa más fuerte:
las trayectorias del browniano tienen, con probabilidad uno,
\emph{variación no acotada} en cualquier intervalo.

\medskip

\noindent\textbf{Variación total de una función.} Recordemos primero la
noción clásica.

\begin{definicion}[Variación total]
\label{def:variacion_total}
Sea $f : [0, T] \to \R$ una función real definida sobre un intervalo
compacto. La \emph{variación total} de $f$ sobre $[0, T]$ es
\begin{equation*}
  V_T(f) \eqdef \sup_{\pi}\, \sum_{k=0}^{n-1} |f(t_{k+1}) - f(t_k)|,
\end{equation*}
donde el supremo se toma sobre todas las particiones $\pi = \{0 = t_0
< t_1 < \cdots < t_n = T\}$ del intervalo. La función $f$ se dice de
\emph{variación acotada} sobre $[0, T]$ si $V_T(f) < \infty$.
\end{definicion}

Las funciones diferenciables con derivada acotada son de variación
acotada (la variación es la integral de $|f'|$), al igual que las
funciones monótonas. Para una función de variación acotada, la
integral de Riemann-Stieltjes $\int_0^T g\, df$ está definida sin
ambigüedades para todo integrando $g$ continuo. Es esta circunstancia
la que falla en el caso del browniano.

\medskip

\noindent\textbf{Argumento heurístico.} Sea $\pi_n = \{0 = t_0 < t_1
< \cdots < t_n = T\}$ la partición uniforme de $[0, T]$ en $n$
subintervalos de longitud $T/n$. Por la condición (ii) de la
Definición~\ref{def:browniano}, cada incremento $W_{t_{k+1}} -
W_{t_k}$ tiene distribución $\mathcal{N}(0,\, T/n)$, y por tanto su
valor absoluto satisface
\begin{equation*}
  \E\bigl[|W_{t_{k+1}} - W_{t_k}|\bigr]
  = \sqrt{\frac{2T}{\pi\, n}}.
\end{equation*}
Sumando sobre los $n$ subintervalos,
\begin{equation*}
  \E\!\left[\,\sum_{k=0}^{n-1} |W_{t_{k+1}} - W_{t_k}|\right]
  = n\, \sqrt{\frac{2T}{\pi\, n}}
  = \sqrt{\frac{2 T n}{\pi}}.
\end{equation*}
La esperanza diverge cuando $n \to \infty$: refinar la partición hace
que la suma de los valores absolutos de los incrementos crezca sin
límite, en lugar de aproximarse a un valor finito como ocurriría con
una función de variación acotada. El razonamiento sugiere, por tanto,
que la trayectoria browniana no puede ser de variación acotada.

El argumento anterior es heurístico (demuestra divergencia en
esperanza, no casi segura) pero muestra claramente el fenómeno. La demostración
rigurosa puede obtenerse, entre otras vías, como corolario del cómputo
de la variación cuadrática que abordamos en la subsección siguiente
(véase la observación al final de dicha subsección).

\begin{teorema}[Variación no acotada de las trayectorias brownianas]
\label{teo:variacion_no_acotada}
Sea $\{W_t\}_{t \geq 0}$ un movimiento browniano estándar y $T > 0$.
Con probabilidad uno, la trayectoria $t \mapsto W_t(\omega)$ tiene
variación no acotada sobre $[0, T]$, esto es, $V_T(W) = \infty$ casi
seguramente.
\end{teorema}

\medskip

\noindent\textbf{Consecuencia: el fracaso de Riemann-Stieltjes.} La
variación no acotada de las trayectorias tiene una consecuencia
inmediata sobre la posibilidad de definir integrales contra el
browniano. La construcción clásica de Riemann-Stieltjes de $\int_0^T
g\, df$ requiere que $f$ sea de variación acotada para que las sumas
\begin{equation*}
  \sum_{k=0}^{n-1} g(\tau_k)\, \bigl(f(t_{k+1}) - f(t_k)\bigr)
\end{equation*}
converjan a un mismo límite con independencia de la elección de los
puntos $\tau_k \in [t_k, t_{k+1}]$. Cuando $f = W$, este resultado
falla: el límite depende de la elección de los $\tau_k$. Construir una integral $\int_0^T g_s\, dW_s$ requerirá, por
tanto, fijar una convención específica para los puntos de
evaluación y verificar que la construcción resultante está bien
definida.

\subsection{Variación cuadrática del movimiento browniano}
\label{subsec:variacion_cuadratica}

Pese a la irregularidad de las trayectorias, la \emph{suma de los
cuadrados} de los incrementos presenta un comportamiento notablemente regular: converge a un límite finito cuando se refina la
partición. Esta cantidad, la \emph{variación cuadrática}, es el
resultado central del capítulo.

\medskip

\begin{teorema}[Variación cuadrática del browniano]
\label{teo:variacion_cuadratica}
Sea $\{W_t\}_{t \geq 0}$ un movimiento browniano estándar y $T > 0$.
Para cada partición $\pi = \{0 = t_0 < t_1 < \cdots < t_n = T\}$ de
$[0, T]$, denotemos su norma por $|\pi| \eqdef \max_{k} (t_{k+1} -
t_k)$ y consideremos la variable aleatoria
\begin{equation*}
  Q_\pi \eqdef \sum_{k=0}^{n-1} (W_{t_{k+1}} - W_{t_k})^2.
\end{equation*}
Entonces, cuando $|\pi| \to 0$,
\begin{equation}\label{eq:convergencia_var_cuadratica}
  Q_\pi \xrightarrow{L^2} T,
\end{equation}
es decir, $\E\bigl[(Q_\pi - T)^2\bigr] \to 0$.
\end{teorema}

La cantidad $Q_\pi$ es, por su definición, una variable
aleatoria (su valor depende de la trayectoria $\omega$). El teorema
afirma que, en el límite de particiones finas, esta variable aleatoria
deja de ser aleatoria: converge a una constante determinista, $T$,
igual a la longitud del intervalo. Es un resultado a primera vista
sorprendente, pues uno esperaría que una suma de cuadrados de
incrementos brownianos (cada uno con varianza no nula) mantuviera
algún grado de aleatoriedad en el límite. La razón por la que esto no
ocurre es que, a medida que la partición se refina, las contribuciones
de cada subintervalo se vuelven cada vez más numerosas y cada vez más
pequeñas; la ley de los grandes números, aplicada al sumatorio,
provoca el promediado a la constante.

\begin{proof}
La estrategia consiste en calcular directamente la esperanza y la
varianza de $Q_\pi$, y mostrar que en el límite la primera tiende a
$T$ y la segunda a cero.

\emph{Paso 1: cálculo de la esperanza.} Por la condición (ii) de la
Definición~\ref{def:browniano}, cada incremento $W_{t_{k+1}} - W_{t_k}$
tiene distribución $\mathcal{N}(0,\, t_{k+1} - t_k)$, y por tanto
\begin{equation*}
  \E\bigl[(W_{t_{k+1}} - W_{t_k})^2\bigr] = t_{k+1} - t_k.
\end{equation*}
Por linealidad de la esperanza,
\begin{equation}\label{eq:esperanza_Q_pi}
  \E[Q_\pi] = \sum_{k=0}^{n-1} (t_{k+1} - t_k) = T.
\end{equation}
La esperanza de $Q_\pi$ es exactamente $T$, sin error: este resultado
es \emph{exacto}, no asintótico, y se cumple para cualquier partición.

\emph{Paso 2: cálculo de la varianza.} Por definición,
\begin{equation*}
  \mathrm{Var}(Q_\pi) = \E\!\left[\bigl(Q_\pi - T\bigr)^2\right]
                      = \E\!\left[\biggl(\sum_{k=0}^{n-1}
                          \bigl[(W_{t_{k+1}} - W_{t_k})^2 - (t_{k+1} -
                          t_k)\bigr]\biggr)^{\!2}\right].
\end{equation*}
Expandiendo el cuadrado,
\begin{align*}
  \mathrm{Var}(Q_\pi)
  &= \sum_{k=0}^{n-1} \E\!\left[\bigl((W_{t_{k+1}} - W_{t_k})^2 -
                       (t_{k+1} - t_k)\bigr)^{\!2}\right] \\
  &\quad + 2\sum_{j < k} \E\!\Bigl[\bigl((W_{t_{j+1}} - W_{t_j})^2 -
                  (t_{j+1} - t_j)\bigr) \\
  &\qquad\qquad\qquad \cdot \bigl((W_{t_{k+1}} - W_{t_k})^2 -
                  (t_{k+1} - t_k)\bigr)\Bigr].
\end{align*}
Los términos cruzados ($j < k$) se anulan: por la independencia de los
incrementos sobre intervalos disjuntos (condición (iii) de la
Definición~\ref{def:browniano}), las variables $(W_{t_{j+1}} -
W_{t_j})^2 - (t_{j+1} - t_j)$ y $(W_{t_{k+1}} - W_{t_k})^2 - (t_{k+1} -
t_k)$ son independientes y, además, ambas tienen esperanza nula
(por~\eqref{eq:esperanza_Q_pi} aplicada a un solo subintervalo). Por
tanto, su producto tiene esperanza cero, y solo sobreviven los
términos diagonales:
\begin{equation*}
  \mathrm{Var}(Q_\pi)
  = \sum_{k=0}^{n-1} \E\!\left[\bigl((W_{t_{k+1}} - W_{t_k})^2 -
                       (t_{k+1} - t_k)\bigr)^{\!2}\right].
\end{equation*}
Para cada sumando, denotemos $\Delta_k \eqdef t_{k+1} - t_k$ y
expandamos:
\begin{align*}
  \E\!\left[\bigl((W_{t_{k+1}} - W_{t_k})^2 - \Delta_k\bigr)^{\!2}\right]
  &= \E\bigl[(W_{t_{k+1}} - W_{t_k})^4\bigr]
   - 2\, \Delta_k\, \E\bigl[(W_{t_{k+1}} - W_{t_k})^2\bigr]
   + \Delta_k^2 \\
  &= 3\, \Delta_k^2 - 2\, \Delta_k^2 + \Delta_k^2 = 2\, \Delta_k^2,
\end{align*}
donde hemos usado que el cuarto momento de una $\mathcal{N}(0,
\Delta_k)$ es $3\, \Delta_k^2$.

Sustituyendo,
\begin{equation}\label{eq:varianza_Q_pi}
  \mathrm{Var}(Q_\pi) = 2 \sum_{k=0}^{n-1} \Delta_k^2.
\end{equation}

\emph{Paso 3: paso al límite.} Para acotar la suma anterior, observamos
que
\begin{equation*}
  \sum_{k=0}^{n-1} \Delta_k^2
  \leq \max_k \Delta_k\, \sum_{k=0}^{n-1} \Delta_k
  = |\pi|\, T.
\end{equation*}
Por consiguiente, $\mathrm{Var}(Q_\pi) \leq 2\, |\pi|\, T \to 0$ cuando
$|\pi| \to 0$. Combinando con $\E[Q_\pi] = T$ del Paso 1,
\begin{equation*}
  \E\bigl[(Q_\pi - T)^2\bigr] = \mathrm{Var}(Q_\pi) \to 0,
\end{equation*}
que es la convergencia $L^2$ enunciada en
\eqref{eq:convergencia_var_cuadratica}.
\end{proof}

\medskip

\noindent\textbf{La regla heurística $(dW_t)^2 = dt$.} El Teorema~\ref{teo:variacion_cuadratica}
admite una interpretación informal particularmente ilustrativa. La afirmación de que
$\sum_k (W_{t_{k+1}} - W_{t_k})^2 \to T$ cuando la partición se refina
puede leerse, formalmente, como
\begin{equation*}
  \int_0^T (dW_t)^2 = T,
\end{equation*}
o bien, en forma diferencial, como la regla
\begin{equation}\label{eq:regla_dW2}
  (dW_t)^2 = dt.
\end{equation}
La identidad~\eqref{eq:regla_dW2} es, evidentemente, una abreviatura informal que carece por sí sola de significado matemático preciso.
Pero captura un hecho fundamental: aunque el ``incremento
infinitesimal'' $dW_t$ es de orden $\sqrt{dt}$ (su desviación
típica), su cuadrado es de orden $dt$, no $dt^2$ como sería para
una función suave. En el cálculo clásico los términos cuadráticos
en $dt$ son despreciables; en el browniano, no.

\medskip

\noindent\textbf{Consecuencia: la variación no acotada.} Cerramos la
sección retomando la afirmación anunciada en la
subsección~\ref{subsec:variacion_no_acotada}: la variación cuadrática
finita y no nula del browniano implica, con probabilidad uno, su
variación no acotada. El argumento es el siguiente. Sobre cualquier
partición $\pi$,
\begin{equation*}
  Q_\pi = \sum_{k=0}^{n-1} (W_{t_{k+1}} - W_{t_k})^2
        \leq \max_k |W_{t_{k+1}} - W_{t_k}|
              \cdot \sum_{k=0}^{n-1} |W_{t_{k+1}} - W_{t_k}|.
\end{equation*}
Si la trayectoria fuera de variación acotada, el segundo factor sería
finito; por la continuidad uniforme de $W$ sobre $[0, T]$ (garantizada
por la condición (iv) de la Definición~\ref{def:browniano}), el primer
factor tendería a cero cuando $|\pi| \to 0$. El producto tendería por
tanto a cero, contradiciendo el Teorema~\ref{teo:variacion_cuadratica}
que asegura $Q_\pi \to T > 0$. Por consiguiente, la trayectoria no
puede ser de variación acotada.

\section{Filtración browniana y martingalas continuas}
\label{sec:filtracion_browniana}

Cerramos el capítulo con la filtración asociada al browniano y los
ejemplos canónicos de martingalas continuas. La sección es breve:
su contenido es esencialmente la adaptación al contexto continuo de
las nociones discretas del Capítulo~\ref{cap:discreto}.

\subsection{Filtración browniana y martingalas asociadas}
\label{subsec:filtracion_y_martingalas}

\noindent\textbf{Filtración natural del browniano.} La filtración
(Definición~\ref{def:filtracion}) se generaliza al contexto continuo
sustituyendo la sucesión $(\F_k)_{k=0}^N$ por una familia
$(\F_t)_{t \geq 0}$ con la misma propiedad de monotonía $\F_s
\subseteq \F_t$ para $s \leq t$.

\begin{definicion}[Filtración natural del browniano]
\label{def:filtracion_browniana}
Sea $\{W_t\}_{t \geq 0}$ un movimiento browniano estándar sobre
$(\Omega, \F, \PP)$. La \emph{filtración natural} de $W$ es la familia
$(\F_t^W)_{t \geq 0}$ definida por
\begin{equation*}
  \F_t^W \eqdef \sigma\bigl(W_s : 0 \leq s \leq t\bigr),
\end{equation*}
es decir, la $\sigma$-álgebra generada por los valores del proceso en el
intervalo $[0, t]$.
\end{definicion}

La filtración natural codifica con precisión la información generada por
el browniano, pero presenta dos limitaciones técnicas que conviene
corregir antes de utilizarla como marco operativo. La primera es que
no contiene necesariamente los conjuntos $\PP$-nulos de $\F$: una
modificación de un proceso adaptado sobre un conjunto de medida nula
puede dejar de ser adaptada, lo que impide trabajar con cuestiones
\emph{casi seguro} con la flexibilidad necesaria. La segunda es que
no se garantiza la \emph{continuidad por la derecha} (es decir, la
identidad $\F_t^W = \bigcap_{s > t} \F_s^W$), propiedad imprescindible para el manejo de tiempos de parada y para la
existencia de versiones continuas de procesos como la integral de
Itô (Capítulo~\ref{cap:itocalculus}).

La solución estándar consiste en pasar a una filtración ligeramente
ampliada que sí cumpla ambas propiedades.

\begin{definicion}[Filtración aumentada del browniano]
\label{def:filtracion_aumentada}
Sea $\mathcal{N} \eqdef \{N \in \F : \PP(N) = 0\}$ la colección de
conjuntos $\PP$-nulos. La \emph{filtración aumentada} del browniano
es la familia $(\F_t)_{t \geq 0}$ definida por
\begin{equation*}
  \F_t \eqdef \sigma\bigl(\F_t^W \cup \mathcal{N}\bigr),
\end{equation*}
es decir, la menor $\sigma$-álgebra que contiene a $\F_t^W$ y a todos
los conjuntos $\PP$-nulos.
\end{definicion}

Esta definición resuelve la primera limitación por construcción.
Es más, también resuelve la segunda de manera automática:
para el browniano, la filtración aumentada satisface por sí sola
la propiedad de continuidad por la derecha. El resultado puede consultarse en \cite{karatzas-shreve}, Proposición~2.7.7, y se
sustenta en la ley cero-uno de Blumenthal.

\begin{proposicion}[Continuidad por la derecha de la filtración
aumentada]
\label{prop:filtracion_continua_derecha}
La filtración $(\F_t)_{t \geq 0}$ de la
Definición~\ref{def:filtracion_aumentada} es continua por la
derecha: $\F_t = \bigcap_{s > t} \F_s$ para todo $t \geq 0$.
\end{proposicion}

En consecuencia, $(\F_t)$ satisface las llamadas \emph{condiciones
habituales}, completitud y continuidad por la derecha, marco
estándar sobre el que se desarrolla la teoría general de procesos
continuos. \textbf{En lo sucesivo, y a lo largo de todo el resto del
trabajo, $(\F_t)$ denotará sin excepción la filtración aumentada del
browniano.} La distinción con la filtración natural $\F_t^W$ no
volverá a explicitarse, pues a efectos informacionales ambas
codifican los mismos eventos: cualquier conjunto de $\F_t$ difiere
de uno de $\F_t^W$ a lo sumo en un conjunto $\PP$-nulo.

\medskip

\begin{definicion}[Martingala en tiempo continuo]
\label{def:martingala_continua}
Sea $(\Omega, \F, (\F_t)_{t \geq 0}, \PP)$ un espacio de probabilidad
filtrado. Un proceso estocástico $\{M_t\}_{t \geq 0}$ es una
\emph{martingala} respecto de $(\F_t)$ si:
\begin{enumerate}[label=\textnormal{(\roman*)}]
  \item $M_t$ es $\F_t$-medible para todo $t \geq 0$ (\emph{adaptación});
  \item $\E[|M_t|] < \infty$ para todo $t \geq 0$ (\emph{integrabilidad});
  \item $\E[M_t \mid \F_s] = M_s$ casi seguramente, para todo $0 \leq s
        \leq t$.
\end{enumerate}
\end{definicion}

Es la traducción literal al contexto continuo de la
Definición~\ref{def:martingala}; las propiedades elementales
(linealidad, esperanza constante en el tiempo) se preservan, con
demostraciones idénticas al caso discreto.

\medskip

\noindent\textbf{Ejemplos canónicos.} Presentamos dos martingalas
continuas asociadas al movimiento browniano. La filtración
considerada es la $(\F_t)$ del browniano fijada en la
Definición~\ref{def:filtracion_aumentada}.

\begin{proposicion}[Martingalas brownianas]
\label{prop:martingalas_brownianas}
Sea $\{W_t\}_{t \geq 0}$ un movimiento browniano estándar. Los
siguientes procesos son martingalas respecto a $(\F_t)$:
\begin{enumerate}[label=\textnormal{(\alph*)}]
  \item el propio browniano, $\{W_t\}_{t \geq 0}$;
  \item para cada $\lambda \in \R$, la \emph{martingala exponencial}
        \begin{equation}\label{eq:martingala_exponencial}
          \mathcal{E}_t^{\lambda} \eqdef
          \exp\!\Bigl(\lambda\, W_t - \tfrac{1}{2}\lambda^2\, t\Bigr).
        \end{equation}
\end{enumerate}
\end{proposicion}

\begin{proof}
La adaptación de los dos procesos a $(\F_t)$ es inmediata, pues son
funciones de $W_t$. La integrabilidad se sigue de las propiedades
gaussianas del browniano. Falta verificar la propiedad (iii) de la
Definición~\ref{def:martingala_continua}.

\emph{(a)} Para $0 \leq s \leq t$, escribimos $W_t = W_s + (W_t - W_s)$,
y por la condición (iii) de la Definición~\ref{def:browniano}, el
incremento $W_t - W_s$ es independiente de $\F_s$. Por consiguiente,
\begin{equation*}
  \E[W_t \mid \F_s] = W_s + \E[W_t - W_s \mid \F_s]
                    = W_s + \E[W_t - W_s] = W_s,
\end{equation*}
ya que $\E[W_t - W_s] = 0$ por la condición (ii) de la
Definición~\ref{def:browniano}.

\emph{(b)} Para la martingala exponencial,
\begin{equation*}
  \mathcal{E}_t^{\lambda}
  = \exp\!\bigl(\lambda W_s - \tfrac{1}{2}\lambda^2 s\bigr)
    \cdot \exp\!\bigl(\lambda(W_t - W_s) - \tfrac{1}{2}\lambda^2(t - s)\bigr)
  = \mathcal{E}_s^{\lambda} \cdot
    \exp\!\bigl(\lambda(W_t - W_s) - \tfrac{1}{2}\lambda^2(t - s)\bigr).
\end{equation*}
El primer factor es $\F_s$-medible y el segundo es independiente de
$\F_s$. Por la propiedad ``sacar lo conocido'' y la independencia,
\begin{equation*}
  \E\bigl[\mathcal{E}_t^{\lambda} \mid \F_s\bigr]
  = \mathcal{E}_s^{\lambda}\,
    \E\!\left[\exp\!\bigl(\lambda(W_t - W_s)\bigr)\right]
    \cdot e^{-\frac{1}{2}\lambda^2(t-s)}.
\end{equation*}
Como $W_t - W_s \sim \mathcal{N}(0, t-s)$, su función generadora de
momentos da $\E[\exp(\lambda(W_t-W_s))] = e^{\frac{1}{2}\lambda^2(t-s)}$,
y el factor exponencial cancela exactamente con el último término, de forma que $\E[\mathcal{E}_t^{\lambda} \mid \F_s] = \mathcal{E}_s^{\lambda}$.
\end{proof}

\medskip

\noindent\textbf{Comentario sobre la martingala exponencial.} El
proceso~\eqref{eq:martingala_exponencial}, conocido también como
\emph{exponencial estocástica} de $W$, jugará un papel importante en el
Capítulo~\ref{cap:itocalculus}: una versión generalizada de la
fórmula~\eqref{eq:martingala_exponencial}, donde $\lambda$ se sustituye
por un proceso adaptado $\theta_t$, es precisamente la derivada de
Radon-Nikodym que define el cambio de medida en el teorema de Girsanov.
La forma exponencial cuadrática del exponente, $-\tfrac{1}{2}\lambda^2
t$, no es un capricho: es exactamente la corrección que hace que la
exponencial tenga esperanza unidad y, por tanto, que la propiedad de
martingala se preserve.

\medskip

Cerrado el estudio del browniano y las martingalas asociadas a él, estamos en condiciones de desarrollar el cálculo estocástico que
ocupará el siguiente capítulo. Las herramientas que allí se introducirán (integral de Itô, fórmula de Itô, teorema de Girsanov, teorema de
representación de martingalas brownianas) se apoyan esencialmente en
las propiedades estudiadas hasta aquí.

% =====================================================================
%   CAPÍTULO 4 — CÁLCULO ESTOCÁSTICO
% =====================================================================
\chapter{Cálculo estocástico}\label{cap:itocalculus}

Este capítulo desarrolla las bases matemáticas del cálculo estocástico sobre las que se construirá la teoría continua de
valoración del Capítulo~\ref{cap:bs}. Sus pilares fundamentales son cuatro:
la \emph{integral estocástica de Itô}, que generaliza al contexto
browniano la integral estocástica discreta del
Capítulo~\ref{cap:discreto}; la \emph{fórmula de Itô}, que sustituye a
la regla de la cadena clásica para procesos brownianos; el \emph{teorema
de Girsanov}, que formaliza el cambio de medida en tiempo continuo; y
el \emph{teorema de representación de martingalas brownianas}, que
constituye el análogo continuo del teorema de representación binomial
de la Sección~\ref{subsec:representacion_binomial}. Cada uno de estos pilares tiene así su correspondencia con el caso discreto del Capítulo~\ref{cap:discreto}: la maquinaria que desarrollamos aquí no
introduce conceptos esencialmente nuevos, sino que traduce ideas ya
conocidas al contexto continuo, donde la irregularidad de las trayectorias del
browniano (Capítulo~\ref{cap:brown}) exige una construcción técnica considerablemente más cuidadosa.

\section{Motivación: la necesidad de una nueva integral}
\label{sec:motivacion_ito}

Antes de emprender la construcción de la integral estocástica, es
imprescindible articular con precisión por qué es necesaria una
construcción \emph{ad hoc} y por qué los métodos clásicos de
integración no son suficientes. La motivación se articula en tres
pasos: primero, mostrar que el proceso del precio del subyacente debe
necesariamente exhibir variación no acotada para que el modelo sea
económicamente sensato; segundo, comprobar que dicha condición hace
fracasar el procedimiento clásico de Riemann-Stieltjes; y tercero,
identificar la elección que conducirá a la definición correcta.

\medskip

\noindent\textbf{El precio del subyacente no puede tener trayectorias
suaves.} La primera observación, anclada en el principio de no
arbitraje, es la siguiente: en un modelo continuo libre de arbitraje,
las trayectorias del precio descontado del subyacente no pueden ser
de variación acotada. El argumento, que recogemos siguiendo
\cite{etheridge}, es constructivo: exhibiremos explícitamente la
oportunidad de arbitraje que surgiría en caso contrario. Trabajamos
con el precio descontado $\widetilde S_t = e^{-rt} S_t$, sobre el que
la condición de autofinanciación adopta una forma particularmente
sencilla; las propiedades de variación de $S$ y $\widetilde S$ son
equivalentes bajo $B_t = e^{rt}$.

La definición rigurosa de cartera autofinanciada en tiempo continuo
se formalizará en el Capítulo~\ref{cap:bs}; aquí solo necesitamos
una de sus consecuencias, análoga directa
de~\eqref{eq:integral_discreta} en el caso discreto: el valor
descontado en $T$ de una cartera autofinanciada partiendo de capital
inicial $V_0$ con estrategia $\phi$ adaptada y de trayectorias
continuas es
\begin{equation}\label{eq:cartera_descontada_continua_BV}
  \widetilde V_T = V_0 + \int_0^T \phi_t\, d\widetilde S_t,
\end{equation}
donde la integral, bajo la hipótesis de variación acotada del
integrador, se interpreta pathwise como integral de
Riemann–Stieltjes clásica.

Sea $F : \R \to \R$ de clase $C^1$ y consideremos la estrategia
$\phi_t = F'(\widetilde S_t)$, partiendo de $V_0 = F(\widetilde S_0)$.
Por el teorema fundamental del cálculo aplicado al integrador
continuo y de variación acotada $\widetilde S$,
\begin{equation*}
  \int_0^T F'(\widetilde S_t)\, d\widetilde S_t
  = F(\widetilde S_T) - F(\widetilde S_0),
\end{equation*}
y sustituyendo en~\eqref{eq:cartera_descontada_continua_BV} se
obtiene $\widetilde V_T = F(\widetilde S_T)$. Bajo la hipótesis de
variación acotada, cualquier cantidad de la forma $F(\widetilde S_T)$
con $F\in C^1$ resulta así replicable a partir de capital inicial
$F(\widetilde S_0)$.

Particularicemos a $F(x) = (x - \widetilde S_0)^2$: la estrategia
$\phi_t = 2\,(\widetilde S_t - \widetilde S_0)$, partiendo de $V_0 =
0$, genera valor nominal en $T$
\begin{equation*}
  V_T = e^{rT}\,\bigl(\widetilde S_T - \widetilde S_0\bigr)^2,
\end{equation*}
no negativo casi seguro y estrictamente positivo sobre el evento
$\{\widetilde S_T \neq \widetilde S_0\}$, de probabilidad
estrictamente positiva en cualquier modelo no trivial. Esto es, genera una
oportunidad de arbitraje partiendo de capital nulo, en
contradicción con la hipótesis fundamental del modelo.

La conclusión, por tanto, es que cualquier modelo realista del precio en
tiempo continuo deberá emplear procesos cuyas trayectorias tengan
variación no acotada. El movimiento browniano introducido en el
Capítulo~\ref{cap:brown}, lejos de ser una excentricidad matemática,
es exactamente el tipo de proceso que la teoría exige.

\medskip

\noindent\textbf{Por qué Riemann-Stieltjes no sirve.} Una vez aceptado
que el integrador en una expresión del tipo $\int_0^T \phi_s\, dW_s$
tendrá variación no acotada, surge inmediatamente el problema técnico:
la construcción clásica de Riemann-Stieltjes deja de funcionar.
Recordemos brevemente por qué.

En la integral de Riemann-Stieltjes, dada una partición $\pi = \{0 =
t_0 < t_1 < \cdots < t_n = T\}$ y puntos $\tau_k \in [t_k, t_{k+1}]$,
se considera la suma
\begin{equation*}
  S_\pi(\phi, f; \tau) \eqdef \sum_{k=0}^{n-1} \phi(\tau_k)\,
                              \bigl(f(t_{k+1}) - f(t_k)\bigr).
\end{equation*}
Si $f$ es de variación acotada y $\phi$ es continua, esta suma converge
a un mismo límite cuando $|\pi| \to 0$, con independencia de la
elección de los puntos $\tau_k$. Este límite común es, por definición,
$\int_0^T \phi\, df$.

Cuando $f$ tiene variación no acotada, en cambio, el resultado se
viene abajo: el límite, si existe, depende de la elección de los
puntos $\tau_k$. El ejemplo más transparente es el del propio browniano
con $\phi = W$, esto es, las sumas que aproximarían a $\int_0^T W_s\,
dW_s$. Consideremos la partición uniforme $t_k = kT/n$ y la elección
del extremo izquierdo ($\tau_k = t_k$), y observemos la identidad
algebraica
\begin{equation*}
  W_{t_k}(W_{t_{k+1}} - W_{t_k})
  = \tfrac{1}{2}\bigl(W_{t_{k+1}}^2 - W_{t_k}^2\bigr)
   - \tfrac{1}{2}(W_{t_{k+1}} - W_{t_k})^2.
\end{equation*}
Sumando para $k = 0, 1, \ldots, n-1$, la primera parte es telescópica
y produce $\tfrac{1}{2}(W_T^2 - W_0^2) = \tfrac{1}{2} W_T^2$, mientras
que la segunda parte tiende, por el
Teorema~\ref{teo:variacion_cuadratica}, a $\tfrac{1}{2} T$ en el
sentido de la convergencia $L^2$. En consecuencia,
\begin{equation}\label{eq:suma_izq}
  \sum_{k=0}^{n-1} W_{t_k}(W_{t_{k+1}} - W_{t_k})
  \xrightarrow{L^2} \tfrac{1}{2}\bigl(W_T^2 - T\bigr).
\end{equation}
La elección análoga del extremo derecho ($\tau_k = t_{k+1}$) produce,
por un cálculo enteramente paralelo en el que el signo del término
cuadrático cambia, el límite $\tfrac{1}{2}(W_T^2 + T)$. Las dos sumas
aproximantes tienden, por tanto, a límites distintos: difieren
exactamente en $T$, que es la variación cuadrática del browniano sobre
$[0, T]$. La causa es ahora evidente: la diferencia entre las dos
sumas es $\sum_k (W_{t_{k+1}} - W_{t_k})^2$, que para una función de
variación acotada tendería a cero, pero que para el browniano
converge a $T$. La patología que rompe Riemann-Stieltjes es exactamente
la propiedad que distingue al browniano de los procesos clásicos.

\medskip

\noindent\textbf{La elección de Itô.} La conclusión del análisis
anterior es clara: para construir una integral $\int_0^T \phi_s\, dW_s$
bien definida, no podemos limitarnos a tomar el límite de las sumas
parciales, sino que debemos fijar una convención específica para los
puntos de evaluación. La elección que adoptaremos, la \emph{elección
de Itô}, consiste en evaluar siempre el integrando en el extremo
\emph{izquierdo} de cada subintervalo,
\begin{equation*}
  \sum_{k=0}^{n-1} \phi_{t_k}\, (W_{t_{k+1}} - W_{t_k}).
\end{equation*}
Esta elección está justificada por dos razones. La primera, de naturaleza matemática, es que la integral resultante
preservará la propiedad de martingala
(Teorema~\ref{teo:martingala_integral_indefinida}). La segunda, de naturaleza financiera, es que el extremo
izquierdo formaliza la noción de \emph{previsibilidad}: la cantidad de
acción que se mantiene durante el período $[t_k, t_{k+1})$ se elige al
inicio del período, en $t = t_k$, antes de saber cómo se moverá el
mercado durante el subintervalo. Es exactamente la misma idea que
articulamos en el caso discreto en la
Sección~\ref{subsec:predecibles_int_disc}, ahora trasladada al
contexto continuo.

\section{Construcción de la integral estocástica de Itô}
\label{sec:construccion_ito}

Una vez expuesta la motivación, pasemos a estudiar la construcción rigurosa de la
integral estocástica $\int_0^T \phi_s\, dW_s$. La estrategia general
sigue un esquema estándar: definimos la integral sobre una clase de
integrandos elementales donde el objeto se reduce a una suma finita;
demostramos sobre dicha clase una identidad fundamental que controla
la norma $L^2$ del resultado; y extendemos la integral por densidad a
una clase suficientemente amplia de integrandos, aprovechando la
completitud de $L^2$. Fijamos a lo largo de toda la sección un horizonte temporal $T
> 0$ y un espacio de probabilidad completo $(\Omega, \F, \PP)$
sobre el que se considera un movimiento browniano estándar
$\{W_t\}_{t \in [0, T]}$ con su filtración aumentada $(\F_t)_{0
\le t \le T}$.

\subsection{Procesos simples e integral elemental}
\label{subsec:integral_simples}

Comenzamos definiendo la clase de integrandos elementales sobre la
que la integral estocástica admite una definición directa, sin
dificultades técnicas.

\begin{definicion}[Proceso simple predecible]
\label{def:proceso_simple}
Un proceso estocástico $\phi : [0, T] \times \Omega \to \R$ es
\emph{simple predecible} —o, abreviadamente, \emph{simple}— si
existe una partición $0 = t_0 < t_1 < \cdots < t_n = T$ y variables
aleatorias $\xi_0, \xi_1, \ldots, \xi_{n-1}$ tales que:
\begin{enumerate}[label=\textnormal{(\roman*)}]
  \item para cada $i = 0, 1, \ldots, n-1$, $\xi_i$ es acotada y
        $\F_{t_i}$-medible;
  \item el proceso $\phi$ admite la representación
        \begin{equation}\label{eq:proceso_simple}
          \phi_t(\omega) = \sum_{i=0}^{n-1} \xi_i(\omega)\,
                           \mathbf{1}_{(t_i, t_{i+1}]}(t).
        \end{equation}
\end{enumerate}
Denotaremos por $\mathcal{S}$ la clase de todos los procesos simples
sobre $[0, T]$.
\end{definicion}

Detengámonos brevemente en la condición (i). La medibilidad de
$\xi_i$ respecto a la $\sigma$-álgebra $\F_{t_i}$ (la información
disponible al inicio del subintervalo $(t_i, t_{i+1}]$, no al final) es la traducción exacta al contexto continuo de la noción de
\emph{previsibilidad}\footnote{La noción rigurosa de proceso
predecible en tiempo continuo se formaliza vía la
$\sigma$-álgebra generada por los procesos adaptados continuos por
la izquierda en $[0, T] \times \Omega$, y coincide para la clase
de los procesos simples con la condición~(i). No será necesario
trabajar con ella explícitamente, pues toda la teoría posterior se
extiende por densidad desde $\mathcal{S}$; véase
\cite{karatzas-shreve}.} introducida en la
Sección~\ref{subsec:predecibles_int_disc} para el caso discreto. El
valor del integrando durante el subintervalo se determina con la
información disponible al inicio del mismo, antes de observar el
incremento del browniano.

\medskip

\begin{definicion}[Integral estocástica de un proceso simple]
\label{def:integral_simple}
Sea $\phi \in \mathcal{S}$ un proceso simple con la representación
\eqref{eq:proceso_simple}. Su \emph{integral estocástica} (o
\emph{integral de Itô}) sobre $[0, T]$ es la variable aleatoria
\begin{equation}\label{eq:integral_simple}
  I(\phi) \eqdef \int_0^T \phi_s\, dW_s
            \eqdef \sum_{i=0}^{n-1} \xi_i\,
                   \bigl(W_{t_{i+1}} - W_{t_i}\bigr).
\end{equation}
\end{definicion}

La fórmula~\eqref{eq:integral_simple} es la versión continua de la
integral estocástica discreta $(H \cdot M)_n$ definida
en~\eqref{eq:integral_discreta}: cada incremento del browniano sobre
un subintervalo es ponderado por el valor que el integrando toma al
inicio de dicho subintervalo. Como los $\xi_i$ son variables aleatorias,
$I(\phi)$ es también una variable aleatoria, y por construcción es
$\F_T$-medible.

\medskip

\noindent\textbf{Buena definición y linealidad.} Antes de continuar,
conviene observar que la integral~\eqref{eq:integral_simple} no
depende de la representación particular utilizada para $\phi$:
refinando la partición, $\xi_i\, (W_{t_{i+1}} - W_{t_i}) = \xi_i\,
(W_s - W_{t_i}) + \xi_i\, (W_{t_{i+1}} - W_s)$, identidad trivial por
la aditividad de los incrementos del browniano, lo cual implica que
dos representaciones cualesquiera (para las que siempre existe un refinamiento común, la unión)
producen la misma integral. Asimismo, la integral es lineal: para
$\phi, \psi \in \mathcal{S}$ y $a, b \in \R$, $I(a\phi + b\psi) = a\,
I(\phi) + b\, I(\psi)$ casi seguramente, lo cual se sigue inmediatamente
de la fórmula~\eqref{eq:integral_simple} tras tomar un refinamiento
común de las particiones.

\subsection{La isometría de Itô}
\label{subsec:isometria_ito}

Establecemos a continuación la \emph{isometría de Itô}: una identidad
que controla la norma $L^2$ de la integral en términos del integrando.

\medskip

\begin{definicion}[Espacio $\mathcal{H}^2$]
\label{def:H2}
Denotamos por $\mathcal{H}^2[0, T]$, o simplemente $\mathcal{H}^2$
cuando el horizonte temporal sea claro por el contexto, el espacio de
los procesos estocásticos $\phi : [0, T] \times \Omega \to \R$ que
satisfacen:
\begin{enumerate}[label=\textnormal{(\roman*)}]
   \item $\phi$ es progresivamente medible\footnote{Es decir, para cada
      $t \in [0, T]$, la restricción de $\phi$ a $[0, t] \times
      \Omega$ es $\mathcal{B}([0, t]) \otimes \F_t$-medible. La
      propiedad implica adaptación y se cumple automáticamente para
      procesos adaptados con trayectorias continuas por la izquierda,
      en particular para los procesos simples de la
      Definición~\ref{def:proceso_simple}.} respecto a la filtración
      $(\F_t)$;
  \item $\displaystyle \E\!\left[\int_0^T \phi_s^2\, ds\right] < \infty.$
\end{enumerate}
Sobre $\mathcal{H}^2$ se considera la norma
\begin{equation}\label{eq:norma_H2}
  \|\phi\|_{\mathcal{H}^2}
  \eqdef \E\!\left[\int_0^T \phi_s^2\, ds\right]^{\!1/2}.
\end{equation}
\end{definicion}

La condición de progresiva medibilidad garantiza la medibilidad
conjunta en $(t, \omega)$ y, en particular, que la integral
$\int_0^T \phi_s^2\, ds$ esté bien definida. Para los procesos
simples (Definición~\ref{def:proceso_simple}) ambas condiciones se
cumplen trivialmente, y por tanto $\mathcal{S} \subseteq
\mathcal{H}^2$.

\medskip

\noindent\textbf{La isometría.}

\begin{teorema}[Isometría de Itô para procesos simples]
\label{teo:isometria_ito_simples}
Sea $\phi \in \mathcal{S}$ un proceso simple. Entonces
\begin{equation}\label{eq:isometria_ito_simples}
  \E\!\left[\bigl(I(\phi)\bigr)^2\right]
  = \E\!\left[\int_0^T \phi_s^2\, ds\right]
  = \|\phi\|_{\mathcal{H}^2}^2.
\end{equation}
Como consecuencia inmediata, $\E[I(\phi)] = 0$.
\end{teorema}

\begin{proof}
Sea $\phi \in \mathcal{S}$ con la representación
\eqref{eq:proceso_simple}, esto es, $\phi_t = \sum_{i=0}^{n-1} \xi_i\,
\mathbf{1}_{(t_i, t_{i+1}]}(t)$, con $\xi_i$ acotadas y
$\F_{t_i}$-medibles. Por la
Definición~\ref{def:integral_simple}, $I(\phi) = \sum_{i=0}^{n-1}
\xi_i\, (W_{t_{i+1}} - W_{t_i})$.

Para abreviar la notación, denotemos $\Delta W_i \eqdef W_{t_{i+1}} -
W_{t_i}$ y $\Delta_i \eqdef t_{i+1} - t_i$.

\emph{Paso 1: cálculo de $\E[I(\phi)^2]$.} Expandiendo el cuadrado de
la suma,
\begin{equation*}
  I(\phi)^2 = \sum_{i=0}^{n-1} \xi_i^2\, (\Delta W_i)^2
            + 2 \sum_{j < i} \xi_j\, \xi_i\, \Delta W_j\, \Delta W_i.
\end{equation*}
Tomando esperanza,
\begin{equation}\label{eq:expansion_cuadrado}
  \E[I(\phi)^2]
  = \sum_{i=0}^{n-1} \E\bigl[\xi_i^2\, (\Delta W_i)^2\bigr]
  + 2 \sum_{j < i} \E\bigl[\xi_j\, \xi_i\, \Delta W_j\, \Delta W_i\bigr].
\end{equation}

\emph{Paso 2: anulación de los términos cruzados.} Sea $j < i$. Las
variables $\xi_j$, $\xi_i$ y $\Delta W_j$ son todas
$\F_{t_i}$-medibles —en efecto, $\xi_j$ es $\F_{t_j}$-medible y
$\F_{t_j} \subseteq \F_{t_i}$, $\xi_i$ es $\F_{t_i}$-medible por
hipótesis, y $\Delta W_j = W_{t_{j+1}} - W_{t_j}$ es $\F_{t_i}$-medible
pues $t_{j+1} \leq t_i$—. Por la propiedad ``sacar lo conocido'' de la
esperanza condicional (Proposición~\ref{prop:propiedades_esperanza_cond},
propiedad (iii)),
\begin{equation*}
  \E\bigl[\xi_j\, \xi_i\, \Delta W_j\, \Delta W_i \,\big|\, \F_{t_i}\bigr]
  = \xi_j\, \xi_i\, \Delta W_j\, \E[\Delta W_i \mid \F_{t_i}].
\end{equation*}
Pero $\Delta W_i$ es independiente de $\F_{t_i}$ por la condición (iii)
de la Definición~\ref{def:browniano}, y tiene media cero, luego
$\E[\Delta W_i \mid \F_{t_i}] = 0$. La esperanza condicional del
producto es cero, y tomando esperanza incondicional, también lo es la
esperanza:
\begin{equation*}
  \E\bigl[\xi_j\, \xi_i\, \Delta W_j\, \Delta W_i\bigr] = 0,
  \qquad \text{para todo } j < i.
\end{equation*}
Los términos cruzados de~\eqref{eq:expansion_cuadrado} se anulan.

\emph{Paso 3: cálculo de los términos diagonales.} Para cada $i$, la
variable $\xi_i^2$ es $\F_{t_i}$-medible y $(\Delta W_i)^2$ es
independiente de $\F_{t_i}$. Aplicando de nuevo ``sacar lo conocido'' y
la independencia,
\begin{equation*}
  \E\bigl[\xi_i^2\, (\Delta W_i)^2\bigr]
  = \E\bigl[\xi_i^2\, \E[(\Delta W_i)^2 \mid \F_{t_i}]\bigr]
  = \E[\xi_i^2]\, \E[(\Delta W_i)^2]
  = \E[\xi_i^2]\, \Delta_i,
\end{equation*}
donde la última igualdad usa que $\Delta W_i \sim \mathcal{N}(0,
\Delta_i)$.

\emph{Paso 4: identificación con la norma $\mathcal{H}^2$.} Por
\eqref{eq:expansion_cuadrado} y los pasos anteriores,
\begin{equation*}
  \E[I(\phi)^2] = \sum_{i=0}^{n-1} \E[\xi_i^2]\, \Delta_i.
\end{equation*}
Por otra parte, sobre cada subintervalo $(t_i, t_{i+1}]$ el proceso
$\phi$ vale $\xi_i$, luego $\phi_s^2 = \xi_i^2$ y
\begin{equation*}
  \int_0^T \phi_s^2\, ds = \sum_{i=0}^{n-1} \xi_i^2\, \Delta_i.
\end{equation*}
Tomando esperanza,
\begin{equation*}
  \E\!\left[\int_0^T \phi_s^2\, ds\right]
  = \sum_{i=0}^{n-1} \E[\xi_i^2]\, \Delta_i,
\end{equation*}
que coincide con $\E[I(\phi)^2]$. La identidad
\eqref{eq:isometria_ito_simples} queda demostrada.

Para la consecuencia $\E[I(\phi)] = 0$, basta tomar esperanza en
$I(\phi) = \sum_i \xi_i\, \Delta W_i$ y aplicar el mismo argumento del
Paso 2 a cada sumando individual: $\E[\xi_i\, \Delta W_i] = \E[\xi_i\,
\E[\Delta W_i \mid \F_{t_i}]] = 0$.
\end{proof}

\medskip

\noindent\textbf{Comentario sobre el papel de la previsibilidad.} La demostración anterior pone de manifiesto con notable claridad la razón matemática por la que la elección de Itô, evaluar el integrando en el extremo izquierdo de cada subintervalo, es la correcta. El paso decisivo es el \emph{Paso 2}: la anulación de los
términos cruzados se obtiene precisamente porque $\xi_i$ es
$\F_{t_i}$-medible mientras que el incremento $\Delta W_i$ es
independiente de $\F_{t_i}$. Si hubiéramos elegido evaluar el
integrando en el extremo derecho, es decir, si las variables fuesen
$\F_{t_{i+1}}$-medibles, la independencia del incremento dejaría de
valer y los términos cruzados ya no se anularían. La isometría es,
pues, una propiedad específica de la integral de Itô, no compartida
por integrales construidas con otras convenciones.

\medskip

\noindent\textbf{Lectura como isometría.} La identidad
\eqref{eq:isometria_ito_simples} admite una lectura geométrica natural. Considerando $\mathcal{S}$ como subconjunto del
espacio de Hilbert $\mathcal{H}^2$ con la norma $\|\cdot\|_{\mathcal{H}^2}$
definida en~\eqref{eq:norma_H2}, y $L^2(\Omega, \F_T, \PP)$ como
espacio de Hilbert con su norma habitual, la
identidad~\eqref{eq:isometria_ito_simples} dice exactamente que el
operador $I : (\mathcal{S},\, \|\cdot\|_{\mathcal{H}^2}) \to
(L^2(\Omega),\, \|\cdot\|_{L^2})$ es lineal y preserva la norma. Esto
es exactamente la noción de \emph{isometría lineal} en el sentido del
análisis funcional.

\subsection{Extensión a $\mathcal{H}^2$ por densidad}
\label{subsec:extension_H2}

La isometría establecida en el Teorema~\ref{teo:isometria_ito_simples}
permite extender la integral desde $\mathcal{S}$ a $\mathcal{H}^2$
por el argumento clásico de densidad: aproximamos el integrando por
procesos simples y usamos la completitud de $L^2(\Omega)$ para pasar
al límite.

\medskip

\noindent\textbf{Densidad de los procesos simples.}

\begin{lema}[Densidad de $\mathcal{S}$ en $\mathcal{H}^2$]
\label{lema:densidad_simples}
La clase $\mathcal{S}$ de procesos simples es densa en $\mathcal{H}^2$
respecto a la norma $\|\cdot\|_{\mathcal{H}^2}$. Es decir, para todo
proceso $\phi \in \mathcal{H}^2$ existe una sucesión $(\phi^{(n)}) \subset
\mathcal{S}$ tal que
\begin{equation*}
  \|\phi^{(n)} - \phi\|_{\mathcal{H}^2} \xrightarrow{n \to \infty} 0.
\end{equation*}
\end{lema}

La demostración del lema, que omitimos, procede en dos etapas. En la
primera se reduce el problema al caso de procesos $\phi$ acotados, vía
truncación. En la segunda, dado un proceso acotado, este se aproxima
construyendo regularizaciones temporales de la forma
\begin{equation*}
  \phi^{(n)}_t \eqdef \sum_{i=0}^{n-1}
                       \biggl(\frac{n}{T} \int_{(i-1) T/n}^{i T/n} \phi_s\,
                              ds\biggr)
                       \mathbf{1}_{(i T/n,\, (i+1) T/n]}(t),
\end{equation*}
que son procesos simples y convergen a $\phi$ en $\mathcal{H}^2$ por
el teorema de continuidad de la traslación en $L^2$. El argumento
detallado puede consultarse en
\cite{karatzas-shreve}, sección 3.2.

\medskip

\noindent\textbf{Construcción de la integral sobre $\mathcal{H}^2$.}
Sea $\phi \in \mathcal{H}^2$. Por el Lema~\ref{lema:densidad_simples}
existe una sucesión $(\phi^{(n)}) \subset \mathcal{S}$ con
$\|\phi^{(n)} - \phi\|_{\mathcal{H}^2} \to 0$. En particular,
$(\phi^{(n)})$ es de Cauchy en $\mathcal{H}^2$. Aplicando la isometría
de Itô (Teorema~\ref{teo:isometria_ito_simples}) y la linealidad a la
diferencia $\phi^{(n)} - \phi^{(m)} \in \mathcal{S}$,
\begin{equation}\label{eq:cauchy_integral}
  \E\!\left[\bigl(I(\phi^{(n)}) - I(\phi^{(m)})\bigr)^{\!2}\right]
  = \E\!\left[\bigl(I(\phi^{(n)} - \phi^{(m)})\bigr)^{\!2}\right]
  = \|\phi^{(n)} - \phi^{(m)}\|_{\mathcal{H}^2}^2 \xrightarrow{n,m \to \infty} 0.
\end{equation}
La sucesión $(I(\phi^{(n)}))$ es, por tanto, de Cauchy en
$L^2(\Omega, \F_T, \PP)$. Como este espacio es completo, existe un
único elemento $L \in L^2(\Omega)$ tal que $I(\phi^{(n)}) \to L$ en
$L^2$.

Falta verificar que el límite $L$ no depende de la sucesión
aproximante particular. Sean $(\phi^{(n)}), (\widetilde{\phi}^{(n)})
\subset \mathcal{S}$ dos sucesiones tales que ambas convergen a $\phi$
en $\mathcal{H}^2$, con límites respectivos $L, \widetilde{L} \in
L^2(\Omega)$. Considerando la sucesión intercalada $\phi^{(1)},
\widetilde{\phi}^{(1)}, \phi^{(2)}, \widetilde{\phi}^{(2)}, \ldots$,
también convergente a $\phi$ en $\mathcal{H}^2$, las imágenes
$I(\phi^{(1)}), I(\widetilde{\phi}^{(1)}), I(\phi^{(2)}), \ldots$
forman una sucesión de Cauchy en $L^2(\Omega)$ por el mismo argumento
de~\eqref{eq:cauchy_integral}, y por tanto convergen a un único
límite. Pero las dos subsucesiones $(I(\phi^{(n)}))$ y
$(I(\widetilde{\phi}^{(n)}))$ deben converger a ese mismo límite, lo
cual fuerza $L = \widetilde{L}$ casi seguramente.

\begin{definicion}[Integral estocástica sobre $\mathcal{H}^2$]
\label{def:integral_H2}
Sea $\phi \in \mathcal{H}^2$ y $(\phi^{(n)}) \subset \mathcal{S}$ una
sucesión cualquiera de procesos simples con $\|\phi^{(n)} -
\phi\|_{\mathcal{H}^2} \to 0$. La \emph{integral estocástica de Itô}
de $\phi$ sobre $[0, T]$ es la variable aleatoria
\begin{equation*}
  I(\phi) \eqdef \int_0^T \phi_s\, dW_s
            \eqdef L^2\text{-}\!\lim_{n \to \infty}\, I(\phi^{(n)}),
\end{equation*}
elemento de $L^2(\Omega, \F_T, \PP)$ definido salvo modificaciones en
conjuntos de medida nula.
\end{definicion}

La definición es consistente: para $\phi \in \mathcal{S}$, tomando la
sucesión constante $\phi^{(n)} = \phi$ se recupera la
Definición~\ref{def:integral_simple}.

\begin{teorema}[Propiedades de la integral estocástica sobre $\mathcal{H}^2$]
\label{teo:propiedades_integral_H2}
El operador $I : \mathcal{H}^2 \to L^2(\Omega, \F_T, \PP)$ definido
por la Definición~\ref{def:integral_H2} satisface las siguientes
propiedades:
\begin{enumerate}[label=\textnormal{(\roman*)}]
  \item \emph{Linealidad:} para $\phi, \psi \in \mathcal{H}^2$ y $a, b
        \in \R$,
        \begin{equation*}
          I(a\phi + b\psi) = a\, I(\phi) + b\, I(\psi)
          \quad \text{c.s.}
        \end{equation*}
  \item \emph{Isometría de Itô:} para $\phi \in \mathcal{H}^2$,
        \begin{equation}\label{eq:isometria_H2}
          \E\!\left[\bigl(I(\phi)\bigr)^{\!2}\right]
          = \E\!\left[\int_0^T \phi_s^2\, ds\right]
          = \|\phi\|_{\mathcal{H}^2}^2.
        \end{equation}
  \item \emph{Esperanza nula:} para $\phi \in \mathcal{H}^2$,
        $\E[I(\phi)] = 0$.
\end{enumerate}
\end{teorema}

\begin{proof}
Las tres propiedades se siguen de las correspondientes propiedades
sobre $\mathcal{S}$ por paso al límite, gracias a que $I$ es continuo
respecto a las normas $\|\cdot\|_{\mathcal{H}^2}$ y $\|\cdot\|_{L^2}$
(continuidad que es a su vez consecuencia de la isometría sobre
$\mathcal{S}$).

A modo de ilustración, demostramos (ii). Sea $\phi \in \mathcal{H}^2$
y $(\phi^{(n)}) \subset \mathcal{S}$ con $\phi^{(n)} \to \phi$ en
$\mathcal{H}^2$. Entonces $I(\phi^{(n)}) \to I(\phi)$ en $L^2(\Omega)$
por la Definición~\ref{def:integral_H2}. La continuidad de la norma
$\|\cdot\|_{L^2}$ da
\begin{equation*}
  \E[I(\phi)^2]
  = \|I(\phi)\|_{L^2}^2
  = \lim_{n\to\infty} \|I(\phi^{(n)})\|_{L^2}^2
  = \lim_{n\to\infty} \E[I(\phi^{(n)})^2].
\end{equation*}
Por el Teorema~\ref{teo:isometria_ito_simples} aplicado a cada
$\phi^{(n)} \in \mathcal{S}$,
\begin{equation*}
  \E[I(\phi^{(n)})^2] = \|\phi^{(n)}\|_{\mathcal{H}^2}^2,
\end{equation*}
y por la continuidad de la norma $\|\cdot\|_{\mathcal{H}^2}$,
$\|\phi^{(n)}\|_{\mathcal{H}^2}^2 \to \|\phi\|_{\mathcal{H}^2}^2$. La
identidad~\eqref{eq:isometria_H2} se obtiene pasando al límite. Las
demostraciones de (i) y (iii) son enteramente análogas.
\end{proof}

La integral así definida es, sin embargo, un objeto \emph{global}
sobre $[0, T]$. En muchos contextos necesitaremos la integral sobre
$[0, t]$ con $t$ variable: esto motiva el estudio del proceso $\{M_t\}$
en la subsección siguiente.

\subsection{La integral indefinida como proceso}
\label{subsec:integral_indefinida}

Para cada $t \in [0, T]$, el proceso truncado $\phi\, \mathbf{1}_{[0,
t]}$ pertenece a $\mathcal{H}^2$ (la restricción de un proceso de
cuadrado integrable es de cuadrado integrable), lo que permite
aplicarle el operador integral $I$ y obtener así una familia
$\{M_t\}_{t \in [0, T]}$ que estudiaremos como proceso indexado por
$t$.

\begin{definicion}[Integral indefinida]
\label{def:integral_indefinida}
Sea $\phi \in \mathcal{H}^2$. Para cada $t \in [0, T]$, definimos
\begin{equation}\label{eq:def_integral_indefinida}
  M_t \eqdef \int_0^t \phi_s\, dW_s
        \eqdef I\!\left(\phi\, \mathbf{1}_{[0, t]}\right).
\end{equation}
La familia $\{M_t\}_{t \in [0, T]}$ así obtenida se denomina la
\emph{integral indefinida} de $\phi$ respecto al browniano $W$.
\end{definicion}

Observemos que para $t = T$ se recupera $M_T = I(\phi)$, el operador
global de la subsección anterior. Asimismo, $M_0 = I(\phi\,
\mathbf{1}_{\{0\}}) = 0$ casi seguramente, pues el proceso $\phi\,
\mathbf{1}_{\{0\}}$ tiene norma $\mathcal{H}^2$ nula.

\medskip

\noindent\textbf{Continuidad de las trayectorias.} Tal como hemos
construido el proceso, cada $M_t$ está definido como un elemento de
$L^2(\Omega)$, esto es, salvo modificaciones en conjuntos de medida
nula. Esta forma de definición no garantiza, por sí sola, que las
trayectorias $t \mapsto M_t(\omega)$ sean funciones continuas: la
familia $\{M_t\}_{t \in [0, T]}$ está indexada por un parámetro
continuo, y la unión no numerable de los conjuntos de medida nula
asociados a cada $M_t$ podría, \emph{a priori}, no ser de medida nula.
Es necesario por tanto verificar la existencia de una \emph{versión}
continua del proceso, esto es, una elección de los representantes en
cada clase de equivalencia que produzca trayectorias continuas casi
seguramente.

\begin{teorema}[Continuidad de la integral indefinida]
\label{teo:continuidad_integral_indefinida}
Sea $\phi \in \mathcal{H}^2$ y $\{M_t\}_{t \in [0, T]}$ el proceso
definido en~\eqref{eq:def_integral_indefinida}. Existe una versión de
este proceso cuyas trayectorias son continuas casi seguramente.
\end{teorema}

La demostración rigurosa, que omitimos, descansa esencialmente sobre la
desigualdad maximal de Doob para martingalas en $L^2$, junto con el
argumento de aproximación por procesos simples ya empleado. Para los
procesos simples, las trayectorias del proceso $M_t$ son continuas por
construcción (son lineales a trozos en $t$); el paso al límite por
densidad, controlado mediante la desigualdad maximal aplicada al
proceso diferencia de dos aproximaciones, garantiza convergencia
uniforme casi segura sobre $[0, T]$, lo cual preserva la continuidad.
La construcción detallada se encuentra en \cite{karatzas-shreve},
sección 3.2.

A partir de aquí, y en lo sucesivo, identificaremos siempre $\{M_t\}$
con su versión continua, sin mención explícita.

\medskip

\noindent\textbf{La propiedad de martingala.} 

\begin{teorema}[Propiedad de martingala de la integral indefinida]
\label{teo:martingala_integral_indefinida}
Sea $\phi \in \mathcal{H}^2$ y $\{M_t\}_{t \in [0, T]}$ la integral
indefinida de $\phi$. Entonces $\{M_t\}$ es una martingala respecto a
la filtración $(\F_t)$, esto es:
\begin{enumerate}[label=\textnormal{(\roman*)}]
  \item $M_t$ es $\F_t$-medible para todo $t \in [0, T]$;
  \item $\E[|M_t|] < \infty$ para todo $t \in [0, T]$;
  \item $\E[M_t \mid \F_s] = M_s$ casi seguramente, para todo $0 \leq s
        \leq t \leq T$.
\end{enumerate}
\end{teorema}

\begin{proof}
\emph{Adaptación e integrabilidad.} Para procesos simples, $M_t$ es
suma finita de productos de variables $\F_t$-medibles y por tanto
$\F_t$-medible. Para $\phi \in \mathcal{H}^2$ general, la convergencia
en $L^2$ de las aproximaciones $M_t^{(n)}$ a $M_t$
(Definición~\ref{def:integral_H2}) implica una subsucesión convergente
casi seguramente, y la $\F_t$-medibilidad se preserva bajo límites
casi seguros. La integrabilidad se sigue de Cauchy-Schwarz:
$\E[|M_t|] \leq \E[M_t^2]^{1/2} = \|\phi\, \mathbf{1}_{[0,
t]}\|_{\mathcal{H}^2} \leq \|\phi\|_{\mathcal{H}^2} < \infty$, donde
la igualdad central usa la isometría de Itô.

\emph{Propiedad de martingala.} Sean $0 \leq s \leq t \leq T$.
Comencemos por el caso de un proceso simple $\phi \in \mathcal{S}$ con
representación $\phi_r = \sum_i \xi_i\, \mathbf{1}_{(t_i, t_{i+1}]}(r)$.
Tomando un refinamiento de la partición que incluya los puntos $s$ y
$t$ —lo cual no altera el proceso $\phi$ ni la integral indefinida—,
podemos suponer sin pérdida de generalidad que $s = t_p$ y $t = t_q$
para ciertos $p \leq q$. Entonces
\begin{equation*}
  M_t - M_s = \int_s^t \phi_r\, dW_r
            = \sum_{i = p}^{q-1} \xi_i\,
              \bigl(W_{t_{i+1}} - W_{t_i}\bigr).
\end{equation*}
Tomando esperanza condicional respecto a $\F_s = \F_{t_p}$ y usando
linealidad,
\begin{equation*}
  \E[M_t - M_s \mid \F_s]
  = \sum_{i = p}^{q-1} \E\bigl[\xi_i\, (W_{t_{i+1}} - W_{t_i})
                               \,\big|\, \F_{t_p}\bigr].
\end{equation*}
Para cada sumando, observemos que $\F_{t_p} \subseteq \F_{t_i}$ pues $i
\geq p$. Aplicando la propiedad de la torre
(Proposición~\ref{prop:propiedades_esperanza_cond}, propiedad (iv))
con $\F_{t_p} \subseteq \F_{t_i}$,
\begin{equation*}
  \E\bigl[\xi_i\, (W_{t_{i+1}} - W_{t_i}) \mid \F_{t_p}\bigr]
  = \E\!\left[\E\bigl[\xi_i\, (W_{t_{i+1}} - W_{t_i})
                       \mid \F_{t_i}\bigr] \,\Big|\, \F_{t_p}\right].
\end{equation*}
La esperanza condicional interna es nula: $\xi_i$ es $\F_{t_i}$-medible
y $W_{t_{i+1}} - W_{t_i}$ es independiente de $\F_{t_i}$ con media
cero, luego sacando lo conocido y usando la independencia,
$\E[\xi_i\, (W_{t_{i+1}} - W_{t_i}) \mid \F_{t_i}] = \xi_i\,
\E[W_{t_{i+1}} - W_{t_i} \mid \F_{t_i}] = 0$. Por consiguiente,
$\E[M_t - M_s \mid \F_s] = 0$, es decir, $\E[M_t \mid \F_s] = M_s$
casi seguramente. La propiedad de martingala queda establecida para
procesos simples.

Para $\phi \in \mathcal{H}^2$ general, tomemos una sucesión
$(\phi^{(n)}) \subset \mathcal{S}$ con $\phi^{(n)} \to \phi$ en
$\mathcal{H}^2$, y sean $M_t^{(n)}$ las correspondientes integrales
indefinidas. Por lo recién demostrado, $\E[M_t^{(n)} \mid \F_s] =
M_s^{(n)}$ casi seguramente. Tomando límite en $L^2(\Omega)$ cuando $n
\to \infty$, $M_t^{(n)} \to M_t$ y $M_s^{(n)} \to M_s$ por la
isometría aplicada a las restricciones a $[0, t]$ y $[0, s]$
respectivamente. Como la esperanza condicional respecto a $\F_s$ es un
operador continuo en $L^2$ —es la proyección ortogonal sobre el
subespacio $L^2(\Omega, \F_s, \PP)$—, la igualdad pasa al límite:
$\E[M_t \mid \F_s] = M_s$ casi seguramente.
\end{proof}

\medskip

\noindent\textbf{La isometría en versión proceso.} Como consecuencia
inmediata de la construcción, la isometría de Itô se generaliza al
proceso indefinido: para todo $t \in [0, T]$,
\begin{equation}\label{eq:isometria_proceso}
  \E[M_t^2] = \E\!\left[\int_0^t \phi_s^2\, ds\right].
\end{equation}
La identidad se sigue aplicando el
Teorema~\ref{teo:propiedades_integral_H2} al proceso restringido
$\phi\, \mathbf{1}_{[0, t]} \in \mathcal{H}^2$. 

\section{La fórmula de Itô}
\label{sec:formula_ito}

Esta sección desarrolla la \emph{fórmula de Itô}, sustituto continuo
de la regla de la cadena clásica y, junto con el teorema de Girsanov,
una de las piezas fundamentales del Capítulo~\ref{cap:bs}.

\subsection{Procesos de Itô y la fórmula de Itô para $f(W_t)$}
\label{subsec:procesos_ito_y_formula}

\begin{definicion}[Proceso de Itô]
\label{def:proceso_ito}
Un proceso estocástico $\{X_t\}_{t \in [0, T]}$ es un \emph{proceso de
Itô} si admite una representación de la forma
\begin{equation}\label{eq:def_proceso_ito}
  X_t = X_0 + \int_0^t a_s\, ds + \int_0^t b_s\, dW_s,
  \qquad t \in [0, T],
\end{equation}
donde $X_0$ es una variable aleatoria $\F_0$-medible, y los procesos
$\{a_t\}$ y $\{b_t\}$ son adaptados a $(\F_t)$ y satisfacen las
condiciones de integrabilidad
\begin{equation*}
  \int_0^T |a_s|\, ds < \infty \text{ casi seguramente},
  \qquad b \in \mathcal{H}^2[0, T].
\end{equation*}
Los procesos $a$ y $b$ se denominan, respectivamente, el
\emph{coeficiente de deriva} y el \emph{coeficiente de difusión} del
proceso $X$.
\end{definicion}

La condición $b \in \mathcal{H}^2[0, T]$ garantiza que la integral
estocástica $\int_0^t b_s\, dW_s$ está bien definida en el sentido del
Teorema~\ref{teo:propiedades_integral_H2}, mientras que la condición
sobre $a$ asegura que la integral $\int_0^t a_s\, ds$ está bien
definida trayectorialmente como integral de Lebesgue.

\medskip

\noindent\textbf{Notación diferencial.} La igualdad
integral~\eqref{eq:def_proceso_ito} se escribe habitualmente, en forma
abreviada, como
\begin{equation}\label{eq:notacion_diferencial}
  dX_t = a_t\, dt + b_t\, dW_t.
\end{equation}
Esta notación es muy cómoda y se utilizará a partir de ahora con
frecuencia, pero conviene insistir en que carece de significado
matemático propio: no es más que una forma compacta de expresar \eqref{eq:def_proceso_ito}. Las identidades de la
forma~\eqref{eq:notacion_diferencial} deben entenderse siempre como
abreviatura de las correspondientes identidades integrales.

\medskip

\noindent\textbf{Enunciado de la fórmula de Itô.}

\begin{teorema}[Fórmula de Itô para $f(W_t)$]
\label{teo:formula_ito_brown}
Sea $f : \R \to \R$ una función de clase $C^2$ con $f$, $f'$ y $f''$
acotadas. Entonces, para todo $t \in [0, T]$, casi seguramente,
\begin{equation}\label{eq:formula_ito_brown}
  f(W_t) = f(0) + \int_0^t f'(W_s)\, dW_s
        + \frac{1}{2} \int_0^t f''(W_s)\, ds.
\end{equation}
Equivalentemente, en notación diferencial,
\begin{equation*}
  d\bigl(f(W_t)\bigr) = f'(W_t)\, dW_t + \tfrac{1}{2} f''(W_t)\, dt.
\end{equation*}
\end{teorema}

\noindent\textbf{Lectura del resultado.} La diferencial de $f(W_t)$
es suma de un término en $dW_t$, análogo al de la regla de la cadena
clásica, y un término extra en $dt$, el \emph{término de Itô}, ausente en el cálculo ordinario. Esta corrección es fruto de la regla heurística $(dW_t)^2 = dt$ adelantada
en~\eqref{eq:regla_dW2}: al desarrollar $f$ en serie de Taylor, el
segundo orden no se desvanece y aporta exactamente $\tfrac{1}{2}
f''(W_t)\, dt$. Las hipótesis de acotación son las técnicamente más
cómodas; la fórmula vale bajo condiciones más débiles vía
localización.

\begin{proof}
Para cada partición $\pi = \{0 = t_0 < t_1 < \cdots < t_n = t\}$ de
$[0, t]$, escribimos la expansión telescópica
\begin{equation}\label{eq:telescopica}
  f(W_t) - f(0)
  = \sum_{k=0}^{n-1} \bigl[f(W_{t_{k+1}}) - f(W_{t_k})\bigr],
\end{equation}
los términos resultantes producirán las dos integrales del lado
derecho de~\eqref{eq:formula_ito_brown}, más un resto despreciable.

\emph{Paso 1: desarrollo de Taylor.} Por la hipótesis $f \in C^2$, para
cada $k$ existe $\theta_k \in [W_{t_k}, W_{t_{k+1}}]$ (o el intervalo
correspondiente con extremos invertidos) tal que
\begin{equation*}
  f(W_{t_{k+1}}) - f(W_{t_k})
  = f'(W_{t_k})\, \Delta W_k
   + \tfrac{1}{2} f''(W_{t_k})\, (\Delta W_k)^2
   + R_k,
\end{equation*}
donde $\Delta W_k \eqdef W_{t_{k+1}} - W_{t_k}$ y el resto se escribe
\begin{equation*}
  R_k = \tfrac{1}{2}\bigl[f''(\theta_k) - f''(W_{t_k})\bigr]\,
        (\Delta W_k)^2.
\end{equation*}
Sumando sobre $k$ y comparando con~\eqref{eq:telescopica},
\begin{equation}\label{eq:descomposicion_telescopica}
  f(W_t) - f(0)
  = \underbrace{\sum_{k=0}^{n-1} f'(W_{t_k})\, \Delta W_k}_{S_1^\pi}
   + \tfrac{1}{2} \underbrace{\sum_{k=0}^{n-1} f''(W_{t_k})\,
                              (\Delta W_k)^2}_{S_2^\pi}
   + \underbrace{\sum_{k=0}^{n-1} R_k}_{R^\pi}.
\end{equation}

\emph{Paso 2: convergencia de $S_1^\pi$.} La suma $S_1^\pi$ es
exactamente la integral estocástica del proceso simple
$\phi^\pi_s \eqdef \sum_k f'(W_{t_k})\, \mathbf{1}_{(t_k, t_{k+1}]}(s)$
respecto a $W$ sobre $[0, t]$. Veamos que $\phi^\pi$ converge a
$f'(W_s)$ en $\mathcal{H}^2[0, t]$ cuando $|\pi| \to 0$. Por la
definición de la norma,
\begin{equation*}
  \|\phi^\pi - f'(W)\|_{\mathcal{H}^2}^2
  = \E\!\left[\int_0^t \bigl(\phi^\pi_s - f'(W_s)\bigr)^{\!2}\, ds\right]
  = \E\!\left[\sum_{k=0}^{n-1} \int_{t_k}^{t_{k+1}}
             \bigl(f'(W_{t_k}) - f'(W_s)\bigr)^{\!2}\, ds\right].
\end{equation*}
Por la continuidad uniforme de $f'$ sobre cualquier compacto y por la
continuidad casi segura de las trayectorias del browniano, el
integrando puntual tiende a cero cuando $|\pi| \to 0$; por la acotación
de $f'$, está dominado por una constante. El teorema de la convergencia
dominada garantiza que la integral converge a cero, esto es, $\phi^\pi
\to f'(W)$ en $\mathcal{H}^2$. Por la isometría de Itô,
\begin{equation*}
  S_1^\pi \xrightarrow{L^2} \int_0^t f'(W_s)\, dW_s.
\end{equation*}

\emph{Paso 3: convergencia de $S_2^\pi$.} Demostraremos que
\begin{equation*}
  S_2^\pi \xrightarrow{L^2} \int_0^t f''(W_s)\, ds
  \quad \text{cuando } |\pi| \to 0.
\end{equation*}
Para ello, descompongamos
\begin{equation*}
  S_2^\pi - \int_0^t f''(W_s)\, ds
  = \underbrace{\sum_{k=0}^{n-1} f''(W_{t_k})\,
                \bigl[(\Delta W_k)^2 - \Delta_k\bigr]}_{A^\pi}
  + \underbrace{\sum_{k=0}^{n-1} f''(W_{t_k})\, \Delta_k
              - \int_0^t f''(W_s)\, ds}_{B^\pi},
\end{equation*}
donde $\Delta_k \eqdef t_{k+1} - t_k$, y demostraremos que tanto
$A^\pi$ como $B^\pi$ tienden a cero en $L^2$.

Para $B^\pi$: la cantidad $\sum_k f''(W_{t_k})\, \Delta_k$ es la suma
de Riemann de la función continua $s \mapsto f''(W_s)$ asociada a la
partición $\pi$ y a la elección $\tau_k = t_k$. Por la continuidad
trayectorial casi segura, converge a $\int_0^t f''(W_s)\, ds$
trayectorialmente cuando $|\pi| \to 0$, y la acotación uniforme de
$f''$ permite aplicar convergencia dominada para obtener convergencia
en $L^2$.

Para $A^\pi$: por un argumento análogo al de la isometría de Itô
(Teorema~\ref{teo:isometria_ito_simples}, Paso 2), usando que
$f''(W_{t_k})$ es $\F_{t_k}$-medible y que $(\Delta W_k)^2 - \Delta_k$
tiene esperanza condicional nula bajo $\F_{t_k}$, los términos
cruzados se anulan al elevar al cuadrado y tomar esperanza. Aplicando
la cota uniforme $|f''| \leq M$ y el cálculo del cuarto momento
gaussiano $\E[((\Delta W_k)^2 - \Delta_k)^2] = 2\, \Delta_k^2$,
\begin{equation*}
  \E[(A^\pi)^2]
  = \sum_{k=0}^{n-1} \E\bigl[(f''(W_{t_k}))^2\bigr]\,
    \E\bigl[((\Delta W_k)^2 - \Delta_k)^2\bigr]
  \leq 2 M^2 \sum_{k=0}^{n-1} \Delta_k^2
  \leq 2 M^2 \cdot |\pi| \cdot t \xrightarrow{|\pi| \to 0} 0.
\end{equation*}

\emph{Paso 4: el resto $R^\pi$ es despreciable.} Recordemos que
\begin{equation*}
  R^\pi = \tfrac{1}{2} \sum_{k=0}^{n-1}
            \bigl[f''(\theta_k) - f''(W_{t_k})\bigr]\, (\Delta W_k)^2,
\end{equation*}
con $\theta_k$ entre $W_{t_k}$ y $W_{t_{k+1}}$. Como $f''$ es uniformemente
continua sobre cualquier compacto y las trayectorias del browniano son
continuas casi seguramente, el factor $|f''(\theta_k) - f''(W_{t_k})|$
tiende uniformemente a cero cuando $|\pi| \to 0$, esto es,
\begin{equation*}
  \varepsilon^\pi \eqdef \max_{0 \leq k \leq n-1}
                          |f''(\theta_k) - f''(W_{t_k})|
                          \xrightarrow{|\pi| \to 0} 0
  \quad \text{casi seguramente}.
\end{equation*}
Por tanto,
\begin{equation*}
  |R^\pi| \leq \tfrac{1}{2}\, \varepsilon^\pi
              \sum_{k=0}^{n-1} (\Delta W_k)^2.
\end{equation*}
La suma del lado derecho converge a $t$ en $L^2$ por el
Teorema~\ref{teo:variacion_cuadratica}, y $\varepsilon^\pi \to 0$ casi
seguramente con cota uniforme $2M$. Combinando ambos hechos (vía
convergencia dominada) se concluye $R^\pi \to 0$ en $L^1$, y por
tanto en probabilidad. Pasando si es preciso a una subsucesión, $R^\pi
\to 0$ casi seguramente.

\emph{Paso 5: conclusión.} Volviendo a la
descomposición~\eqref{eq:descomposicion_telescopica} y combinando los
límites obtenidos en los pasos anteriores (tras pasar a una
subsucesión común que asegure todas las convergencias casi
seguramente),
\begin{equation*}
  f(W_t) - f(0) = \int_0^t f'(W_s)\, dW_s
                + \tfrac{1}{2} \int_0^t f''(W_s)\, ds
                + 0,
\end{equation*}
casi seguramente, que es la identidad~\eqref{eq:formula_ito_brown}.
\end{proof}

\medskip

\subsection{Versión general y ejemplos de aplicación}
\label{subsec:formula_ito_general}

\medskip

\begin{teorema}[Fórmula de Itô para procesos de Itô]
\label{teo:formula_ito_general}
Sea $\{X_t\}$ un proceso de Itô con dinámica $dX_t = a_t\, dt + b_t\,
dW_t$, y sea $f : [0, T] \times \R \to \R$ una función de clase
$C^{1,2}$, esto es, continuamente diferenciable en su primera
variable y dos veces continuamente diferenciable en la segunda, con
las derivadas pertinentes acotadas sobre el rango de $X$. Entonces el
proceso $\{f(t, X_t)\}_{t \in [0, T]}$ es también un proceso de Itô,
y satisface
\begin{equation}\label{eq:formula_ito_general}
  d\bigl(f(t, X_t)\bigr)
  = \biggl[\partial_t f(t, X_t) + a_t\, \partial_x f(t, X_t)
           + \tfrac{1}{2}\, b_t^2\, \partial_{xx} f(t, X_t)\biggr]\, dt
   + b_t\, \partial_x f(t, X_t)\, dW_t.
\end{equation}
Equivalentemente, en forma integral,
\begin{multline*}
  f(t, X_t) = f(0, X_0)
      + \int_0^t \!\Bigl[\partial_t f(s, X_s) + a_s\, \partial_x f(s, X_s)
                          + \tfrac{1}{2}\, b_s^2\, \partial_{xx} f(s, X_s)\Bigr]
                          ds \\
      + \int_0^t b_s\, \partial_x f(s, X_s)\, dW_s.
\end{multline*}
\end{teorema}

La demostración de este teorema, que omitimos, sigue la misma
estrategia que la del Teorema~\ref{teo:formula_ito_brown}: expansión
telescópica seguida de un desarrollo de Taylor de orden 2 en la variable
espacial, ahora con un término adicional de orden 1 en la variable
temporal. La identificación de los límites es enteramente análoga; los
detalles pueden consultarse en \cite{karatzas-shreve}, sección 3.3.

\medskip

\noindent\textbf{Las reglas mnemónicas de Itô.} La estructura
de~\eqref{eq:formula_ito_general} se memoriza fácilmente operando con
$dX_t = a_t\, dt + b_t\, dW_t$ como con un objeto diferencial
ordinario, salvo por las reglas de multiplicación
\begin{equation}\label{eq:reglas_ito}
  (dt)^2 = 0, \qquad dt\, dW_t = 0, \qquad (dW_t)^2 = dt,
\end{equation}
que codifican respectivamente que $dt$ es de orden superior a
$\sqrt{dt}$ y la variación cuadrática del browniano del
Teorema~\ref{teo:variacion_cuadratica}. Aplicando un desarrollo de
Taylor formal a $f(t, X_t)$ y usando estas reglas para descartar
términos despreciables, se recupera exactamente
\eqref{eq:formula_ito_general}.

\medskip

\noindent\textbf{Tres ejemplos.} Ilustramos el uso de la fórmula con
tres ejemplos que reaparecerán en lo sucesivo.

\medskip

\noindent\textit{Ejemplo 1: el cálculo de $\int_0^t W_s\, dW_s$.}
Aplicamos el Teorema~\ref{teo:formula_ito_brown} a la función $f(x) =
\tfrac{1}{2} x^2$, cuyas derivadas son $f'(x) = x$ y $f''(x) = 1$.
Estrictamente esta función no satisface las hipótesis de acotación,
pero el resultado se mantiene por argumentos de localización. La fórmula da
\begin{equation*}
  \tfrac{1}{2} W_t^2 = 0 + \int_0^t W_s\, dW_s
                       + \tfrac{1}{2} \int_0^t 1\, ds,
\end{equation*}
es decir,
\begin{equation}\label{eq:integral_W_dW}
  \int_0^t W_s\, dW_s = \tfrac{1}{2}\bigl(W_t^2 - t\bigr).
\end{equation}
Recuperamos así, por la vía sistemática del cálculo de Itô, el
resultado que en la motivación
(Sección~\ref{sec:motivacion_ito}) habíamos obtenido por el cálculo
directo de las sumas. La fórmula~\eqref{eq:integral_W_dW} muestra
también, vía el Teorema~\ref{teo:martingala_integral_indefinida}, que
$\{W_t^2 - t\}$ es una martingala (resultado que enuncia el papel
de esta cantidad ya anticipado en la
Proposición~\ref{prop:martingalas_brownianas}).

\medskip

\noindent\textit{Ejemplo 2: el movimiento browniano geométrico.}
Consideremos el proceso
\begin{equation*}
  S_t \eqdef S_0\, \exp\bigl(\nu t + \sigma\, W_t\bigr),
\end{equation*}
donde $S_0 > 0$, $\nu \in \R$ y $\sigma > 0$. Aplicamos el
Teorema~\ref{teo:formula_ito_general} a la función $f(t, x) = S_0\,
\exp(\nu t + \sigma x)$ y al proceso $X_t = W_t$ (que es proceso de Itô
con $a \equiv 0$, $b \equiv 1$). Las derivadas de $f$ son
\begin{equation*}
  \partial_t f = \nu\, f, \qquad
  \partial_x f = \sigma\, f, \qquad
  \partial_{xx} f = \sigma^2\, f.
\end{equation*}
Sustituyendo en la fórmula,
\begin{equation*}
  dS_t = \bigl[\nu\, S_t + 0 + \tfrac{1}{2} \sigma^2\, S_t\bigr]\, dt
        + \sigma\, S_t\, dW_t,
\end{equation*}
es decir,
\begin{equation}\label{eq:gbm_sde}
  \boxed{\; dS_t = \mu\, S_t\, dt + \sigma\, S_t\, dW_t, \qquad
  \mu \eqdef \nu + \tfrac{1}{2}\sigma^2. \;}
\end{equation}
La identidad~\eqref{eq:gbm_sde} es la \emph{ecuación diferencial
estocástica del movimiento browniano geométrico}, modelo canónico
del precio del subyacente en Black-Scholes
(Capítulo~\ref{cap:bs}).

\medskip

\noindent\textit{Ejemplo 3: la martingala exponencial generalizada.}
Sea $\{\theta_t\}_{t \in [0, T]}$ un proceso adaptado con $\theta \in
\mathcal{H}^2$. Consideremos el proceso $\{L_t\}$ definido por
\begin{equation}\label{eq:martingala_exp_general}
  L_t \eqdef \exp\!\biggl(- \int_0^t \theta_s\, dW_s
                           - \tfrac{1}{2} \int_0^t \theta_s^2\, ds\biggr).
\end{equation}
Aplicamos la fórmula de Itô para identificar su SDE. Definamos $Y_t
\eqdef -\int_0^t \theta_s\, dW_s - \tfrac{1}{2} \int_0^t \theta_s^2\,
ds$, que es un proceso de Itô con $dY_t = -\tfrac{1}{2} \theta_t^2\,
dt - \theta_t\, dW_t$. El proceso $L_t = \exp(Y_t)$ es función del
proceso de Itô $Y$ vía $f(y) = e^y$, con $f'(y) = f''(y) = e^y = L_t$.
Por el Teorema~\ref{teo:formula_ito_general} aplicado a $f(Y_t)$, donde el coeficiente de deriva de $Y$ es $-\tfrac{1}{2}\theta_t^2$ y
el de difusión es $-\theta_t$,
\begin{equation*}
  dL_t = L_t\, \biggl[-\tfrac{1}{2}\theta_t^2 +
                       \tfrac{1}{2}\, (-\theta_t)^2\biggr]\, dt
        + L_t\, (-\theta_t)\, dW_t.
\end{equation*}
Los dos términos del corchete cancelan exactamente, y queda
\begin{equation}\label{eq:sde_martingala_exp}
  dL_t = -\theta_t\, L_t\, dW_t.
\end{equation}
La cancelación que produce el término de deriva nulo no es accidental:
la corrección $-\tfrac{1}{2}\int \theta_s^2\, ds$ que aparece en el
exponente de~\eqref{eq:martingala_exp_general} está calibrada
precisamente para que esto ocurra. La consecuencia es que $\{L_t\}$ es
un proceso sin deriva, expresable como integral estocástica de
$-\theta_t\, L_t$ respecto a $W$. Por el
Teorema~\ref{teo:martingala_integral_indefinida}, si $\theta\, L \in
\mathcal{H}^2$, entonces $\{L_t\}$ es una martingala.

El proceso~\eqref{eq:martingala_exp_general} generaliza la
martingala exponencial $\mathcal{E}_t^\lambda$ del
Capítulo~\ref{cap:brown}, sustituyendo el parámetro $\lambda$ por un
proceso adaptado $\theta_t$. Reaparecerá en
§\ref{sec:girsanov} como derivada de Radon-Nikodym del cambio
de medida de Girsanov.

\medskip

\section{El teorema de Girsanov}
\label{sec:girsanov}

En el modelo binomial del Capítulo~\ref{cap:discreto}, el cambio
$\PP \to \Q$ se construyó explícitamente modificando las
probabilidades del random walk. En el contexto continuo, lo que se
modifica es la medida sobre el espacio de trayectorias, y la
pregunta es bajo qué medida $\Q$ un proceso con deriva $\theta_t$
respecto a $W$ se convierte en martingala. La respuesta la
proporciona el \emph{teorema de Girsanov}.

\subsection{Cambio de medida en tiempo continuo}
\label{subsec:cambio_medida}

Antes de enunciar el teorema, conviene introducir el lenguaje y los
conceptos sobre los que descansará: equivalencia de medidas y derivada
de Radon-Nikodym en el contexto continuo.

\medskip

\begin{definicion}[Medidas equivalentes]
\label{def:medidas_equivalentes}
Dos medidas de probabilidad $\PP$ y $\Q$ sobre un espacio medible
$(\Omega, \F)$ se dicen \emph{equivalentes}, y se denota $\Q \sim \PP$,
si para todo $A \in \F$,
\begin{equation*}
  \PP(A) = 0 \iff \Q(A) = 0.
\end{equation*}
\end{definicion}

La equivalencia de medidas es el grado mínimo de compatibilidad que
necesitaremos: dos medidas equivalentes asignan probabilidad nula
exactamente a los mismos sucesos, aunque pueden valorar de manera
distinta los demás. En particular, las afirmaciones del tipo ``casi
seguramente'' son intercambiables entre las dos.

\medskip

\noindent\textbf{Derivada de Radon-Nikodym.}

\begin{teorema}[Radon-Nikodym, versión equivalente]
\label{teo:radon_nikodym}
Sean $\PP$ y $\Q$ dos medidas de probabilidad equivalentes sobre
$(\Omega, \F)$. Existe una única (salvo igualdad casi segura) variable
aleatoria $L : \Omega \to (0, \infty)$ que es $\F$-medible,
$\PP$-integrable, y satisface
\begin{equation}\label{eq:def_RN}
  \Q(A) = \int_A L\, d\PP, \qquad A \in \F.
\end{equation}
La variable $L$ se denomina \emph{derivada de Radon-Nikodym de $\Q$
respecto a $\PP$}, y se denota $L = d\Q/d\PP$.
\end{teorema}

La demostración puede consultarse en \cite{williams}.
La identidad~\eqref{eq:def_RN} se puede reformular, equivalentemente, en
términos de esperanzas: para cualquier variable aleatoria $X$ tal que
$L X$ sea $\PP$-integrable,
\begin{equation}\label{eq:cambio_esperanza_RN}
  \E_\Q[X] = \E_\PP[L X],
\end{equation}
fórmula que determina cómo se calculan esperanzas bajo $\Q$ a partir de
esperanzas bajo $\PP$.

\begin{observacion}[El cambio de medida como método general]
\label{obs:cambio_medida_general}
La identidad~\eqref{eq:cambio_esperanza_RN} nos muestra una idea que va más allá del marco de la valoración: permite calcular esperanzas bajo una
medida $\Q$ operando bajo otra medida $\PP$, sin más que reponderar
mediante la densidad $L$. Sustituir una medida por otra equivalente para
convertir un cálculo difícil en uno asequible es un recurso de uso
general en probabilidad; en simulación, sin ir más lejos, es el
fundamento del \emph{importance sampling}, que reduce la varianza de las
estimaciones Monte Carlo de sucesos poco probables. El teorema de
Girsanov, que abordamos a continuación, es la versión de este método
adaptada al movimiento browniano.
\end{observacion}

\medskip

\noindent\textbf{El caso filtrado: la densidad como proceso.} En el
contexto que nos interesa, recordemos que $(\Omega, \F)$ está dotado además de una
filtración $(\F_t)_{t \in [0, T]}$, la filtración browniana aumentada. Si $\PP$
y $\Q$ son equivalentes sobre $(\Omega, \F_T)$, podemos restringirlas
a cada $\F_t$ con $t < T$, y obtener así una familia de derivadas de
Radon-Nikodym
\begin{equation}\label{eq:proceso_densidad}
  L_t \eqdef \frac{d(\Q\!\restriction_{\F_t})}{d(\PP\!\restriction_{\F_t})},
  \qquad t \in [0, T].
\end{equation}
La derivada de Radon-Nikodym se convierte de esta forma en un \emph{proceso} en
lugar de una variable aleatoria. La propiedad estructural fundamental
de este proceso es la siguiente.

\begin{proposicion}[Propiedad de martingala del proceso densidad]
\label{prop:martingala_densidad}
Sean $\PP$ y $\Q$ medidas equivalentes sobre $(\Omega, \F_T)$, y
$\{L_t\}_{t \in [0, T]}$ el proceso definido en~\eqref{eq:proceso_densidad}.
Entonces $\{L_t\}$ es una $\PP$-martingala respecto a $(\F_t)$ con
$\E_\PP[L_t] = 1$ para todo $t$.
\end{proposicion}

\begin{proof}
Adaptación e integrabilidad son inmediatas: $L_t$ es $\F_t$-medible
por construcción, y $\E_\PP[L_t] = \Q(\Omega) = 1$ por
\eqref{eq:def_RN} aplicada con $A = \Omega$.

Para la propiedad de martingala, sean $0 \leq s < t \leq T$ y $A \in
\F_s$. Como $\F_s \subseteq \F_t$, también $A \in \F_t$. Aplicando
\eqref{eq:def_RN} dos veces,
\begin{equation*}
  \int_A L_t\, d\PP = \Q(A) = \int_A L_s\, d\PP.
\end{equation*}
Por la unicidad de la esperanza condicional
(Definición~\ref{def:esperanza_condicional}), $\E_\PP[L_t \mid \F_s] =
L_s$ casi seguramente.
\end{proof}

\medskip

\noindent\textbf{Construcción de medidas equivalentes vía martingalas
exponenciales.} Nosotros seguiremos el camino inverso: dada una
$\PP$-martingala positiva $\{L_t\}$ con $\E_\PP[L_T] = 1$, podemos
\emph{definir} una nueva medida $\Q$ sobre $(\Omega, \F_T)$ mediante
\begin{equation}\label{eq:def_Q_continuo}
  \Q(A) \eqdef \int_A L_T\, d\PP, \qquad A \in \F_T.
\end{equation}
La integrabilidad y la condición $\E_\PP[L_T] = 1$ garantizan que $\Q$
es una medida de probabilidad, y la positividad estricta de
$L_T$ asegura $\Q \sim \PP$. La derivada de Radon-Nikodym es entonces
$d\Q/d\PP = L_T$, y el proceso densidad recupera, por la
Proposición~\ref{prop:martingala_densidad}, la martingala $\{L_t\}$ de
partida.

La pregunta, por tanto, es: ¿qué clase de martingalas
positivas con valor uno en $t = 0$ podemos construir? El
Ejemplo~3 de la Sección~\ref{subsec:formula_ito_general} ofrece la
respuesta natural: los procesos exponenciales
\begin{equation}\label{eq:L_t_recordada}
  L_t = \exp\!\biggl(-\int_0^t \theta_s\, dW_s
                      - \tfrac{1}{2} \int_0^t \theta_s^2\, ds\biggr),
\end{equation}
asociados a un proceso adaptado $\theta$, son candidatos canónicos. Por
construcción son positivos y satisfacen $L_0 = 1$. Hemos verificado
también que satisfacen la SDE $dL_t = -\theta_t\, L_t\, dW_t$, esto es,
no tienen deriva. Por el Teorema~\ref{teo:martingala_integral_indefinida},
si $\theta\, L \in \mathcal{H}^2$, entonces $\{L_t\}$ es una
$\PP$-martingala. Y en tal caso $\E_\PP[L_T] = \E_\PP[L_0] = 1$, como
exige~\eqref{eq:def_Q_continuo}.

\medskip

\noindent\textbf{La condición de Novikov.} La condición $\theta\, L \in
\mathcal{H}^2$ es delicada de verificar en la práctica, pues involucra
al propio proceso $L$ que estamos definiendo. Existe un criterio
estándar más manejable, conocido como \emph{condición de Novikov},
que garantiza que el proceso~\eqref{eq:L_t_recordada} es efectivamente
una martingala.

\begin{teorema}[Condición de Novikov]
\label{teo:novikov}
Sea $\{\theta_t\}_{t \in [0, T]}$ un proceso adaptado tal que
\begin{equation}\label{eq:novikov}
  \E_\PP\!\left[\exp\!\biggl(\tfrac{1}{2} \int_0^T \theta_s^2\, ds\biggr)\right]
  < \infty.
\end{equation}
Entonces el proceso $\{L_t\}$ definido en~\eqref{eq:L_t_recordada} es
una $\PP$-martingala.
\end{teorema}

La demostración de este resultado escapa del
ámbito del presente trabajo y puede consultarse en
\cite{karatzas-shreve}, sección 3.5. La utilizaremos como herramienta
para garantizar, en la práctica, la condición de martingala. En las
aplicaciones que abordaremos en el Capítulo~\ref{cap:bs}, el proceso
$\theta$ resultará ser determinista (de hecho constante), y la
condición de Novikov se verificará de manera trivial.

\subsection{El teorema de Girsanov}
\label{subsec:girsanov_teorema}

\noindent\textbf{La caracterización de Lévy.}

\begin{teorema}[Caracterización de Lévy del movimiento browniano]
\label{teo:caracterizacion_levy}
Sea $\{X_t\}_{t \in [0, T]}$ un proceso estocástico continuo y
adaptado a una filtración $(\F_t)$ sobre un espacio de probabilidad
$(\Omega, \F, \PP)$, con $X_0 = 0$ casi seguramente.
Si:
\begin{enumerate}[label=\textnormal{(\roman*)}]
  \item $\{X_t\}$ es una martingala respecto a $(\F_t)$ bajo $\PP$;
  \item el proceso $\{X_t^2 - t\}$ es también una martingala respecto a
        $(\F_t)$ bajo $\PP$;
\end{enumerate}
entonces $\{X_t\}$ es un movimiento browniano estándar bajo $\PP$.
\end{teorema}

La demostración de este teorema, tal y como aquí viene formulada, puede encontrarse en \cite{etheridge}, sección 4.3.

La condición (ii) equivale, para martingalas continuas, a $[X]_t =
t$. La caracterización dice por tanto: una martingala continua con
$X_0 = 0$ y variación cuadrática $t$ es necesariamente un browniano.

\medskip

\noindent\textbf{El teorema de Girsanov.} Pasamos al enunciado del
resultado central.

\begin{teorema}[Girsanov]
\label{teo:girsanov}
Sea $\{W_t\}_{t \in [0, T]}$ un movimiento browniano estándar sobre
$(\Omega, \F, (\F_t), \PP)$, y sea $\{\theta_t\}_{t \in [0, T]}$ un
proceso adaptado que satisface la condición de Novikov
\eqref{eq:novikov}. Definamos $\{L_t\}$ por~\eqref{eq:L_t_recordada} y
la medida $\Q$ sobre $(\Omega, \F_T)$ por
$\Q(A) = \int_A L_T\, d\PP$. Entonces:
\begin{enumerate}[label=\textnormal{(\roman*)}]
  \item $\Q$ es una medida de probabilidad equivalente a $\PP$, con
        $d\Q/d\PP = L_T$;
  \item el proceso definido por
        \begin{equation}\label{eq:def_W_tilde}
          \widetilde{W}_t \eqdef W_t + \int_0^t \theta_s\, ds,
          \qquad t \in [0, T],
        \end{equation}
        es un movimiento browniano estándar sobre $(\Omega, \F, (\F_t),
        \Q)$.
\end{enumerate}
\end{teorema}

\begin{proof}
La afirmación (i) está esencialmente demostrada en la subsección
anterior. Por la condición de Novikov y el
Teorema~\ref{teo:novikov}, $\{L_t\}$ es una $\PP$-martingala
estrictamente positiva con $L_0 = 1$, y por tanto $\E_\PP[L_T] = 1$.
La aplicación $A \mapsto \int_A L_T\, d\PP$ define entonces una medida
de probabilidad sobre $(\Omega, \F_T)$, equivalente a $\PP$ por la
positividad estricta de $L_T$, y con derivada de Radon-Nikodym $L_T$
por construcción.

Para la afirmación (ii) aplicamos la
caracterización de Lévy (Teorema~\ref{teo:caracterizacion_levy}) bajo
la medida $\Q$. Es necesario verificar que $\widetilde{W}$ satisface las
hipótesis del teorema: continuidad trayectorial, $\widetilde{W}_0
= 0$, propiedad de martingala bajo $\Q$, y la propiedad de que
$\{\widetilde{W}_t^2 - t\}$ sea $\Q$-martingala.

\emph{Paso 1: continuidad y valor inicial.} La continuidad de
$\{\widetilde{W}_t\}$ es inmediata: $W_t$ es continuo casi seguramente
(condición (iv) de la Definición~\ref{def:browniano}) y la integral
$\int_0^t \theta_s\, ds$ es continua en $t$ por construcción. La suma
es por tanto continua. Asimismo, $\widetilde{W}_0 = W_0 + 0 = 0$.

\emph{Paso 2: $\{\widetilde{W}_t\}$ es $\Q$-martingala.} La estrategia
es la siguiente: por un lema clásico que ahora detallamos, un proceso
$\{Y_t\}$ es $\Q$-martingala si y solo si $\{Y_t\, L_t\}$ es
$\PP$-martingala. Aplicaremos esta caracterización a $Y_t =
\widetilde{W}_t$.

\noindent Demostremos primero el lema. Si $\{Y_t L_t\}$ es $\PP$-martingala,
entonces para $0 \leq s \leq t$ y $A \in \F_s$,
\begin{equation*}
  \int_A Y_t\, L_t\, d\PP = \int_A Y_s\, L_s\, d\PP.
\end{equation*}
El lado izquierdo, por~\eqref{eq:def_Q_continuo} aplicada con
restricción a $\F_t$ (recordemos que $L_t$ es la derivada
de Radon-Nikodym de $\Q\!\restriction_{\F_t}$ respecto a
$\PP\!\restriction_{\F_t}$), es $\int_A Y_t\, d\Q$. Análogamente, el
lado derecho es $\int_A Y_s\, d\Q$. Por la unicidad de la esperanza
condicional, $\E_\Q[Y_t \mid \F_s] = Y_s$, esto es, $\{Y_t\}$ es
$\Q$-martingala. La implicación recíproca se obtiene invirtiendo el
argumento.

Aplicado el lema, basta demostrar que $\{\widetilde{W}_t\, L_t\}$ es
$\PP$-martingala. Calculemos su SDE bajo $\PP$ aplicando la fórmula
del producto, que es un caso particular de la fórmula de Itô general
(Teorema~\ref{teo:formula_ito_general}) aplicada a $f(x, y) = xy$ con
$X_t = \widetilde{W}_t$ e $Y_t = L_t$ como argumentos. Las
diferenciales de los procesos son
\begin{equation*}
  d\widetilde{W}_t = \theta_t\, dt + dW_t, \qquad
  dL_t = -\theta_t\, L_t\, dW_t,
\end{equation*}
donde la primera identidad se sigue de la
definición~\eqref{eq:def_W_tilde}, y la segunda fue establecida
en~\eqref{eq:sde_martingala_exp}. La fórmula del producto produce
\begin{equation*}
  d(\widetilde{W}_t\, L_t)
  = L_t\, d\widetilde{W}_t + \widetilde{W}_t\, dL_t
   + d\widetilde{W}_t\, dL_t,
\end{equation*}
donde el último término es el producto formal de las dos
diferenciales, evaluado mediante las reglas de
Itô~\eqref{eq:reglas_ito}. Calculémoslo:
\begin{equation*}
  d\widetilde{W}_t\, dL_t
  = (\theta_t\, dt + dW_t) \cdot (-\theta_t\, L_t\, dW_t)
  = -\theta_t^2\, L_t\, dt\, dW_t - \theta_t\, L_t\, (dW_t)^2.
\end{equation*}
Por las reglas de Itô, $dt\, dW_t = 0$ y $(dW_t)^2 = dt$, de modo que
$d\widetilde{W}_t\, dL_t = -\theta_t\, L_t\, dt$. Sustituyendo,
\begin{align*}
  d(\widetilde{W}_t\, L_t)
  &= L_t\,(\theta_t\, dt + dW_t)
   + \widetilde{W}_t\, (-\theta_t\, L_t\, dW_t)
   - \theta_t\, L_t\, dt \\
  &= \underbrace{(L_t\, \theta_t - \theta_t\, L_t)}_{= 0}\, dt
   + L_t\, dW_t - \widetilde{W}_t\, \theta_t\, L_t\, dW_t \\
  &= L_t\, (1 - \theta_t\, \widetilde{W}_t)\, dW_t.
\end{align*}
La cancelación del término en $dt$ es el cálculo crucial de la
demostración: muestra que $\{\widetilde{W}_t\, L_t\}$ es un proceso
sin deriva, expresable como integral estocástica del proceso $L_t (1 -
\theta_t\, \widetilde{W}_t)$ respecto a $W$. Por el
Teorema~\ref{teo:martingala_integral_indefinida}, asumiendo
condiciones suficientes de integrabilidad sobre $\theta$, garantizadas
en última instancia por la condición de Novikov, $\{\widetilde{W}_t\,
L_t\}$ es una $\PP$-martingala. Por el lema previo, $\{\widetilde{W}_t\}$
es una $\Q$-martingala.

\emph{Paso 3: $\{\widetilde{W}_t^2 - t\}$ es $\Q$-martingala.} Por la
caracterización del Paso 2, basta demostrar que $\{(\widetilde{W}_t^2
- t)\, L_t\}$ es $\PP$-martingala. Aplicando la fórmula de Itô a
$f(x) = x^2$ con $X_t = \widetilde{W}_t$ y usando $(d\widetilde{W}_t)^2
= dt$ por las reglas de Itô, se obtiene $d(\widetilde{W}_t^2 - t) = 2\,
\widetilde{W}_t\, \theta_t\, dt + 2\, \widetilde{W}_t\, dW_t$.
Multiplicando por $L_t$ y aplicando la fórmula del producto exactamente
como en el Paso 2, los términos en $dt$ se cancelan por la misma
razón de fondo (la corrección $-\tfrac{1}{2} \int \theta_s^2\, ds$
en el exponente de $L_t$ está calibrada para producir esta cancelación),
y queda
\begin{equation*}
  d\bigl((\widetilde{W}_t^2 - t)\, L_t\bigr)
  = (\text{términos en } dW_t),
\end{equation*}
sin término en $dt$. El proceso $(\widetilde{W}_t^2 - t)\, L_t$ es,
por tanto, $\PP$-martingala bajo las condiciones de integrabilidad
oportunas, y por el lema de equivalencia $\{\widetilde{W}_t^2 - t\}$
es $\Q$-martingala.

\emph{Conclusión.} El proceso $\{\widetilde{W}_t\}$ es continuo,
satisface $\widetilde{W}_0 = 0$, es $\Q$-martingala, y $\{\widetilde{W}_t^2
- t\}$ es también $\Q$-martingala. Por la caracterización de Lévy
(Teorema~\ref{teo:caracterizacion_levy}), $\widetilde{W}$ es un
movimiento browniano estándar bajo $\Q$.
\end{proof}

\medskip

\noindent\textbf{Lectura del teorema.} El teorema de Girsanov es un resultado de gran interés a efectos de nuestro trabajo: el cambio de medida codificado
por la martingala exponencial $\{L_t\}$ \emph{absorbe la deriva
adicional $\theta_t$ del proceso $W_t + \int_0^t \theta_s\, ds$ hasta
hacerla desaparecer}. Bajo la nueva medida $\Q$, lo que era un proceso
con deriva $\theta_t$ se vuelve un browniano puro. La operación es
reversible: aplicando el teorema con $-\theta$ se recupera $W$ a partir
de $\widetilde{W}$.

\medskip

\noindent\textbf{Reformulación en términos de procesos de Itô.} El uso
en la práctica del teorema en el Capítulo~\ref{cap:bs} pasará por la
siguiente reformulación. Sea $\{X_t\}$ un proceso de Itô con dinámica
\begin{equation*}
  dX_t = a_t\, dt + b_t\, dW_t,
\end{equation*}
y supongamos $b_t > 0$ casi seguramente. Definamos $\theta_t \eqdef
a_t / b_t$, y supongamos que satisface la condición de Novikov.
Aplicando el Teorema~\ref{teo:girsanov} con este $\theta$, bajo la
medida $\Q$ resultante el proceso $\widetilde{W}$ definido
en~\eqref{eq:def_W_tilde} es un browniano. Sustituyendo $dW_t =
d\widetilde{W}_t - \theta_t\, dt$ en la SDE de $X$,
\begin{equation*}
  dX_t = a_t\, dt + b_t\, (d\widetilde{W}_t - \theta_t\, dt)
       = (a_t - b_t\, \theta_t)\, dt + b_t\, d\widetilde{W}_t.
\end{equation*}
Sustituyendo $\theta_t = a_t/b_t$, el coeficiente de $dt$ se anula, y
\begin{equation}\label{eq:girsanov_operativo}
  \boxed{\; dX_t = b_t\, d\widetilde{W}_t. \;}
\end{equation}
Bajo la medida $\Q$, el proceso $X$ pierde su deriva y se convierte
en una integral estocástica pura respecto al nuevo browniano
$\widetilde{W}$. En particular, por el
Teorema~\ref{teo:martingala_integral_indefinida}, $\{X_t\}$ es una
$\Q$-martingala.

\section{Teorema de representación de martingalas brownianas}
\label{sec:representacion_martingalas}

Esta sección presenta la generalización al contexto continuo del
teorema de representación binomial
(subsección~\ref{subsec:representacion_binomial}). A diferencia del
caso discreto, donde el integrando se construía explícitamente
resolviendo un sistema lineal nodo a nodo, la versión continua se
apoya en argumentos de análisis funcional y \emph{no proporciona
una fórmula explícita}: garantiza la existencia del integrando, no
su construcción.

\medskip

Trabajamos a lo largo de la
sección sobre el espacio de probabilidad filtrado $(\Omega, \F,
(\F_t)_{t \in [0, T]}, \PP)$, donde $(\F_t)$ es la filtración del
browniano $\{W_t\}$ fijada en la Definición~\ref{def:filtracion_aumentada}. Es esencial subrayar que el teorema requiere trabajar exactamente
con esta filtración (generada por el browniano más los conjuntos
nulos): si la filtración de base contuviera información estricta
adicional, ajena al browniano, el resultado dejaría de valer. Esta restricción no será problemática en las aplicaciones del
Capítulo~\ref{cap:bs}, donde la única fuente de aleatoriedad será
precisamente el browniano que modela el subyacente.

\begin{teorema}[Representación de variables $L^2$]
\label{teo:representacion_l2}
Sea $X$ una variable aleatoria $\F_T$-medible con $\E[X^2] < \infty$.
Existe un único proceso $\phi \in \mathcal{H}^2[0, T]$ tal que
\begin{equation}\label{eq:representacion_l2}
  X = \E[X] + \int_0^T \phi_s\, dW_s
  \qquad \text{casi seguramente}.
\end{equation}
\end{teorema}

La demostración escapa al alcance del presente trabajo, pero podemos
articular con detalle la estrategia, que es la misma que ya empleamos
al construir la integral estocástica en la
Sección~\ref{subsec:extension_H2}: identificar un subespacio
suficientemente amplio de $L^2(\Omega, \F_T, \PP)$ donde la
representación es directa, y extenderla por densidad e isometría al
espacio completo.

Concretamente, se considera el subespacio
\[
  K \eqdef \biggl\{ M_0 + \int_0^T \phi_s\, dW_s :
                    M_0 \in \R,\ \phi \in \mathcal{H}^2[0, T] \biggr\}
  \subseteq L^2(\Omega, \F_T, \PP),
\]
formado precisamente por las variables aleatorias que admiten la
representación buscada. La prueba a partir de aquí consiste en dos hechos.

El primero es que $K$ es un subespacio \emph{cerrado} de $L^2(\Omega,
\F_T, \PP)$. La verificación es directa: si $(Y_n) \subset K$
converge en $L^2$ a un cierto $Y$, las constantes $\E[Y_n]$ convergen
a $\E[Y]$ por continuidad de la esperanza, y los integrandos
$\phi^{(n)}$ asociados forman una sucesión de Cauchy en
$\mathcal{H}^2[0, T]$ por la \emph{isometría de It\^o}
(Teorema~\ref{teo:propiedades_integral_H2}); su límite $\phi$
proporciona la representación de $Y$, así que $Y \in K$.

El segundo es que las \emph{exponenciales estocásticas} de la forma
\begin{equation*}
  \mathcal{E}(f)
  \eqdef \exp\!\biggl(\int_0^T f(s)\, dW_s
                       - \tfrac{1}{2} \int_0^T f(s)^2\, ds\biggr),
\end{equation*}
con $f : [0, T] \to \R$ determinista y de cuadrado integrable,
pertenecen a $K$ y, además, generan un subespacio denso de
$L^2(\Omega, \F_T, \PP)$. La pertenencia a $K$ es un cálculo directo
con la fórmula de It\^o aplicada a la SDE $d\mathcal{E}(f)_t =
\mathcal{E}(f)_t\, f(t)\, dW_t$, que en su forma integral da
$\mathcal{E}(f) = 1 + \int_0^T \mathcal{E}(f)_s\, f(s)\, dW_s$. La
densidad, en cambio, es la parte verdaderamente delicada del
argumento: se demuestra mostrando que ninguna $Z \in L^2(\Omega,
\F_T, \PP)$ puede ser ortogonal a todas las exponenciales, lo que se
reduce, vía un cambio analítico de parámetro, a un argumento de
inyectividad de la transformada de Fourier de las distribuciones
gaussianas finito-dimensionales del browniano.

Combinando ambos hechos, $K$ es un subespacio cerrado que contiene
un subconjunto denso, y por tanto coincide con todo $L^2(\Omega,
\F_T, \PP)$. La unicidad del integrando $\phi$ se sigue, una vez
más, de la isometría de It\^o. Para una exposición detallada nos
remitimos a \cite{etheridge}, sección~4.6.

\medskip

\noindent\textbf{Versión para martingalas.} El
Teorema~\ref{teo:representacion_l2} admite la siguiente formulación
general, que es con la que trabajaremos en el
Capítulo~\ref{cap:bs}.

\begin{teorema}[Representación de martingalas brownianas]
\label{teo:representacion_martingalas}
Sea $\{M_t\}_{t \in [0, T]}$ una martingala respecto a la filtración $(\F_t)$ del browniano, con $M_t \in L^2$ para todo $t \in [0,
T]$. Existe un único proceso $\phi \in \mathcal{H}^2[0, T]$ tal que
\begin{equation}\label{eq:representacion_martingalas}
  M_t = M_0 + \int_0^t \phi_s\, dW_s,
  \qquad t \in [0, T],
  \quad \text{casi seguramente}.
\end{equation}
\end{teorema}

\begin{proof}
Aplicando el Teorema~\ref{teo:representacion_l2} a la variable
$M_T$, existe un único $\phi \in \mathcal{H}^2$ con
\begin{equation*}
  M_T = M_0 + \int_0^T \phi_s\, dW_s
\end{equation*}
casi seguramente, donde hemos usado $\E[M_T] = M_0$ por la propiedad
de martingala. Para extender la identidad a $t < T$, tomamos
esperanza condicional respecto a $\F_t$ en ambos lados. El lado
izquierdo da $\E[M_T \mid \F_t] = M_t$ por la propiedad de martingala
de $M$. El lado derecho da $M_0 + \E[\int_0^T \phi_s\, dW_s \mid \F_t]
= M_0 + \int_0^t \phi_s\, dW_s$, ya que la integral indefinida hasta
$T$ es martingala (Teorema~\ref{teo:martingala_integral_indefinida}) y
por tanto su esperanza condicional en $\F_t$ es la integral hasta
$t$. La identidad~\eqref{eq:representacion_martingalas} queda
demostrada. La unicidad se sigue de la unicidad en el
Teorema~\ref{teo:representacion_l2}.
\end{proof}

\medskip

\noindent\textbf{Lectura financiera del integrando.} La interpretación
financiera del integrando $\phi$ en el
Teorema~\ref{teo:representacion_martingalas} es enteramente análoga a
la del caso discreto. Si $\{M_t\}$ es la martingala asociada al valor
descontado de un derivado, construida en el Capítulo~\ref{cap:bs} como
$M_t = \E_\Q[e^{-rT}\, C_T \mid \F_t]$ (siguiendo el mismo patrón que en \eqref{eq:def_M_replicante_discreto}), entonces el proceso $\phi_t$
que aparece en su representación es la \emph{cantidad de subyacente}
que el portfolio replicante debe mantener en el instante $t$. Es la
generalización al contexto continuo de la fórmula
explícita~\eqref{eq:phi_explicito} del modelo binomial, con la salvedad
ya señalada: en el caso continuo $\phi$ existe pero no se construye
explícitamente.

\section{Síntesis: del modelo binomial al cálculo estocástico}
\label{sec:sintesis_cap4}

Cerramos el capítulo estableciendo la correspondencia entre las
herramientas del modelo binomial y las del cálculo estocástico
continuo.

\medskip

\noindent\textbf{La correspondencia entre los dos contextos.} Cada una
de las cuatro piezas del cálculo estocástico continuo tiene su
contraparte en el caso discreto. La \emph{integral estocástica de
Itô}, $\int_0^t \phi_s\, dW_s$, es la versión continua de la integral
estocástica discreta $(\phi \cdot \widetilde S)_n$
de~\eqref{eq:integral_discreta}: ambas expresan la ganancia acumulada
de una estrategia $\phi$ predecible, y la condición de previsibilidad, formulada en discreto como $\F_{k-1}$-medibilidad y en continuo como
medibilidad respecto a $\F_{t_i}$ (sobre los procesos simples), es la
que hace que la integral preserve la propiedad de martingala
(Teoremas~\ref{teo:integral_preserva_martingala}
y~\ref{teo:martingala_integral_indefinida}). La \emph{fórmula de Itô}
es la versión continua de la regla de la cadena clásica, con la
novedad del término extra $\tfrac{1}{2}\, b_t^2\, \partial_{xx} f\,
dt$ en~\eqref{eq:formula_ito_general}, que no tiene contraparte en el
caso discreto y refleja la variación cuadrática no
nula del browniano.

El \emph{teorema de Girsanov} es la versión continua de la
construcción explícita de la medida $\Q$ en el modelo
binomial~\eqref{eq:def_Q_multi}. En ambos casos, el cambio de medida
está controlado por una densidad $L = d\Q/d\PP$. La diferencia es que
en discreto $L$ se construye combinatorialmente, asignando a cada
trayectoria una probabilidad $q^j(1-q)^{N-j}$, mientras que en
continuo $L_T$ es la martingala exponencial $\exp(-\int \theta\, dW -
\tfrac{1}{2}\int \theta^2\, ds)$. El \emph{teorema de representación
de martingalas brownianas} (Teorema~\ref{teo:representacion_martingalas})
es la versión continua del teorema de representación binomial
(Teorema~\ref{teo:representacion_binomial}). En el caso discreto, la
demostración fue constructiva y disponemos de una fórmula explícita
del integrando; en el caso continuo, en cambio, la demostración apela
a argumentos abstractos de densidad y completitud, y no proporciona
expresión explícita: saber cómo construirlo será, en cada problema
concreto, una cuestión adicional que se resolverá mediante la fórmula
de Itô.

\medskip

\noindent\textbf{La receta en tres pasos, en versión continua.} Las
cuatro piezas anteriores son los ingredientes con los que obtendremos
la solución del problema de valoración en el contexto continuo. La
estrategia, fiel al esquema del modelo binomial discreto
(subsección~\ref{subsec:tres_pasos_discreto}), seguirá los tres mismos
pasos:

\begin{enumerate}[label=\textbf{Paso \arabic*.}, leftmargin=*]
  \item \emph{Cambio de medida.} Identificar la medida $\Q$ bajo la
        cual el precio descontado del subyacente es martingala. Esta
        identificación, que en el caso discreto se hizo por
        construcción explícita, se hará en el caso continuo aplicando
        el \emph{teorema de Girsanov}.
  \item \emph{Construcción del proceso de valor descontado.} Definir
        el proceso $\{M_t\}$ asociado al \emph{payoff} descontado del
        derivado vía $M_t = \E_\Q[e^{-rT}\, C_T \mid \F_t]$. Por la
        propiedad de la torre, $\{M_t\}$ es $\Q$-martingala. Este paso
        es idéntico al del caso discreto.
  \item \emph{Aplicación del teorema de representación.} Aplicar el
        \emph{teorema de representación de martingalas brownianas} para
        obtener un proceso $\phi$ tal que $M_t = M_0 + \int_0^t \phi_s\,
        d\widetilde{W}_s$, donde $\widetilde{W}$ es el browniano bajo
        $\Q$ que produce Girsanov. El integrando $\phi$ codifica la
        cantidad de subyacente del portfolio replicante.
\end{enumerate}

La estructura es la misma que en el caso discreto: solo cambia la
sofisticación técnica de las herramientas. El
Capítulo~\ref{cap:bs} aplica esta receta al browniano geométrico
como subyacente y obtiene, como resultado central, la fórmula de
Black-Scholes para la opción \emph{call} europea.

% =====================================================================
%   CAPÍTULO 5 — EL MODELO DE BLACK-SCHOLES
% =====================================================================
\chapter{El modelo de Black-Scholes}\label{cap:bs}

Llegamos al objetivo último del trabajo: aplicar el aparato del
cálculo estocástico al \emph{modelo de Black-Scholes} y obtener,
como resultado central, la fórmula clásica de valoración de la
opción \emph{call} europea. La estructura del capítulo sigue el mismo patrón que el Capítulo~\ref{cap:discreto}: hipótesis del modelo, ejecución de
la receta de tres pasos, particularización al \emph{call}, y
recuperación de la EDP clásica.

\section{El modelo de Black-Scholes}
\label{sec:modelo_bs}

\subsection{Hipótesis del mercado y dinámica del subyacente}
\label{subsec:hipotesis_modelo_bs}

\noindent\textbf{Las hipótesis del modelo.} El modelo de Black-Scholes
describe un mercado idealizado articulado a partir de las siguientes
hipótesis estructurales:

\begin{enumerate}[label=\textnormal{(H\arabic*)}, leftmargin=*]
  \item El mercado consta de dos activos: un \emph{subyacente} arriesgado
        cuyo precio es un proceso estocástico $\{S_t\}_{t \in [0, T]}$,
        y un \emph{bono libre de riesgo} con tipo de interés constante
        $r > 0$.
  \item El subyacente sigue un movimiento browniano geométrico (GBM): su
        dinámica está gobernada por la ecuación diferencial estocástica
        \begin{equation}\label{eq:bs_dinamica}
          dS_t = \mu\, S_t\, dt + \sigma\, S_t\, dW_t,
          \qquad S_0 > 0,
        \end{equation}
        donde $\mu \in \R$ es el rendimiento esperado y $\sigma > 0$ es
        la \emph{volatilidad}, ambos constantes, y $\{W_t\}$ es un
        movimiento browniano estándar sobre un espacio de probabilidad
        filtrado $(\Omega, \F, (\F_t), \PP)$ con $(\F_t)$ la filtración
        aumentada del browniano.
  \item El mercado es \emph{libre de fricciones}: no hay costes de
        transacción, los activos son infinitamente divisibles, las
        ventas en corto están permitidas y el préstamo y el endeudamiento
        se realizan al mismo tipo $r$.
  \item El mercado es \emph{libre de arbitraje}: no existe ninguna
        estrategia que, partiendo de capital nulo, produzca con
        probabilidad estrictamente positiva un valor terminal estrictamente positivo y nunca negativo.
\end{enumerate}

Las hipótesis (H1) y (H3)–(H4) son traducción literal al contexto
continuo de las hipótesis del modelo binomial del
Capítulo~\ref{cap:discreto}: la estructura del mercado, la idealización
y la condición de no arbitraje se mantienen intactas. La novedad es (H2): la elección concreta de la dinámica del subyacente,
que en el caso discreto era el random walk binomial multiplicativo y
en el caso continuo es el GBM.

\medskip

\noindent\textbf{Forma explícita del subyacente.} La SDE
\eqref{eq:bs_dinamica} admite una solución cerrada cuya expresión ya
obtuvimos como Ejemplo~2 de la
Sección~\ref{subsec:formula_ito_general}: el GBM con SDE
\eqref{eq:bs_dinamica} es exactamente
\begin{equation}\label{eq:bs_S_explicito}
  S_t = S_0\, \exp\!\left(\Bigl(\mu - \tfrac{1}{2}\sigma^2\Bigr) t
                          + \sigma\, W_t\right),
  \qquad t \in [0, T].
\end{equation}
La identidad~\eqref{eq:bs_S_explicito} muestra que $S_t$ es estrictamente positivo casi seguramente (propiedad económicamente exigible al
precio de un activo), y que su logaritmo $\log S_t$ tiene distribución
gaussiana $\mathcal{N}\bigl(\log S_0 + (\mu - \tfrac{1}{2}\sigma^2) t,\,
\sigma^2 t\bigr)$. Esta propiedad de log-normalidad será la pieza
clave del cálculo explícito de la fórmula de Black-Scholes en la
Sección~\ref{sec:formula_bs}.

\medskip

\noindent\textbf{El bono libre de riesgo.} Con tipo de interés constante
$r$, el precio del bono libre de riesgo evoluciona deterministicamente
mediante
\begin{equation}\label{eq:bs_bono}
  B_t = e^{rt},
  \qquad t \in [0, T],
\end{equation}
que es la solución a la EDO ordinaria $dB_t = r\, B_t\, dt$ con
$B_0 = 1$, sin término
estocástico (el bono no tiene riesgo).

\subsection{Portfolio en tiempo continuo}
\label{subsec:portfolio_continuo}

\begin{definicion}[Portfolio en tiempo continuo]
\label{def:portfolio_continuo}
Un \emph{portfolio} (o \emph{estrategia}) sobre el horizonte $[0, T]$
es un par de procesos estocásticos $(\phi_t, \psi_t)_{t \in [0, T]}$
donde:
\begin{enumerate}[label=\textnormal{(\roman*)}]
  \item $\phi_t$ representa la cantidad de subyacente en el portfolio
        en el instante $t$;
  \item $\psi_t$ representa la cantidad de bono libre de riesgo en el
        portfolio en el instante $t$;
  \item ambos procesos son adaptados a la filtración $(\F_t)$ del
        browniano y satisfacen las condiciones técnicas de
        integrabilidad
        \begin{equation*}
          \phi \in \mathcal{H}^2[0, T], \qquad
          \int_0^T |\psi_s|\, ds < \infty \text{ casi seguramente}.
        \end{equation*}
\end{enumerate}
El \emph{valor} del portfolio en el instante $t$ se denota $V_t$ y
viene dado por
\begin{equation}\label{eq:valor_portfolio}
  V_t \eqdef \phi_t\, S_t + \psi_t\, B_t.
\end{equation}
\end{definicion}

Las condiciones de integrabilidad sobre $(\phi, \psi)$ son las que
garantizan que las integrales que aparecerán a continuación, la integral
estocástica $\int \phi\, dS$ y la integral ordinaria $\int \psi\, dB$, están bien definidas en el sentido del Capítulo~\ref{cap:itocalculus}.
La adaptación de $\phi$ y $\psi$ a $(\F_t)$ es la traducción al
contexto continuo de la previsibilidad discreta
(Sección~\ref{subsec:predecibles_int_disc}): el portfolio se decide con
la información disponible antes del movimiento del mercado.

\subsection{La condición \emph{self-financing}}
\label{subsec:self_financing_continuo}

La pieza estructural más delicada del modelo continuo es la traducción
al contexto continuo de la condición de \emph{auto-financiación}
introducida en el caso discreto en la
Sección~\ref{subsec:tres_pasos_discreto}. Recordemos su contenido:
una estrategia binomial era \emph{self-financing} si en cada nodo del
árbol se cumplía
\begin{equation}\label{eq:self_financing_discreto_recordatorio}
  \phi_{k+1}\, S_k + \psi_{k+1}\, B_k = \phi_k\, S_k + \psi_k\, B_k,
\end{equation}
es decir, si la cantidad de capital disponible al inicio del período
$k+1$, tras reasignar el portfolio, coincide con la que estaba
disponible al final del período $k$. Reordenando los términos,
\eqref{eq:self_financing_discreto_recordatorio} equivale a
\begin{equation*}
  S_k\, (\phi_{k+1} - \phi_k) + B_k\, (\psi_{k+1} - \psi_k) = 0,
\end{equation*}
y consecuentemente, sumando sobre los períodos hasta el instante $n$,
\begin{equation*}
  V_n - V_0
  = \sum_{k=0}^{n-1} \bigl[\phi_{k+1}\, (S_{k+1} - S_k)
                            + \psi_{k+1}\, (B_{k+1} - B_k)\bigr].
\end{equation*}
El valor del portfolio cambia exclusivamente debido a los movimientos
del mercado, no a aportes ni retiros de capital.

\medskip

\begin{definicion}[Estrategia \emph{self-financing}]
\label{def:self_financing_continuo}
Una estrategia $(\phi_t, \psi_t)_{t \in [0, T]}$ se dice
\emph{self-financing} si su valor $V_t = \phi_t\, S_t + \psi_t\, B_t$
satisface
\begin{equation}\label{eq:self_financing_continuo}
  dV_t = \phi_t\, dS_t + \psi_t\, dB_t,
  \qquad t \in [0, T],
\end{equation}
o, en forma integral,
\begin{equation}\label{eq:self_financing_integral}
  V_t = V_0 + \int_0^t \phi_s\, dS_s + \int_0^t \psi_s\, dB_s,
  \qquad t \in [0, T].
\end{equation}
\end{definicion}

La identidad~\eqref{eq:self_financing_continuo} es la traducción
exacta al contexto continuo de la condición discreta:
articula que el cambio del valor del portfolio en el instante $t$
proviene exclusivamente de los cambios en los precios del subyacente
y del bono, ponderados respectivamente por las cantidades $\phi_t$ y
$\psi_t$ que se mantienen en cada instante.

\medskip

\noindent\textbf{Equivalencia con la formulación primitiva.} La
Definición~\ref{def:self_financing_continuo} es la formulación
operativa de la condición self-financing. Conviene verificar su
equivalencia con la formulación original, que exige coste nulo en la reasignación
infinitesimal del portfolio.

\begin{proposicion}[Equivalencia de las dos formulaciones]
\label{prop:equivalencia_self_financing}
Sea $(\phi_t, \psi_t)$ una estrategia con $V_t = \phi_t\, S_t +
\psi_t\, B_t$. Las siguientes condiciones son equivalentes:
\begin{enumerate}[label=\textnormal{(\roman*)}]
  \item $dV_t = \phi_t\, dS_t + \psi_t\, dB_t$;
  \item $S_t\, d\phi_t + B_t\, d\psi_t + d\langle \phi, S\rangle_t = 0$,
\end{enumerate}
donde $\langle \phi, S\rangle_t$ denota la covariación cuadrática entre
los procesos $\phi$ y $S$.
\end{proposicion}

\begin{proof}
Aplicamos la fórmula del producto (caso particular de la fórmula de
Itô, Teorema~\ref{teo:formula_ito_general}) al producto $\phi_t\, S_t$:
\begin{equation*}
  d(\phi_t\, S_t) = \phi_t\, dS_t + S_t\, d\phi_t + d\langle \phi, S\rangle_t.
\end{equation*}
Análogamente, para $\psi_t\, B_t$, donde $B_t$ tiene variación acotada,
luego $\langle \psi, B\rangle_t = 0$,
\begin{equation*}
  d(\psi_t\, B_t) = \psi_t\, dB_t + B_t\, d\psi_t.
\end{equation*}
Sumando ambas identidades, obtenemos la diferencial del valor:
\begin{equation*}
  dV_t = d(\phi_t S_t + \psi_t B_t)
       = \phi_t\, dS_t + \psi_t\, dB_t
       + \bigl[S_t\, d\phi_t + B_t\, d\psi_t + d\langle \phi, S\rangle_t\bigr].
\end{equation*}
La condición (i) equivale, por consiguiente, a la anulación del
término entre corchetes, que es exactamente la condición (ii).
\end{proof}

\section{Aplicación de la receta de tres pasos}
\label{sec:tres_pasos_continuo}

Ejecutamos en esta sección la receta de tres pasos articulada en la
subsección~\ref{subsec:tres_pasos_discreto} para el caso discreto,
ahora en el contexto continuo.

\subsection{Planteamiento del problema}
\label{subsec:planteamiento_continuo}

\noindent\textbf{Derivados europeos.}

\begin{definicion}[Derivado europeo]
\label{def:derivado_europeo}
Un \emph{derivado europeo} con vencimiento $T$ sobre el subyacente
$\{S_t\}$ es un contrato que paga al tenedor, en el instante $t = T$,
una cantidad $C_T = h(S_T)$, donde $h : (0, \infty) \to \R$ es una
función medible, el \emph{payoff} del derivado, que satisface la
condición de integrabilidad
\begin{equation}\label{eq:integrabilidad_payoff}
  \E_\PP\bigl[h(S_T)^2\bigr] < \infty.
\end{equation}
\end{definicion}

La condición~\eqref{eq:integrabilidad_payoff} es la mínima necesaria
para que los procesos definidos a partir del \emph{payoff} estén en
los espacios $L^2$ requeridos por el teorema de representación de
martingalas brownianas. Para los \emph{payoffs} estándar, incluida la
\emph{call} europea $h(S_T) = (S_T - K)^+$ que abordaremos en la
Sección~\ref{sec:formula_bs}, la condición se verifica
automáticamente por la log-normalidad de $S_T$.

\medskip

\noindent\textbf{El problema de valoración.} El problema central es
determinar, en cada instante $t \in [0, T]$, un precio justo $V(t,
S_t)$ para el derivado, esto es, un precio compatible con la hipótesis
de no arbitraje (H4). El procedimiento a seguir, al igual que en el esquema discreto, pasa
por la noción de \emph{estrategia replicante}:

\begin{definicion}[Estrategia replicante]
\label{def:estrategia_replicante}
Una \emph{estrategia replicante} para el derivado europeo $h(S_T)$ es
una estrategia \emph{self-financing} $(\phi_t, \psi_t)_{t \in [0, T]}$
cuyo valor terminal coincide con el \emph{payoff}, esto es,
\begin{equation*}
  V_T = \phi_T\, S_T + \psi_T\, B_T = h(S_T)
  \qquad \text{casi seguramente}.
\end{equation*}
\end{definicion}

Por el argumento de no arbitraje del caso discreto, si la estrategia
replicante existe, su valor $V_t$ en cada instante es necesariamente
el precio del derivado en $t$. El problema se reduce, por tanto, a
construir una estrategia replicante.

\subsection{Paso 1: cambio de medida vía Girsanov}
\label{subsec:paso1_girsanov}

Sustituimos la medida del mercado $\PP$ por una medida equivalente
$\Q$ bajo la cual el precio descontado del subyacente sea martingala.
La identificación se realiza vía el teorema de Girsanov.\footnote{Reaparece aquí, en su forma continua, la idea anticipada en la
Observación~\ref{obs:cambio_medida_general}: pasar a una medida
equivalente (la neutral al riesgo) para simplificar el problema. Lo
que en el modelo binomial fue una reponderación explícita de las
probabilidades sobre las trayectorias del árbol es ahora el cambio de
medida que establece el teorema de Girsanov.}

\medskip

\noindent\textbf{La dinámica del precio descontado.} Para articular el
cambio de medida, conviene primero calcular la SDE del proceso de
precio descontado
\begin{equation}\label{eq:precio_descontado_def}
  \widetilde S_t \eqdef e^{-rt}\, S_t,
  \qquad t \in [0, T].
\end{equation}
Aplicamos la fórmula de Itô (Teorema~\ref{teo:formula_ito_general}) a
la función $f(t, x) = e^{-rt}\, x$ y al proceso $X_t = S_t$. Las
derivadas relevantes son
\begin{equation*}
  \partial_t f = -r\, e^{-rt}\, x, \qquad
  \partial_x f = e^{-rt}, \qquad
  \partial_{xx} f = 0.
\end{equation*}
Con $a_t = \mu\, S_t$, $b_t = \sigma\, S_t$ los coeficientes de la SDE
\eqref{eq:bs_dinamica}, la fórmula de Itô da
\begin{align*}
  d\widetilde S_t
  &= \bigl[-r\, e^{-rt}\, S_t + e^{-rt}\, \mu\, S_t + 0\bigr]\, dt
   + e^{-rt}\, \sigma\, S_t\, dW_t \\
  &= (\mu - r)\, e^{-rt}\, S_t\, dt + \sigma\, e^{-rt}\, S_t\, dW_t.
\end{align*}
Como $e^{-rt} S_t = \widetilde S_t$, obtenemos
\begin{equation}\label{eq:sde_S_descontado}
  \boxed{\; d\widetilde S_t
        = (\mu - r)\, \widetilde S_t\, dt
        + \sigma\, \widetilde S_t\, dW_t. \;}
\end{equation}
La SDE~\eqref{eq:sde_S_descontado} muestra que el precio descontado es
un nuevo proceso de Itô con coeficiente de deriva $(\mu - r)\,
\widetilde S_t$ (proporcional al exceso de rendimiento del subyacente
sobre el tipo libre de riesgo) y, por tanto, $\{\widetilde S_t\}$ no es
martingala bajo $\PP$. El objetivo del cambio de medida es eliminar esa deriva.

\medskip

\noindent\textbf{La elección de $\theta$.} La SDE
\eqref{eq:sde_S_descontado} encaja exactamente en el esquema del teorema de Girsanov reformulado en términos de procesos de Itô,
ec.~\eqref{eq:girsanov_operativo}: tenemos un proceso con dinámica
$dX_t = a_t\, dt + b_t\, dW_t$ y queremos eliminar la deriva $a_t$
mediante la elección
\begin{equation*}
  \theta_t \eqdef \frac{a_t}{b_t}
              = \frac{(\mu - r)\, \widetilde S_t}
                     {\sigma\, \widetilde S_t}
              = \frac{\mu - r}{\sigma}.
\end{equation*}
La cantidad
\begin{equation}\label{eq:precio_riesgo}
  \theta \eqdef \frac{\mu - r}{\sigma},
\end{equation}
constante en el modelo de Black-Scholes, recibe el nombre clásico de
\emph{precio de mercado del riesgo}: es exactamente la cantidad que el
teorema de Girsanov requiere para anular la deriva del precio
descontado.

\medskip

\noindent\textbf{Verificación de la condición de Novikov.} Como
$\theta$ es constante, la condición de
Novikov~\eqref{eq:novikov} se verifica trivialmente:
\begin{equation*}
  \E_\PP\!\left[\exp\!\biggl(\tfrac{1}{2}\int_0^T \theta^2\, ds\biggr)\right]
  = \exp\!\biggl(\tfrac{1}{2}\theta^2\, T\biggr) < \infty.
\end{equation*}
Estamos, por tanto, en condiciones de aplicar el
Teorema~\ref{teo:girsanov} con $\theta_t \equiv \theta$.

\medskip

\noindent\textbf{Aplicación del teorema de Girsanov.} El proceso de
densidad asociado a $\theta$ es, conforme a~\eqref{eq:L_t_recordada},
\begin{equation}\label{eq:L_T_BS}
  L_T = \exp\!\biggl(-\theta\, W_T - \tfrac{1}{2}\,\theta^2\, T\biggr).
\end{equation}
Definimos la medida $\Q$ sobre $(\Omega, \F_T)$ mediante
\begin{equation*}
  \Q(A) \eqdef \int_A L_T\, d\PP, \qquad A \in \F_T.
\end{equation*}
Por el Teorema~\ref{teo:girsanov}, $\Q$ es una medida de probabilidad
equivalente a $\PP$, con $d\Q/d\PP = L_T$, y el proceso
\begin{equation}\label{eq:W_tilde_BS}
  \widetilde W_t \eqdef W_t + \theta\, t,
  \qquad t \in [0, T],
\end{equation}
es un movimiento browniano estándar bajo $\Q$.

\medskip

\noindent\textbf{Caracterización de $\Q$ como medida martingala.}
Sustituyendo $dW_t = d\widetilde W_t - \theta\, dt$ en la SDE
\eqref{eq:sde_S_descontado} del precio descontado:
\begin{equation*}
  d\widetilde S_t = (\mu - r)\, \widetilde S_t\, dt
                  + \sigma\, \widetilde S_t\, (d\widetilde W_t - \theta\, dt).
\end{equation*}
Recordando $\theta = (\mu - r)/\sigma$, los términos en $dt$ cancelan:
\begin{equation}\label{eq:S_descontado_Q_martingala}
  \boxed{\; d\widetilde S_t = \sigma\, \widetilde S_t\, d\widetilde W_t. \;}
\end{equation}
La identidad~\eqref{eq:S_descontado_Q_martingala} muestra que, bajo
$\Q$, el precio descontado del subyacente es solución de una SDE sin
deriva: es una integral estocástica pura respecto al browniano
$\widetilde W$. Así, por el
Teorema~\ref{teo:martingala_integral_indefinida}, $\{\widetilde S_t\}$
es una $\Q$-martingala. La medida $\Q$ recibe, por consiguiente, el
nombre de \emph{medida martingala equivalente} o \emph{medida neutral
al riesgo}.

\medskip

\noindent\textbf{La dinámica del subyacente bajo $\Q$.} Conviene
también explicitar la dinámica del propio subyacente $S_t$ (no
descontado) bajo la nueva medida. Sustituyendo $dW_t = d\widetilde W_t
- \theta\, dt$ en~\eqref{eq:bs_dinamica},
\begin{equation}\label{eq:S_dinamica_Q}
  \boxed{\; dS_t = r\, S_t\, dt + \sigma\, S_t\, d\widetilde W_t. \;}
\end{equation}
Bajo $\Q$, el subyacente sigue siendo un GBM, pero con \emph{tendencia
$r$ en lugar de $\mu$}. Esta es la traducción al contexto continuo
del fenómeno que ya pudimos ver en el modelo binomial: la operación de
cambio de medida modifica la dinámica del subyacente de modo que su rendimiento esperado coincida con el del bono libre de riesgo, de ahí el nombre \emph{neutral al riesgo}.

\noindent La forma explícita del subyacente bajo $\Q$, análoga
a~\eqref{eq:bs_S_explicito}, se obtiene resolviendo la
SDE~\eqref{eq:S_dinamica_Q} mediante la misma técnica del Ejemplo~2 de
la Sección~\ref{subsec:formula_ito_general}:
\begin{equation}\label{eq:S_explicito_Q}
  S_t = S_0\, \exp\!\left(\Bigl(r - \tfrac{1}{2}\sigma^2\Bigr) t
                          + \sigma\, \widetilde W_t\right).
\end{equation}
La identidad~\eqref{eq:S_explicito_Q} será la herramienta clave para los
cálculos explícitos de la Sección~\ref{sec:formula_bs}. Bajo $\Q$, el
$\log S_T$ tiene distribución gaussiana
$\mathcal{N}\bigl(\log S_0 + (r - \tfrac{1}{2}\sigma^2) T,\, \sigma^2
T\bigr)$.

\subsection{Paso 2: la $\Q$-martingala asociada al \emph{payoff} descontado}
\label{subsec:paso2_martingala}

Identificada la medida neutral al riesgo $\Q$, el segundo paso de la receta, en analogía con el caso discreto
(subsección~\ref{subsec:tres_pasos_discreto}), consiste en construir
un proceso $\Q$-martingala asociado al \emph{payoff}.

\medskip

\noindent\textbf{Definición del proceso.} Sea $h(S_T)$ el \emph{payoff}
del derivado. Definimos el proceso $\{M_t\}_{t \in [0, T]}$ mediante
\begin{equation}\label{eq:def_M_continuo}
  M_t \eqdef \E_\Q\!\bigl[e^{-rT}\, h(S_T) \,\big|\, \F_t\bigr],
  \qquad t \in [0, T].
\end{equation}
La interpretación financiera de $M_t$ es clara: representa la mejor
estimación, condicionada a la información disponible en el instante
$t$, del valor descontado del \emph{payoff} terminal, calculada bajo
la medida neutral al riesgo $\Q$.

\medskip

\noindent\textbf{Buena definición e integrabilidad cuadrada.} Para los
\emph{payoffs} de los derivados estándar, incluida la \emph{call}
europea $(S_T - K)^+$, la condición $h(S_T) \in L^2(\Q)$ se verifica
gracias a la log-normalidad del subyacente bajo $\Q$
(ec.~\ref{eq:S_explicito_Q}) y a la equivalencia entre $\PP$ y $\Q$
(véase \cite{karatzas-shreve}, sección 5.8, para los detalles).
En consecuencia, $M_t \in L^2(\Q)$ para todo $t$, condición necesaria
para aplicar el teorema de representación de martingalas brownianas
en el paso siguiente.

\medskip

\noindent\textbf{Propiedad de martingala.}

\begin{proposicion}[$\{M_t\}$ es $\Q$-martingala]
\label{prop:M_es_Q_martingala}
El proceso $\{M_t\}$ definido en~\eqref{eq:def_M_continuo} es una
$\Q$-martingala respecto a la filtración $(\F_t)$.
\end{proposicion}

\begin{proof}
La adaptación de $M_t$ a $\F_t$ es inmediata por la propia definición
de la esperanza condicional. La integrabilidad cuadrada se sigue de
$h(S_T) \in L^2(\Q)$ y de la desigualdad de Jensen aplicada a la
esperanza condicional: $\E_\Q[M_t^2] \leq \E_\Q\bigl[
\E_\Q[(e^{-rT} h(S_T))^2 \mid \F_t]\bigr] = \E_\Q[(e^{-rT} h(S_T))^2]
< \infty$.

Para la propiedad de martingala, sean $0 \leq s \leq t \leq T$.
Aplicando la propiedad de la torre
(Proposición~\ref{prop:propiedades_esperanza_cond}, propiedad (iv))
con $\F_s \subseteq \F_t$,
\begin{equation*}
  \E_\Q[M_t \mid \F_s]
  = \E_\Q\!\left[\,\E_\Q\bigl[e^{-rT}\, h(S_T) \,\big|\, \F_t\bigr]
                 \,\Big|\, \F_s\right]
  = \E_\Q\!\bigl[e^{-rT}\, h(S_T) \,\big|\, \F_s\bigr]
  = M_s.
\end{equation*}
\end{proof}

\subsection{Paso 3: representación, replicación y \emph{self-financing}}
\label{subsec:paso3_replicacion}

\noindent\textbf{Aplicación del teorema de representación.} La
$\Q$-martingala $\{M_t\}$ está adaptada a la filtración del browniano
$\widetilde W$ ya que, por ~\eqref{eq:W_tilde_BS}, la filtración $(\F_t)$ generada por $W$ y aumentada con nulos es la misma que la generada por $\widetilde W$; además, satisface $M_t \in L^2(\Q)$ para todo $t$ por la
Proposición~\ref{prop:M_es_Q_martingala}. Estamos, por consiguiente,
en condiciones de aplicar el
Teorema~\ref{teo:representacion_martingalas} bajo la medida $\Q$ y
respecto al browniano $\widetilde W$: existe un único proceso
$\widetilde\phi \in \mathcal{H}^2[0, T]$ (donde $\mathcal{H}^2$ se
entiende ahora respecto a $\Q$) tal que
\begin{equation}\label{eq:representacion_M}
  M_t = M_0 + \int_0^t \widetilde\phi_s\, d\widetilde W_s,
  \qquad t \in [0, T].
\end{equation}

\medskip

\noindent\textbf{Identificación del portfolio.} El integrando
$\widetilde\phi_t$ proporcionado por~\eqref{eq:representacion_M} es la
pieza que permite construir la estrategia replicante. La idea es la
siguiente: queremos definir $(\phi_t, \psi_t)$ de modo que el valor
descontado del portfolio, $\widetilde V_t \eqdef e^{-rt}\, V_t$, coincida
con $M_t$ en cada instante. Para ello, observemos que si tal identidad
se cumple, entonces por~\eqref{eq:representacion_M} la diferencial de
$\widetilde V_t$ sería $d\widetilde V_t = \widetilde\phi_t\,
d\widetilde W_t$. Por otro lado, si la estrategia es \emph{self-financing}
y $\widetilde V_t = e^{-rt} V_t$, un cálculo directo (que precisaremos
después) da
\begin{equation*}
  d\widetilde V_t = \phi_t\, d\widetilde S_t
                  = \phi_t\, \sigma\, \widetilde S_t\, d\widetilde W_t,
\end{equation*}
donde la última igualdad usa~\eqref{eq:S_descontado_Q_martingala}.
Comparando ambas expresiones para $d\widetilde V_t$, identificamos
$\widetilde\phi_t = \phi_t\, \sigma\, \widetilde S_t$, esto es,
\begin{equation}\label{eq:phi_BS}
  \boxed{\;
    \phi_t \eqdef \frac{\widetilde\phi_t}{\sigma\, \widetilde S_t}
              = \frac{\widetilde\phi_t\, e^{rt}}{\sigma\, S_t}. \;}
\end{equation}
La cantidad de bono $\psi_t$ se define entonces forzando la igualdad
del valor del portfolio con el valor descontado dado por la
representación,
\begin{equation}\label{eq:psi_BS}
  \boxed{\;
    \psi_t \eqdef \frac{V_t - \phi_t\, S_t}{B_t},
    \qquad \text{con } V_t \eqdef e^{rt}\, M_t. \;}
\end{equation}
Así, $V_t = \phi_t\, S_t + \psi_t\, B_t$ por construcción y, en
particular, $V_T = e^{rT}\, M_T = e^{rT}\, e^{-rT}\, h(S_T) = h(S_T)$:
la estrategia verificaría, por tanto, la condición de replicación
(Definición~\ref{def:estrategia_replicante}). En consecuencia, basta verificar que $(\phi_t, \psi_t)$ es
\emph{self-financing} en el sentido de la
Definición~\ref{def:self_financing_continuo}.

\medskip

\noindent\textbf{Verificación de la condición \emph{self-financing}.}

\begin{teorema}[La estrategia es \emph{self-financing}]
\label{teo:replicante_self_financing}
La estrategia $(\phi_t, \psi_t)$ definida
por~\eqref{eq:phi_BS} y~\eqref{eq:psi_BS} es \emph{self-financing}, esto
es, satisface
\begin{equation*}
  dV_t = \phi_t\, dS_t + \psi_t\, dB_t.
\end{equation*}
\end{teorema}

\begin{proof}
La estrategia consiste en calcular ambos lados de la identidad y ver que coinciden.

\emph{Paso 1: cálculo de $dV_t$.} Por construcción, $V_t = e^{rt}\,
M_t$. Aplicamos la fórmula de Itô (Teorema~\ref{teo:formula_ito_general})
a la función $g(t, m) = e^{rt}\, m$ y al proceso $X_t = M_t$. Las
derivadas son $\partial_t g = r\, e^{rt}\, m$, $\partial_m g = e^{rt}$
y $\partial_{mm} g = 0$. Como $\{M_t\}$ tiene
diferencial~\eqref{eq:representacion_M}, con deriva nula y coeficiente de
difusión $\widetilde\phi_t$ respecto a $\widetilde W$, la fórmula da
\begin{equation}\label{eq:dV_calculado}
  dV_t = r\, e^{rt}\, M_t\, dt + e^{rt}\, \widetilde\phi_t\, d\widetilde W_t
       = r\, V_t\, dt + e^{rt}\, \widetilde\phi_t\, d\widetilde W_t.
\end{equation}

\emph{Paso 2: cálculo de $\phi_t\, dS_t + \psi_t\, dB_t$.}
Para $\phi_t\, dS_t$, sustituimos~\eqref{eq:S_dinamica_Q} en la
definición~\eqref{eq:phi_BS}:
\begin{align*}
  \phi_t\, dS_t
  &= \frac{\widetilde\phi_t\, e^{rt}}{\sigma\, S_t}\,
     \bigl(r\, S_t\, dt + \sigma\, S_t\, d\widetilde W_t\bigr) \\
  &= \frac{r\, \widetilde\phi_t\, e^{rt}}{\sigma}\, dt
   + \widetilde\phi_t\, e^{rt}\, d\widetilde W_t.
\end{align*}
Para $\psi_t\, dB_t$, usando $dB_t = r\, B_t\, dt$ y la
definición~\eqref{eq:psi_BS} de $\psi_t$:
\begin{equation*}
  \psi_t\, dB_t = r\, \psi_t\, B_t\, dt
              = r\, (V_t - \phi_t\, S_t)\, dt.
\end{equation*}
Sumando ambas contribuciones,
\begin{equation*}
  \phi_t\, dS_t + \psi_t\, dB_t
  = \frac{r\, \widetilde\phi_t\, e^{rt}}{\sigma}\, dt
   + \widetilde\phi_t\, e^{rt}\, d\widetilde W_t
   + r\, V_t\, dt - r\, \phi_t\, S_t\, dt.
\end{equation*}
Sustituyendo $\phi_t\, S_t = \widetilde\phi_t\, e^{rt} / \sigma$
(tras despejar de~\eqref{eq:phi_BS}), los términos
$\frac{r\, \widetilde\phi_t\, e^{rt}}{\sigma}\, dt$ y $-r\, \phi_t\,
S_t\, dt$ se cancelan y queda
\begin{equation}\label{eq:dV_self_financing}
  \phi_t\, dS_t + \psi_t\, dB_t
  = r\, V_t\, dt + \widetilde\phi_t\, e^{rt}\, d\widetilde W_t.
\end{equation}

\emph{Paso 3: comparación.} Las
expresiones~\eqref{eq:dV_calculado} y~\eqref{eq:dV_self_financing}
coinciden término a término. La estrategia $(\phi_t, \psi_t)$ es,
por consiguiente, \emph{self-financing}.
\end{proof}

\medskip

\noindent\textbf{El resultado central.} Combinando los tres pasos,
estamos en condiciones de enunciar el resultado principal de la sección.

\begin{teorema}[Valoración de derivados europeos en el modelo de Black-Scholes]
\label{teo:valoracion_BS}
Sea $h(S_T)$ un \emph{payoff} europeo que satisface la condición de
integrabilidad~\eqref{eq:integrabilidad_payoff}. Entonces existe una
única estrategia replicante \emph{self-financing} $(\phi_t, \psi_t)$
para el derivado, con valor del portfolio en cada instante $t \in [0, T]$
dado por
\begin{equation}\label{eq:formula_valoracion}
  \boxed{\;
    V(t, S_t) = \E_\Q\!\bigl[e^{-r(T - t)}\, h(S_T) \,\big|\, \F_t\bigr]. \;}
\end{equation}
Por el argumento de no arbitraje, $V(t, S_t)$ es el precio justo del
derivado en el instante $t$.
\end{teorema}

\begin{proof}
La existencia se sigue de la construcción articulada en los Pasos 1, 2
y 3, junto con la verificación
\emph{self-financing} del Teorema~\ref{teo:replicante_self_financing}.
La fórmula~\eqref{eq:formula_valoracion} se obtiene de
$V_t = e^{rt}\, M_t = e^{rt}\, \E_\Q[e^{-rT}\, h(S_T) \mid \F_t]
= \E_\Q[e^{-r(T-t)}\, h(S_T) \mid \F_t]$, usando la
definición~\eqref{eq:def_M_continuo} de $M_t$ y la $\F_t$-medibilidad
del factor $e^{rt}$.

La unicidad se sigue de la unicidad del integrando $\widetilde\phi$ en
el teorema de representación
(Teorema~\ref{teo:representacion_martingalas}): si dos estrategias
replicantes \emph{self-financing} produjeran el mismo valor terminal
$h(S_T)$, sus respectivas martingalas $\E_\Q[e^{-rT}\, h(S_T) \mid \F_t]$
coincidirían, su representación sería única, y por
\eqref{eq:phi_BS}–\eqref{eq:psi_BS} las dos estrategias serían
necesariamente iguales.
\end{proof}

\medskip

\noindent\textbf{Lectura del resultado.} La
fórmula~\eqref{eq:formula_valoracion} es la traducción al continuo
del Teorema~\ref{teo:valoracion_binomial}: el precio en cada
instante es la esperanza condicional, bajo la medida neutral al
riesgo, del \emph{payoff} descontado.

\medskip

Hemos resuelto así el problema general de valoración. La fórmula
\eqref{eq:formula_valoracion} expresa, para cualquier \emph{payoff}
europeo $h$, el precio del derivado como una esperanza condicional
bajo $\Q$. La siguiente sección particulariza esta fórmula al caso de
la opción \emph{call} europea, $h(S_T) = (S_T - K)^+$, y obtiene
mediante un cálculo directo la fórmula clásica de Black-Scholes.

\section{La fórmula de Black-Scholes para la opción \emph{call} europea}
\label{sec:formula_bs}

Particularizamos ahora la fórmula del
Teorema~\ref{teo:valoracion_BS} al \emph{payoff} de la \emph{call}
europea, $h(S_T) = (S_T - K)^+$, obteniendo con ello la fórmula clásica de
Black-Scholes (1973) con sus coeficientes $d_1$ y $d_2$.

\medskip

\noindent\textbf{La opción \emph{call} europea.} Recordemos brevemente
el contrato. La opción \emph{call} europea con \emph{strike} $K > 0$
y vencimiento $T$ es el derivado europeo cuyo \emph{payoff} en $T$ es
\begin{equation}\label{eq:payoff_call}
  h(S_T) = (S_T - K)^+ \eqdef \max(S_T - K, 0).
\end{equation}
Su tenedor adquiere el derecho (no la obligación) de comprar el
subyacente al precio $K$ en $T$: si en ese instante $S_T > K$ ejerce
el derecho y obtiene un beneficio $S_T - K$; si $S_T \leq K$ no lo
ejerce y el \emph{payoff} es nulo. La condición de
integrabilidad~\eqref{eq:integrabilidad_payoff} se verifica
automáticamente para este \emph{payoff} pues, $0 \leq (S_T - K)^+ \leq
S_T$ y los momentos del GBM son finitos a todo orden, de modo que el
Teorema~\ref{teo:valoracion_BS} es aplicable.

\medskip

\noindent\textbf{Estrategia del cálculo.} Calcularemos primero el precio del
\emph{call} en $t = 0$, esto es, $V(0, S_0) = \E_\Q[e^{-rT}\, (S_T -
K)^+]$, y obtendremos al final la generalización al caso $t > 0$ mediante un argumento sobre la propiedad markoviana de $\{S_t\}$. Para $t = 0$, la estrategia consiste en descomponer la
esperanza a partir de la función indicador de ejercicio:
\begin{equation}\label{eq:descomposicion_call}
  \E_\Q\bigl[(S_T - K)^+\bigr]
  = \E_\Q\bigl[(S_T - K)\, \mathbf{1}_{\{S_T > K\}}\bigr]
  = \underbrace{\E_\Q\bigl[S_T\, \mathbf{1}_{\{S_T > K\}}\bigr]}_{(\mathrm{I})}
   - K\, \underbrace{\Q(S_T > K)}_{(\mathrm{II})}.
\end{equation}

\noindent\textbf{Expresión de $S_T$ bajo $\Q$.}
Por~\eqref{eq:S_explicito_Q} con $t = T$,
\begin{equation}\label{eq:S_T_Q}
  S_T = S_0\, \exp\!\left(\Bigl(r - \tfrac{1}{2}\sigma^2\Bigr) T
                          + \sigma\, \widetilde W_T\right),
\end{equation}
donde $\widetilde W_T \sim \mathcal{N}(0, T)$ bajo $\Q$. Reescribiendo
$\widetilde W_T = \sqrt{T}\, Z$ con $Z \sim \mathcal{N}(0, 1)$ bajo
$\Q$,
\begin{equation*}
  S_T = S_0\, \exp\!\left(\Bigl(r - \tfrac{1}{2}\sigma^2\Bigr) T
                          + \sigma\, \sqrt{T}\, Z\right).
\end{equation*}
La desigualdad $S_T > K$ equivale, tomando logaritmos, a
\begin{equation*}
  \log\!\left(\frac{S_0}{K}\right)
  + \Bigl(r - \tfrac{1}{2}\sigma^2\Bigr) T
  + \sigma\sqrt{T}\, Z > 0,
\end{equation*}
es decir, $Z > -d_2$, donde
\begin{equation}\label{eq:def_d2}
  d_2 \eqdef \frac{\log(S_0/K) + \bigl(r - \tfrac{1}{2}\sigma^2\bigr) T}
                  {\sigma\, \sqrt{T}}.
\end{equation}

\noindent\textbf{Cálculo del término $(\mathrm{II})$.} Por la simetría
de la distribución gaussiana, $(\mathrm{II}) = \Q(Z > -d_2) = N(d_2)$,
donde $N(x) \eqdef (2\pi)^{-1/2} \int_{-\infty}^x e^{-y^2/2}\, dy$ es
la función de distribución de la normal estándar.

\medskip

\noindent\textbf{Cálculo del término $(\mathrm{I})$.} Tenemos
\begin{equation*}
  (\mathrm{I}) = \E_\Q\bigl[S_T\, \mathbf{1}_{\{Z > -d_2\}}\bigr]
              = S_0\, e^{(r - \frac{1}{2}\sigma^2) T}\,
                \E_\Q\!\left[e^{\sigma\sqrt{T}\, Z}\,
                              \mathbf{1}_{\{Z > -d_2\}}\right].
\end{equation*}
Para evaluar la esperanza, escribimos explícitamente la integral
gaussiana, usando que $Z \sim \mathcal{N}(0, 1)$, y completamos el
cuadrado en el exponente:
\begin{align*}
  \E_\Q\!\left[e^{\sigma\sqrt{T}\, Z}\, \mathbf{1}_{\{Z > -d_2\}}\right]
  &= \frac{1}{\sqrt{2\pi}} \int_{-d_2}^{\infty}
       e^{\sigma\sqrt{T}\, z}\, e^{-z^2/2}\, dz \\
  &= \frac{1}{\sqrt{2\pi}} \int_{-d_2}^{\infty}
       \exp\!\left(-\tfrac{1}{2}\bigl(z - \sigma\sqrt{T}\bigr)^2
                   + \tfrac{1}{2}\sigma^2 T\right) dz \\
  &= e^{\frac{1}{2}\sigma^2 T} \cdot
     \frac{1}{\sqrt{2\pi}} \int_{-d_2}^{\infty}
       \exp\!\left(-\tfrac{1}{2}\bigl(z - \sigma\sqrt{T}\bigr)^2\right) dz.
\end{align*}
Hacemos el cambio de variable $u \eqdef z - \sigma\sqrt{T}$, con $du =
dz$ y nuevos límites $u > -d_2 - \sigma\sqrt{T}$:
\begin{equation*}
  \E_\Q\!\left[e^{\sigma\sqrt{T}\, Z}\, \mathbf{1}_{\{Z > -d_2\}}\right]
  = e^{\frac{1}{2}\sigma^2 T}\,
    \frac{1}{\sqrt{2\pi}} \int_{-d_2 - \sigma\sqrt{T}}^{\infty}
       e^{-u^2/2}\, du
  = e^{\frac{1}{2}\sigma^2 T}\, N(d_2 + \sigma\sqrt{T}),
\end{equation*}
donde la última igualdad usa de nuevo la simetría de la gaussiana.
Definiendo
\begin{equation}\label{eq:def_d1}
  d_1 \eqdef d_2 + \sigma\sqrt{T}
       = \frac{\log(S_0/K) + \bigl(r + \tfrac{1}{2}\sigma^2\bigr) T}
              {\sigma\, \sqrt{T}},
\end{equation}
y sustituyendo en la expresión de $(\mathrm{I})$, el factor
$e^{\frac{1}{2}\sigma^2 T}$ cancela exactamente con el correspondiente término en la
deriva $r - \tfrac{1}{2}\sigma^2$ del exponente de $S_T$, y queda
\begin{equation}\label{eq:termino_I}
  (\mathrm{I}) = S_0\, e^{rT}\, N(d_1).
\end{equation}

\medskip

\noindent\textbf{La fórmula de Black-Scholes.} Sustituyendo
$(\mathrm{I})$ y $(\mathrm{II})$ en~\eqref{eq:descomposicion_call} y
multiplicando por el factor de descuento $e^{-rT}$, obtenemos la
fórmula central:

\begin{teorema}[Fórmula de Black-Scholes para la \emph{call} europea]
\label{teo:formula_BS}
El precio en $t = 0$ de la opción \emph{call} europea con \emph{strike}
$K$ y vencimiento $T$ en el modelo de Black-Scholes es
\begin{equation}\label{eq:formula_BS}
  \boxed{\;
    V(0, S_0) = S_0\, N(d_1) - K\, e^{-rT}\, N(d_2),
  \;}
\end{equation}
donde
\begin{equation*}
  d_1 = \frac{\log(S_0/K) + \bigl(r + \tfrac{1}{2}\sigma^2\bigr) T}
              {\sigma\, \sqrt{T}},
  \qquad
  d_2 = d_1 - \sigma\sqrt{T}
      = \frac{\log(S_0/K) + \bigl(r - \tfrac{1}{2}\sigma^2\bigr) T}
              {\sigma\, \sqrt{T}},
\end{equation*}
y $N(\cdot)$ es la función de distribución de la normal estándar.
\end{teorema}

\medskip

\noindent\textbf{Generalización al caso $t > 0$.} La fórmula
\eqref{eq:formula_BS} se ha derivado en $t = 0$. La generalización al
caso $t \in (0, T)$ se obtiene teniendo en cuenta la siguiente observación: el
proceso $\{S_t\}$ bajo $\Q$ es markoviano (pues es solución de la SDE
\eqref{eq:S_dinamica_Q}, que tiene coeficientes que solo dependen del estado
actual) y, por tanto, la esperanza condicional
$\E_\Q[(S_T - K)^+ \mid \F_t]$ depende exclusivamente de $S_t$, no de
toda la trayectoria pasada. Por la propiedad de Markov, calcular esta
esperanza condicional equivale a calcular la esperanza incondicional
del problema "iniciado" en $t$ con valor inicial $S_t$ y duración $T -
t$. Así, la derivación que hemos hecho para $t = 0$, repetida con $S_0$ sustituido por
$S_t$ y $T$ sustituido por $T - t$, nos da
\begin{equation}\label{eq:formula_BS_t}
  \boxed{\;
    V(t, S_t) = S_t\, N(d_1(t, S_t)) - K\, e^{-r(T-t)}\, N(d_2(t, S_t)),
  \;}
\end{equation}
donde
\begin{align*}
  d_1(t, S_t) &= \frac{\log(S_t/K) + \bigl(r + \tfrac{1}{2}\sigma^2\bigr)(T-t)}
                       {\sigma\, \sqrt{T-t}}, \\
  d_2(t, S_t) &= d_1(t, S_t) - \sigma\sqrt{T-t}.
\end{align*}
La fórmula~\eqref{eq:formula_BS_t} expresa el precio del \emph{call}
en cualquier instante $t \in [0, T)$ como función del precio actual
del subyacente $S_t$ y del tiempo restante $T - t$. Se trata de la \emph{función
de valor} del derivado, a la que volveremos en la siguiente sección al derivar la EDP de Black-Scholes.

\section{La EDP de Black-Scholes}
\label{sec:edp_bs}

Cerramos el capítulo con la deducción de la
EDP clásica de Black-Scholes a partir del cálculo de martingalas.

\medskip

La derivación que hemos articulado en las secciones precedentes (vía
cambio de medida, construcción de la $\Q$-martingala asociada al
\emph{payoff} y aplicación del teorema de representación) es la \emph{vía
de las martingalas} de la teoría de Black-Scholes. Existe una
derivación clásica alternativa, históricamente anterior, que parte de
la idea de \emph{cobertura instantánea} y conduce a una ecuación
diferencial en derivadas parciales. La presente subsección recupera esta segunda
formulación a partir de la fórmula~\eqref{eq:formula_BS_t}, mostrando
la consistencia entre ambas vías.

\medskip

\noindent\textbf{Aplicación de Itô a la función de valor.} La función
de valor del \emph{call} europeo, $V(t, S_t)$, dada
por~\eqref{eq:formula_BS_t}, es de clase $C^{1,2}$ sobre el dominio
$[0, T) \times (0, \infty)$. Aplicamos la fórmula de Itô
(Teorema~\ref{teo:formula_ito_general}) a $V(t, S_t)$, siendo $S_t$
proceso de Itô bajo $\Q$ con dinámica~\eqref{eq:S_dinamica_Q}:
\begin{equation}\label{eq:Ito_V}
  dV(t, S_t)
  = \biggl[\partial_t V + r\, S_t\, \partial_S V
           + \tfrac{1}{2}\, \sigma^2\, S_t^2\, \partial_{SS} V\biggr]\, dt
   + \sigma\, S_t\, \partial_S V\, d\widetilde W_t.
\end{equation}

\medskip

\noindent\textbf{Comparación con la dinámica del portfolio replicante.}
Por otro lado, en la subsección~\ref{subsec:paso3_replicacion}
demostramos que el valor del portfolio replicante satisface
(combinando \eqref{eq:dV_calculado} y la identificación
\eqref{eq:phi_BS}):
\begin{equation}\label{eq:dinamica_V_portfolio}
  dV_t = r\, V_t\, dt + \sigma\, S_t\, \phi_t\, d\widetilde W_t,
\end{equation}
donde $\phi_t = \widetilde\phi_t\, e^{rt} / (\sigma\, S_t)$ es la
cantidad de subyacente. Como $V(t, S_t) = V_t$ por construcción, las
expresiones~\eqref{eq:Ito_V} y~\eqref{eq:dinamica_V_portfolio} deben
coincidir término a término.

\emph{Igualando los coeficientes de $d\widetilde W_t$:}
$\sigma\, S_t\, \partial_S V = \sigma\, S_t\, \phi_t$, de donde
$\phi_t = \partial_S V$.

\emph{Igualando los coeficientes de $dt$:}
\begin{equation*}
  \partial_t V + r\, S_t\, \partial_S V
  + \tfrac{1}{2}\, \sigma^2\, S_t^2\, \partial_{SS} V
  = r\, V.
\end{equation*}
Como esta identidad debe valer para todo $(t, S_t)$ con $S_t$
recorriendo $(0, \infty)$ con probabilidad positiva, se traduce en una
ecuación diferencial parcial sobre la función $V(t, S)$:

\begin{teorema}[Ecuación de Black-Scholes]
\label{teo:edp_BS}
La función de valor $V(t, S)$ del \emph{call} europeo en el modelo de
Black-Scholes satisface la ecuación diferencial parcial
\begin{equation}\label{eq:EDP_BS}
  \boxed{\;
    \partial_t V + r\, S\, \partial_S V
    + \tfrac{1}{2}\, \sigma^2\, S^2\, \partial_{SS} V = r\, V,
  \;}
\end{equation}
sobre el dominio $(0, T) \times (0, \infty)$, con la condición de
contorno terminal
\begin{equation}\label{eq:condicion_terminal_BS}
  V(T, S) = (S - K)^+.
\end{equation}
\end{teorema}

La identidad~\eqref{eq:EDP_BS} es la \emph{ecuación de Black-Scholes}
clásica, en su forma original publicada en 1973. Por lo demostrado, la
función definida por~\eqref{eq:formula_BS_t} es solución de este
problema de contorno; un argumento de unicidad muestra que es
de hecho la única solución suficientemente regular
(Véase \cite{karatzas-shreve}, sección 5.8).

\medskip

\noindent\textbf{La conexión vía Feynman-Kac.} La equivalencia entre
las dos vías de derivación, martingalas y EDP, puede expresarse de forma abstracta a partir de la \emph{fórmula de Feynman-Kac}, que establece la
identidad entre la solución de la
EDP~\eqref{eq:EDP_BS}–\eqref{eq:condicion_terminal_BS} y la fórmula
probabilística~\eqref{eq:formula_valoracion}. Su demostración rigurosa
queda fuera del alcance de este trabajo y puede consultarse en
\cite{shreve}, capítulo 6.

% =====================================================================
%   CAPÍTULO 6 — CONCLUSIONES
% =====================================================================
\chapter{Conclusiones}\label{cap:conclusiones}
\section{Recapitulación del trabajo}
\label{sec:recapitulacion}

El trabajo se ha centrado en el problema de \emph{valoración
del derivado europeo}, abordado en su versión discreta
(Capítulo~\ref{cap:discreto}, modelo binomial) y en su versión
continua (Capítulos~\ref{cap:brown} a~\ref{cap:bs}, modelo de
Black-Scholes). En ambos contextos se observa la misma estructura organizada en tres pasos: cambio de medida hacia una medida neutral
al riesgo, construcción de la $\Q$-martingala asociada al
\emph{payoff} descontado, y aplicación de un teorema de
representación. La traducción entre ambos contextos no es trivial (combinatorial y constructiva en discreto, vía Girsanov y argumentos
abstractos en continuo), pero la estructura lógica es idéntica.

La relación estructural entre el modelo binomial discreto y el modelo
de Black-Scholes admite, además de esta analogía conceptual, una formulación cuantitativa precisa: la fórmula de Black-Scholes es
\emph{el límite} de la fórmula binomial multiperíodo cuando el número
de pasos tiende a infinito. Si en el árbol binomial sobre $[0, T]$
con $N$ períodos de longitud $\Delta t = T/N$ se definen los
parámetros de subida y bajada como
\begin{equation*}
  u = e^{\mu\Delta t+\sigma\sqrt{\Delta t}}, \qquad
  d = e^{\mu\Delta t-\sigma\sqrt{\Delta t}}, \qquad
  q = \frac{e^{r \Delta t} - d}{u - d},
\end{equation*}
y se hace $N \to \infty$, el random walk binomial converge en
distribución al GBM bajo $\Q$ (aplicando el teorema central
del límite), y el precio del \emph{call} obtenido por recursión finita
en el árbol converge a la fórmula de Black-Scholes
\eqref{eq:formula_BS}. Análogamente, la cantidad $\phi_k^A$ del
portfolio replicante discreto, dada por la fórmula
explícita~\eqref{eq:phi_explicito}, converge a $N(d_1)$. La
demostración rigurosa de estas convergencias puede consultarse en \cite{etheridge}, capítulo 2. La fórmula
de Black-Scholes no es, pues, un objeto independiente del modelo
binomial, sino su límite cuando los pasos temporales se hacen
infinitamente pequeños.

\section{Reflexión sobre el enfoque}
\label{sec:reflexion_enfoque}

La decisión estructural fundamental del trabajo ha sido abordar
Black-Scholes vía martingalas en lugar de vía EDP. Ambas
formulaciones son equivalentes, y su consistencia se expresa a
través de Feynman-Kac (Sección~\ref{sec:edp_bs}); si bien la vía
de las martingalas permite presentar el paso del modelo discreto
al continuo como una extensión natural, conservando las ideas
conceptuales del Capítulo~\ref{cap:discreto}. Esta transición,
menos evidente en una presentación basada en EDP, justifica el enfoque adoptado.


Conviene subrayar, por último, que el aparato desarrollado para abordar
esta vía (el cambio de medida, el teorema de Girsanov y la
representación de martingalas) posee un interés que va más allá del
problema de valoración: son herramientas centrales de la probabilidad y
del análisis estocástico modernos. El problema financiero ha sido,
de alguna forma, el hilo que ha motivado la introducción de
esa maquinaria.

\section{Líneas de continuación}
\label{sec:lineas_continuacion}

El trabajo se ha limitado a opciones europeas sobre un único subyacente
en el modelo de Black-Scholes con parámetros constantes. Cada una de
estas restricciones puede relajarse de diversas formas:

\begin{enumerate}[leftmargin=*]
  \item \emph{Opciones americanas}: incorporando la noción de
        tiempo de parada óptima al esquema desarrollado.
  \item \emph{Modelos más realistas}: volatilidad estocástica
        (Heston), saltos (Merton), tipos de interés estocásticos
        (Vasicek, Hull-White) y combinaciones de estos.
  \item \emph{Métodos numéricos}: simulación Monte Carlo y diferencias
        finitas para \emph{payoffs} sin fórmula cerrada (asiáticas,
        \emph{barrier}, \emph{basket options}, etc.).
  \item \emph{Calibración y volatilidad implícita}: estimación práctica
        de $\sigma$ a partir de precios de mercado de opciones líquidas.
\end{enumerate}

Cada una presupone los fundamentos desarrollados en el presente
trabajo y constituye un capítulo independiente de la matemática
financiera contemporánea.

% =====================================================================
%   BIBLIOGRAFÍA
% =====================================================================
\begin{thebibliography}{99}

\bibitem{bachelier}
L. Bachelier.
\emph{Théorie de la spéculation}.
Annales scientifiques de l'École Normale Supérieure, 3\textsuperscript{e}
série, 17, 21--86, 1900.

\bibitem{black-scholes}
F. Black y M. Scholes.
The pricing of options and corporate liabilities.
\emph{Journal of Political Economy}, 81 (3), 637--654, 1973.

\bibitem{etheridge}
A. Etheridge.
\emph{A Course in Financial Calculus}.
Cambridge University Press, Cambridge, 2002.

\bibitem{karatzas-shreve}
I. Karatzas y S. E. Shreve.
\emph{Brownian Motion and Stochastic Calculus}.
2\textsuperscript{a} edición, Graduate Texts in Mathematics, vol.~113.
Springer, Nueva York, 1991.

\bibitem{merton}
R. C. Merton.
Theory of rational option pricing.
\emph{Bell Journal of Economics and Management Science}, 4 (1),
141--183, 1973.

\bibitem{oksendal}
B.~\O{}ksendal.
\emph{Stochastic Differential Equations: An Introduction with
Applications}.
6\textsuperscript{a} edición, Universitext. Springer, Berlín, 2003.

\bibitem{shreve}
S. E. Shreve.
\emph{Stochastic Calculus for Finance II: Continuous-Time Models}.
Springer Finance. Springer, Nueva York, 2004.

\bibitem{williams}
D. Williams.
\emph{Probability with Martingales}.
Cambridge Mathematical Textbooks. Cambridge University Press,
Cambridge, 1991.

\end{thebibliography}

\end{document}
